% ====================================================================
% graphite_ica_ch1.tex — Chapter 1
%   흑연 음극 dQ/dV(ICA) 상전이 피크의 물리:
%   피크 {위치·너비·높이}가 {온도·C-rate·극판 전위}에 어떻게 반응하는가를
%   열역학·동역학·전기화학·통계로 유도해 실무 피팅에 쓸 닫힌 식에 도달한다.
% ====================================================================
\documentclass[11pt,a4paper]{article}
\usepackage[margin=22mm]{geometry}
\usepackage{kotex}
\usepackage{amsmath,amssymb,mathtools,bm}
\usepackage{booktabs,longtable,array}
\usepackage{enumitem}
\usepackage{amsthm}
\usepackage{hyperref}
\usepackage{fancyhdr}
\usepackage{lastpage}
\usepackage{setspace}
\setstretch{1.12}
\setlength{\parskip}{0.45em}\setlength{\parindent}{0pt}\setlength{\headheight}{14pt}
\hypersetup{colorlinks=true, linkcolor=blue, urlcolor=blue, citecolor=blue,
  pdftitle={흑연 음극 dQ/dV 이론: 평형·동역학·충방전 히스테리시스 — Chapter 1}}
\pagestyle{fancy}\fancyhf{}
\lhead{흑연 음극 $dQ/dV$ 이론 — Ch.1}\rhead{\thepage/\pageref{LastPage}}
\renewcommand{\headrulewidth}{0.4pt}

\newtheorem*{keybox}{핵심}
\newtheorem*{stagebox}{흑연 staging 배경}
\newtheorem*{boundbox}{유효 범위}
\newtheorem*{flagbox}{비고}
\newtheorem*{rigorbox}{심화}
\newtheorem*{verifybox}{검산}

\newcommand{\dd}{\mathrm{d}}
\newcommand{\eff}{\mathrm{eff}}
\newcommand{\eq}{\mathrm{eq}}
\newcommand{\app}{\mathrm{app}}
\newcommand{\drive}{\mathrm{drive}}
\newcommand{\cell}{\mathrm{cell}}
\newcommand{\bg}{\mathrm{bg}}
\newcommand{\ext}{\mathrm{ext}}
\newcommand{\OCV}{\mathrm{OCV}}
\newcommand{\sstat}{\mathrm{ss}}
\newcommand{\peak}{\mathrm{peak}}
\newcommand{\dis}{\mathrm{dis}}
\newcommand{\chg}{\mathrm{chg}}
\newcommand{\hys}{\mathrm{hys}}
\newcommand{\pol}{\mathrm{pol}}
\newcommand{\kin}{\mathrm{kin}}
\newcommand{\obs}{\mathrm{obs}}

\renewcommand{\theequation}{1.\arabic{equation}}
\renewcommand{\thesection}{1.\arabic{section}}
\title{\textbf{\LARGE Chapter 1}\\[0.35em]\Large 흑연 음극 $dQ/dV$ 이론: 평형 peak $\cdot$ 동역학 $\cdot$ 충방전 히스테리시스}
\author{}
\date{}

% TOC: widen number boxes (two-digit section numbers like 1.10 overflow article's 1.5em default)
\makeatletter
\renewcommand*\l@section[2]{%
  \ifnum \c@tocdepth >\z@
    \addpenalty\@secpenalty
    \addvspace{1.0em \@plus\p@}%
    \setlength\@tempdima{2.3em}%
    \begingroup
      \parindent \z@ \rightskip \@pnumwidth
      \parfillskip -\@pnumwidth
      \leavevmode \bfseries
      \advance\leftskip\@tempdima
      \hskip -\leftskip
      #1\nobreak\hfil
      \nobreak\hb@xt@\@pnumwidth{\hss #2\kern-\p@\kern\p@}\par
    \endgroup
  \fi}
\renewcommand*\l@subsection{\@dottedtocline{2}{2.3em}{3.2em}}
\makeatother

\begin{document}
\maketitle

% ====================================================================
\section*{서론}
\addcontentsline{toc}{section}{서론}
% ====================================================================
리튬이온전지 흑연 음극의 \textbf{incremental capacity analysis}(ICA)는 충방전 곡선 $Q(V)$ 의 미분
$\dd Q/\dd V$ 를 본다 \cite{bloom2005,dubarry2012}. 흑연은 리튬과 단계적 staging 상전이(stage
$4\!\to\!3\!\to\!2\mathrm L\!\to\!2\!\to\!1$)를 거치며, 여기서 $2\mathrm L$ 은 stage-2 의 희박(dilute) 단계를
가리킨다 \cite{dahn1991,ohzuku1993,funabiki1999}. 각 상전이가 일어나는 전위에서 $\dd Q/\dd V$ 에는 \textbf{peak}
가 선다. 왜 상전이가 peak 로 나타나는지부터 분명히 하자. 일차 상전이 동안 두 상이 공존하면 전위는 거의 평탄하게
머무는데, 그동안에도 전하는 계속 흘러 들어간다. 좁은 $\dd V$ 구간에 큰 $\dd Q$ 가 몰리는 것이므로, 미분
$\dd Q/\dd V$ 는 그 전위에서 치솟는다. 따라서 ``전위 평탄 구간 $=$ ICA peak''가 본 장 전체가 다루는 대상의
정의가 된다. 관측되는 유효 전이는 $j=1,\dots,N_p$ 개이고 흑연에서는 $N_p\approx3\sim4$ 이다. 인접 전이의 중심
전위 차가 전이폭 이하이면 두 peak 가 융합해 하나로 보이므로, 유효 전이 수는 결정학적 stage 수보다 적을 수 있다.
각 전이마다 전환율과 peak 가 따로 있다.

본 장이 설명하려는 \textbf{관측}은 다음이다.
\begin{quote}\itshape
$\dd Q/\dd V$ 의 상전이 peak 는 그 \textbf{위치$\cdot$너비$\cdot$높이}가 \textbf{(a) 온도, (b) C-rate(전류),
(c) 극판 전위} 세 가지에 따라 변한다. 특히 \emph{온도가 낮을수록 peak 의 꼬리(tail)가 길게 늘어져 다음 peak
와 겹치고}, 높을수록 짧게 끝난다. \textbf{나아가 같은 전이가 충전과 방전에서 서로 다른 전위에 peak 를 남긴다}
(\emph{충방전 히스테리시스}) — 이 전위 갈림은 전류를 $0$ 으로 줄여도 \emph{사라지지 않는} 열역학적 성분을 갖는다는
점에서, 단순한 분극(전류$\times$저항)과 구별된다.
\end{quote}

\textbf{이 장의 목표}는 이 관측을 화학열역학$\cdot$전기화학$\cdot$반응속도론$\cdot$통계와 미적분의 물리로,
\textbf{중간 단계를 빠짐없이 밝혀} 설명하는 \emph{닫힌 논리식}을 세우는 것이다. 그 식에는 두 가지를 요구한다.
첫째, 실데이터 피팅에 바로 쓸 수 있어야 한다. 둘째, 그 유도가 1차 원리에서 빈틈없이 닫혀야 한다. 이 두 요구가
본 장의 모든 절을 관통하는 기준이며, 마지막 절들(\S\ref{sec:master}$\cdot$\S\ref{sec:falsify})에서 피팅 절차와
반증 가능성으로 각각 회수된다.

\textbf{본론은 방전이고, 도착점은 방전 한 줄 피팅식이다.} 본 장의 주 전개는 \emph{방전}(탈리튬화) $\dd Q/\dd V$
다. 전하 보존이 정의하는 내부 전위에서 출발해 평형 peak, 정규용액 상호작용, 동역학 꼬리, 입자 분포를 차례로
닫으면, 방전 ICA 전체가 한 줄 피팅식(\S\ref{sec:master} 의 식~\eqref{eq:master})으로 모인다. 여기까지가 본 장의
본론이다. 충방전 히스테리시스는 그 본론 위에 세우는 \emph{확장}이다. 방전 전개에서 이미 도입한 정규용액 물리
(\S\ref{sec:regsol})가 상분리 분기를 만들기 때문에, 방전식의 전이 중심과 분극 부호만 방향별로 바꾸면 충$\cdot$방전
양방향 통합식(식~\eqref{eq:hys_master})이 방전 파라미터에 전이당 $\{\Omega_j,\gamma_j\}$ 둘만 더해 닫힌다
(\S\ref{sec:hys_branch} 이하; 충전 꼬리의 동역학만 방향별로 따로 닫는다). 어느 측정
(저율 OCV$\cdot$전류 차단$\cdot$다율$\cdot$다온도$\cdot$충방전 gap)이 어느 파라미터를 닫는지는
\S\ref{sec:master} 가 절차로 매듭짓는다. 그에 이르는 논리의 뼈대는 다음 순서다:
\begin{quote}
전하 보존이 내부 전위 $V_n$ 을 정의(\S\ref{sec:stage} — 모든 식의 전위축) $\to$ 평형 상승부(\S\ref{sec:eqpeak},
이상 극한) $\to$ \textbf{정규용액 상호작용$\cdot$상분리}(\S\ref{sec:regsol}) $\to$ 전위가 낮추는 동역학 꼬리
(\S\ref{sec:lag}$\cdot$\S\ref{sec:barrier}) $\to$ 입자 분포가 만드는 늘어짐(\S\ref{sec:dist}) $\to$ 여러 전이의
합산(\S\ref{sec:overlap}) $\to$ [확장] \textbf{충방전 분기}(\S\ref{sec:hys_branch}) $\to$ 통합식(\S\ref{sec:master}).
\end{quote}
이 사슬의 핵심 spine 은 정규용액 상호작용 $\Omega_j$ 하나가 peak 모양과 충방전 히스테리시스를 \textbf{동시에}
지배한다는 사실이다. $\Omega_j$ 는 방전 본론에서는 peak 폭을 좁히는 보정으로 들어오고(\S\ref{sec:regsol}), 확장에서는
분기 gap 을 만드는 원천이 된다(\S\ref{sec:hys_branch}). 따로 보이는 두 현상이 한 파라미터로 닫히는 것은 이 때문이다.
독자는 매 절 끝에서 ``이 결과가 최종 피팅식의 어느 항이 되는가''를 확인하게 된다.

\begin{stagebox}
흑연의 리튬화는 층간 갤러리를 단번에 채우지 않고 \emph{stage} 순서 $4\to3\to2\mathrm L\to2\to1$ 로 단계적으로
채운다 \cite{dahn1991,ohzuku1993}. stage 번호 $n$ 은 대략 ``리튬으로 채운 한 층마다 빈 갤러리 $n-1$ 개''를 뜻한다.
채울수록 $n$ 이 줄어 stage 1$=\mathrm{LiC_6}$(완전 채움)$\cdot$stage 2$=\mathrm{LiC_{12}}$ 가 되며, stage
3$\cdot$4$\cdot$2L 은 더 희박한 조성이다(조성은 근사). \textbf{각 stage 전이가 $\dd Q/\dd V$ 에 peak 를 하나씩
남긴다.} 인접 전이의 중심 전위 차가 전이폭 이하이면 융합해 관측 peak 수가 $N_p\approx3\sim4$ 가 되고, 이것이 본
장이 전이 $j=1,\dots,N_p$ 로 다루는 대상이다. 위 채움 순서는 \emph{리튬화(충전)} 방향임에 주의한다. 본 장의 주
전개는 \emph{방전}$=$탈리튬화 기준이므로 진행은 그 역순이다 — 점유율 $\theta_j$ 가 $1\to0$ 으로 줄고 stage 는
역방향으로 풀린다. 충방전 양방향의 분기는 \S\ref{sec:hys_branch} 이하에서 같은 골격 위에 세운다.
\end{stagebox}

\tableofcontents
\newpage

% ====================================================================
\section{기호와 규약 (Notation \& Conventions)}\label{sec:notation}
% ====================================================================
본문에 들어가기 전에 기호$\cdot$단위$\cdot$부호 규약을 정리한다. 각 양은 처음 쓰이는 곳에서 다시 물리적으로
동기 부여할 것이므로, 여기 표는 필요할 때 돌아와 찾아보는 사전으로 쓰면 된다.

\textbf{전이 인덱스 규약.} 상전이는 $j=1,\dots,N_p$ 로 여럿이다. $N_p$ 는 관측되는 유효 전이 수($\approx3\sim4$)
이며, 전하 보존식~\eqref{eq:charge} 의 합 상한이 된다. 각 전이에는 Li 점유율 $\theta_j$ 와 전환율(진행률)
$\xi_j\equiv1-\theta_j$ 를 둔다 — 방전(탈리튬화)이 진행하면 $\theta_j$ 가 $1\to0$ 으로 줄고 $\xi_j$ 가 $0\to1$
로 커진다. 그 평형값은 $\xi_{\eq,j}$, 완화 속도상수는 $k_j$, 정/역 방향 속도는 $r_j^\pm$, 지연은
$r_j\equiv\xi_{\eq,j}-\xi_j$, 꼬리 길이는 $L_{q,j}$ 로 적는다. 한 전이만 논하는 문맥에서는 $j$ 를 생략할 수
있으나, 박스로 강조된 최종 결과식에는 항상 $j$ 를 단다.

\textbf{부호 규약.} 진행 방향 부호 $s\in\{+1,-1\}$ 는 전이의 방전 방향을 표시하며, 본 장은 방전 기준 $s=+1$ 로
고정한다. 전류 부호 $s_I\in\{+1,-1\}$ 는 방전 시 분극이 apparent 전위를 올리는 방향을 $+$ 로 잡는 규약이고, 같은
방전 기준에서 $s_I=+1$ 이다. 충방전 양방향을 다룰 때 쓰는 방향 부호 $\sigma_d=\pm1$ 은 \S\ref{sec:master} 에서
도입한다. 셋의 역할 분담을 미리 요약해 두면 다음과 같다: $s$ 는 평형 곡선의 정의 안에 $+1$ 로 고정된 채 남는다 —
평형 곡선은 진행 방향에 불변인 상태함수이기 때문이다. 방향에 따라 실제로 바뀌는 것은 $\sigma_d$ 가 맡으며, 그
작용처는 셋이다 — 분극 부호, 분기 중심의 선택, 그리고 꼬리의 지지$\cdot$감쇠 방향(\S\ref{sec:master}).

\begin{longtable}{p{0.15\textwidth} p{0.13\textwidth} p{0.64\textwidth}}
\toprule 기호 & 단위 & 의미 \\ \midrule \endhead
$\theta_j$ & --- & 전이 $j$ 의 Li 점유율(분율), 방전 시 $1\to0$ \\
$\xi_j,\ \xi_{\eq,j}$ & --- & 전이 $j$ 의 전환율(진행률 $=1-\theta_j$)과 그 평형값, $0\le\xi_j\le1$ \\
$Q_j$ & C & 전이 $j$ 가 완전 진행 시 담는 용량 ($>0$) \\
$Q_\cell$ & C & 셀 기준 총 방전 용량 \\
$Q_\bg(V_n,T)$ & C & staging 으로 분리되지 않는 배경 누적 용량 \\
$C_\bg$ & C/V & 배경 미분용량 $=\partial Q_\bg/\partial V_n$ (chemical capacitance) \\
$q$ & --- & 무차원 용량 좌표 $=Q_\ext/Q_\cell=\int|I|\,\dd t/Q_\cell$ \\
$V_n$ & V & \emph{내부} 음극 전위 — 전하 보존식의 해 \\
$V_\app$ & V & 측정되는 apparent 전위 $=V_n+s_I|I|R_n$ \\
$V_\drive$ & V & 상변이 속도를 구동하는 전위; 기본 $V_\drive=V_n$ \\
$V_{n,\OCV}$ & V & 평형($\xi_j=\xi_{\eq,j}$)에서의 전하 보존식 해 \\
$U_j(T),\ U_j^0$ & V & 전이 $j$ 의 평형 중심 전위, 그 표준값 \\
$w_j$ & V & 평형 전이 폭(이상 극한서 $=RT/F$; 일반은 effective width) \\
$R_n$ & $\Omega$ & 음극 유효 분극 저항 \\
$\mu_\mathrm{Li},\ \mu^0$ & J/mol & 삽입 리튬의 \emph{점유율} 화학퍼텐셜과 그 표준값 \\
$\Omega_j$ & J/mol & 정규용액 상호작용 에너지 척도($>2RT$ 면 상분리$\cdot$히스) \\
$g_j(\xi),\ g_j''$ & J/mol & 전이 $j$ 의 정규용액 몰 자유에너지와 그 곡률 $\dd^2g_j/\dd\xi^2$(\S\ref{sec:regsol}) \\
$T_{c,j},\ \xi_{s,j}^\pm$ & K;\,--- & 전이 $j$ 임계온도 $=\Omega_j/2R$; spinodal 진행률 $=\tfrac12\pm\tfrac12\sqrt{1-2RT/\Omega_j}$ \\
$\gamma_j$ & --- & 분기 축소 인자(다입자) $0\le\gamma_j\le1$ ($1$=spinodal 상한, $0$=가역) \\
$\Delta U_j^\hys$ & V & 충방전 평형 분기 gap $=\tfrac{2}{sF}[\Omega_j u_j-2RT\,\mathrm{artanh}\,u_j]$, $u_j=\sqrt{1-2RT/\Omega_j}$ \\
$U_j^{\dis},\,U_j^{\chg}$ & V & 방전$\cdot$충전 분기 중심 $=U_j\pm\tfrac12\gamma_j\Delta U_j^\hys$ \\
$\mathcal A_j$ & J/mol & 전이 $j$ 의 구동력(affinity) $=sF(V_\drive-U_j)$($\Omega\ne0$ 일반화는 \S\ref{sec:barrier} 주) \\
$\Delta G_{a,j}$ & J/mol & 활성화 자유에너지 $=\Delta H_{a,j}-T\Delta S_{a,j}$ \\
$\Delta H_{a,j},\Delta S_{a,j}$ & {\footnotesize J/mol;\,J/(mol\,K)} & 활성화 엔탈피$\cdot$엔트로피(\emph{속도}에 들어감) \\
$\Delta G_{\eff,j}$ & J/mol & 유효 활성화 자유에너지 $=\Delta G_{a,j}-\chi\mathcal A_j$ \\
$\Delta S_j$ & J/(mol\,K) & 전이 $j$ 의 \emph{삽입(리튬화) 반응} 엔트로피($\partial U_j/\partial T=\Delta S_j/(sF)$; 활성화 $\Delta S_{a,j}$ 와 \emph{다른 양}) \\
$k_j,\ k_0$ & 1/s & 완화 속도상수 $=r_j^++r_j^-$, Eyring prefactor $k_0=k_BT/h$ \\
$k_j^{\mathrm{fwd}}$ & 1/s & 꼬리 영역 forward 근사 $\simeq r_j^+$($\mathcal A_j\gtrsim RT$; \S\ref{sec:lag}) — 직렬 합성 $k_j^\eff$ 와 다른 양 \\
$k_j^{\eff},k_{j,\lim},$ $k_{j,\mathrm{act}}$ & 1/s & 병목 포함 유효 속도 $1/k_j^{\eff}=1/k_{j,\mathrm{act}}+1/k_{j,\lim}$ ($k_{j,\mathrm{act}}\!\equiv\!k_j$=활성화 완화속도, $k_{j,\lim}$=직렬 율속) \\
$r_j^\pm$ & 1/s & 전이 $j$ 의 정/역 방향 속도 \\
$\xi_{\sstat,j},\ \bar r_j$ & --- & 정상점 $=r_j^+/(r_j^++r_j^-)$; 가중 평균 지연 $=\int\rho_G\,r_j\,\dd G$ \\
$\chi$ & --- & 전달 계수(transfer coefficient) $0\le\chi\le1$ (전기화학 표준 표기 $\alpha$; 본 장 $\chi$) \\
$\lambda$ & J/mol & Marcus 재구성 에너지 \\
$\rho_G(G),\ \bar G$ & mol/J;\,J/mol & 활성화 배리어의 \emph{용량-분율} 밀도(단위 정규화 $\int\rho_G\,\dd G=1$; 총 용량 $Q_j$ 는 별도 곱)와 우세 집단 대표값 \\
$\sigma_G$ & J/mol & 배리어 분포의 표준편차 \\
$q_a,\ V_a$ & ---,\ V & 꼬리 시작 좌표와 그 전위 $V_a=V_n(q_a)$ \\
$r_{a,j}$ & --- & 꼬리 시작점 잔류 지연 $=r_j(q_a)=\xi_{\eq,j}(q_a)-\xi_j(q_a)$, $0<r_{a,j}\le\xi_{\eq,j}(q_a)$(포화창서 상한 $\simeq1$; \S\ref{sec:synth}) \\
$L_{q,j}$ & --- & 용량축 꼬리 특성 길이 $=|I|/(Q_\cell k_j)$ (공율속 시 $k_j\!\to\!k_j^{\eff}$; \S\ref{sec:barrier}) \\
$L_{V,j}$ & V & 전압축 꼬리 길이 $=|\dd V_n/\dd q|_{q_a}L_{q,j}$ \\
$R,F,k_B,h,N_A$ & {\footnotesize J/(mol\,K), C/mol, J/K, J\,s, mol$^{-1}$} & 기체상수, Faraday, Boltzmann, Planck, Avogadro ($R=k_BN_A$) \\
\bottomrule
\end{longtable}

% ====================================================================
\section{전하 보존이 정하는 내부 전위 $V_n$}\label{sec:stage}
% ====================================================================
상전이를 구동하는 것은 음극 \emph{내부}의 전위 $V_n$ 이다. 그런데 회로에서 읽는 전압은 이 내부 전위와 같지
않다. 전류가 흐르는 동안에는 분극에 의한 전위 강하가 측정값에 함께 섞여 들어오기 때문이다. 그렇다면 $V_n$ 은
어디서 오는가. 본 장의 답은, $V_n$ 이 외부에서 주어지는 입력이 아니라 \textbf{전하 보존이라는 한 등식의 해}라는
것이다. 이 절은 그 전위 $V_n$ 의 정의를 세운다. 이후의 모든 절 — 특히 ``전위가 배리어를 낮춘다''는 동역학 논의
(\S\ref{sec:barrier}) — 가 이 전위축 위에서 전개된다.

\textbf{전하 보존식.} 외부 회로로 흘린 방전 전하 $Q_\cell q$ 는 음극의 \emph{상태 변화}로 정확히 계상된다.
방전은 탈리튬화로 Li 를 빼내는 과정이므로, 전하를 어딘가에 ``저장''한다는 그림이 아니라 전하 \emph{수지}의 그림이
맞다. 계상의 근거는 Faraday 법칙이다. 외부로 전자 하나가 나갈 때마다 음극에서 Li$^+$ 하나가 빠져나가므로, 흘린
전하와 $F\times$(빠진 Li 몰수)는 항등으로 같다. 빠진 Li 는 \emph{서로소인 두 종류의 자리}에서 나오며, 그래서
합이 교차항 없이 선형으로 갈라진다. 첫째 몫은 뚜렷한 staging 전이로 분리되는 $\sum_j Q_j\xi_j$ 다 — 전이 $j$ 가
진행률 $\xi_j=1-\theta_j$ 만큼 진행하면, 곧 점유율 $\theta_j$ 가 그만큼 줄면, 그 용량 기여는 $Q_j\xi_j$ 가 된다.
둘째 몫은 그렇게 분리되지 않는 배경 $Q_\bg(V_n,T)$ 다. 외부로 흘린 전하가 이 두 몫의 합과 같다는 전하 수지(보존)가
\begin{equation}
\boxed{\;Q_\cell\,q=Q_\bg(V_n,T)+\sum_{j=1}^{N_p}Q_j\,\xi_j\;.}
\label{eq:charge}
\end{equation}
좌변의 $q=\int|I|\,\dd t/Q_\cell$ 는 궤적을 따라 단조 증가하는 외부 입력, 곧 아는 값이다. 우변은 내부 상태 —
$V_n$ 과 $\{\xi_j\}$ — 의 함수다. $\xi_j$ 는 방전에서 $0\to1$ 로 진행하며, 충전 분기는 진행이 역순일 뿐 식의
형태는 같다(충방전 양방향의 처리는 \S\ref{sec:hys_branch} 이하). 주어진 $(q,T,\{\xi_j\})$ 에서 이 식을 만족하는
$V_n$ 이 곧 내부 전위다. 즉 $V_n$ 은 외부 lookup 테이블에서 읽는 값이 아니라 보존식의 해이며, 이는 다공성 전극의
implicit 화학퍼텐셜 관계와 같은 구조다 \cite{mckinnon1983,newman2004}.

\textbf{해의 유일성(배경 미분용량).} 식~\eqref{eq:charge} 가 $V_n$ 에 대해 유일해를 가지려면 우변이 $V_n$ 의
단조 함수여야 한다. 배경 용량의 기울기
\begin{equation}
C_\bg\equiv\frac{\partial Q_\bg}{\partial V_n}\ \ge\ \epsilon_Q\ >\ 0
\label{eq:cbg}
\end{equation}
가 양이면 충분하다. $C_\bg>0$ 의 물리적 출처는 다음과 같다. 배경 몫은 staging 으로 묶이지 않는 자유도 — 희박
고용체 영역의 연속 삽입, 표면과 결함 자리, 전기이중층 — 인데, 이들이 단상 안정 조건
$\partial\mu/\partial(\text{조성})>0$ 을 만족하는 한 ``전위가 낮아질수록 Li 가 더 들어가는'' 단조 응답을 보이므로
미분용량이 양수가 된다. 이는 동어반복이 아니라 열역학 안정성의 귀결이다. 더 정확히 말하면, 우변이 (i) $V_n$ 에
연속이고 (ii) 순증가($C_\bg>0$)이며 (iii) 목표값이 우변의 치역 안에 있을 때 중간값 정리에 의해 해가 존재하고
유일하다 \cite{bazant2013}. 균일 하한 $\epsilon_Q>0$ 까지 두면 역함수가 Lipschitz 가 되어 수치 해의 조건수가
좋아진다. 다만 유일성 자체에는 점별 $C_\bg>0$ 로 족하다 — $\epsilon_Q$ 는 유일성 조건이 아니라 수치 안정성
조건이다.

\textbf{한 시점의 풀이 vs 궤적 — 순환 고리의 정직한 처리.} 먼저 무엇이 순환인지부터 명시한다. 다음 절에서
보겠지만 $\xi_{\eq,j}$ 는 $V_n$ 의 함수다. 그런데 그 $V_n$ 자체가 $\xi_j$ 를 입력으로 받는 식~\eqref{eq:charge}
의 해다. 이 고리를 모순 없이 닫는 길은 \emph{미분 경로의 구분}이다. 한 순간 — 용량 좌표 $q$ 를 고정한 시점 —
에는 그때의 동역학 상태 $\{\xi_j(q)\}$ 를 고정한 채 식~\eqref{eq:charge} 를 $V_n$ 에 대해 푼다. 여기서
$\{\xi_j(q)\}$ 는 임의 입력이 아니라 완화 운동방정식(\S\ref{sec:lag})이 정해 둔 궤적의 현재 상태다. $\xi_j$ 가
고정되면 우변은 $Q_\bg(V_n,T)+\text{(상수)}$ 형태이므로 $C_\bg>0$ 단조성으로 $V_n$ 이 닫힌다. 그러나 \emph{관측}되는
$\dd Q/\dd V_n$ 은 한 순간이 아니라 궤적을 따라 보는 양이다. 따라서 $\xi_j$ 와 $V_n$ 이 모두 $q$ 의 함수임을
살려, 연쇄율 $\dd Q/\dd V_n=(\dd Q/\dd q)\,/\,(\dd V_n/\dd q)$ 로 닫아야 한다. $\xi_j$ 를 $V_n$ 의 \emph{직접}
함수로 보고 ``$\dd\xi_j/\dd V_n$''을 그대로 쓸 수 있는 것은 \emph{평형 분기}에 한한다(아래 문단). 평형과 비평형의
미분 경로를 이렇게 구분해야 self-consistent 고리 $V_n\to\xi_{\eq}\to V_n$ 이 모순 없이 닫힌다. 분모
$\dd V_n/\dd q\to0$ 인 이상적 plateau 에서 연쇄율이 발산하는 것은 결함이 아니다. 그것이 바로 서론의 ``전위 평탄
구간 $=$ peak'' 그 자체이며, 무한히 날카로운 이상 peak 를 유한 폭으로 만드는 것은 뒤에서 더할 broadening — 평형 폭
(\S\ref{sec:eqpeak})$\cdot$분포(\S\ref{sec:dist})$\cdot$동역학(\S\ref{sec:lag}) — 이다. 연쇄율은 $V_n(q)$ 가
구간 단조인 곳에서 정의된다.

\textbf{평형 극한 $=$ OCV, 그리고 유일성의 한계.} 평형 또는 충분히 느린 조건에서는 $\xi_j=\xi_{\eq,j}(V_n,T)$
(다음 절)이고, 식~\eqref{eq:charge} 의 해가 \emph{평형 음극 전위} $V_{n,\OCV}(q,T)$ 다. 여기서는 \S\ref{sec:regsol}
의 결과 두 가지만 미리 빌려 쓴다. 첫째, 평형 기울기는 $\dd\xi_{\eq,j}/\dd V_n=sF/g_j''$ 이어서 자유에너지
\emph{곡률} $g_j''$ 의 부호가 기울기의 부호를 정한다. 둘째, 상호작용이 강하면($\Omega_j>2RT$) $g_j''<0$ 인 구간
(spinodal)이 생긴다. 이를 받아 우변의 기울기 $C_\bg+\sum_j Q_j\,\dd\xi_{\eq,j}/\dd V_n$ 을 영역별로 본다.
\textbf{단상 영역}($\Omega_j<2RT$, $g_j''>0$)에서는 $\dd\xi_{\eq,j}/\dd V_n>0$ 이라 우변이 더 가파른 단조 증가가
되고, 해는 유일하다. 임계 $\Omega_j=2RT$ 에서는 $\xi=\tfrac12$ 한 점에서 $g''=0$ 이 되어 기울기가 발산하지만
유일성 자체는 유지된다. \textbf{그러나 상분리 영역}($\Omega_j>2RT$)은 사정이 다르다. spinodal 불안정 구간
$\xi_{s,j}^-<\xi<\xi_{s,j}^+$ 에서만 $g_j''<0$ 이고, 거기서 $\dd\xi_{\eq,j}/\dd V_n=sF/g_j''<0$ 이므로
$V_{\eq,j}(\xi)$ 가 비단조 loop 를 그린다(\S\ref{sec:regsol}; 양 어깨는 $g_j''>0$ 이라 양의 기울기다). 우변이
비단조가 되면 한 $V_n$ 에 여러 해가 가능해진다. 이 \emph{다가성}이 곧 충$\cdot$방전 분기, 즉 히스테리시스
(\S\ref{sec:hys_branch})의 뿌리다. 정리하면 ``OCV $=$ 전하 보존식의 평형 해''라는 명제는 단상에서는 유일해로
성립하고, 상분리 구간에서는 이력 — 충전에서 왔는가 방전에서 왔는가 — 이 고르는 \emph{분기에 대한 조건부} 명제로
읽어야 한다. 전역 유일성은 성립하지 않는다. 한편 전류가 흐르면 $\xi_j$ 가 평형을 따라가지 못해(\S\ref{sec:lag})
$V_n$ 이 OCV 에서 이탈한다. 이 이탈이 뒤에서 다룰 꼬리의 바탕이다.

\textbf{측정 전위와 구동 전위.} 측정되는 전위는 내부 전위에 분극(저항 강하)이 더해진 값이다:
\begin{equation}
V_\app=V_n+s_I|I|R_n.
\label{eq:vapp}
\end{equation}
상전이의 \emph{속도}를 구동하는 것은 그중 구동 전위 $V_\drive$ 다. 본 장의 기본형은 $V_\drive=V_n$ 으로, 계면
charge-transfer 과전압을 내부 전위에 흡수한 lumped 근사다. 이 근사는 상변이의 율속이 계면 전자전달이 아니라 고상
내 재배열$\cdot$핵생성에 지배될 때 타당하다. 과전압을 별도 항으로 분리하는 일반형은, 그 가정이 깨지는 조건과 함께
이후 장에서 다룬다. 세 전위 $V_n$(내부 해)$\cdot$$V_\app$(측정)$\cdot$$V_\drive$(구동)는 표기만 다른 같은 양이
아니라 서로 다른 측정과 역할을 갖는, 물리가 다른 양이다. 본 장 전체에서 이 구분을 일관되게 유지한다.

\begin{verifybox}
\textbf{Worked example — 보존식이 $V_n$ 을 정하는 한 장면.} 감을 잡기 위해 $N_p=3$,
$Q_1=Q_2=Q_3=0.3\,Q_\cell$, 배경 몫 $0.1\,Q_\cell$ 의 가상 음극을 보자. 방전이 $q=0.5$ 까지 진행했고 그
시점의 동역학 상태가 $\xi_1=1.00$, $\xi_2=0.55$, $\xi_3=0.02$ 라 하자. 전이 몫의 합은
$0.3\times(1.00+0.55+0.02)=0.471\,Q_\cell$ 이므로, 식~\eqref{eq:charge} 는 배경이
$Q_\bg(V_n,T)=0.029\,Q_\cell$ 을 담도록 $V_n$ 을 요구한다. $C_\bg>0$ 의 단조성에 의해 그런 $V_n$ 은 정확히
하나다. 이것이 ``$V_n$ 은 보존식의 해''라는 명제의 구체적 모습이며, 피팅 알고리즘(\S\ref{sec:master})이
미측정 조건을 예측할 때 평형 해 $V_{n,\OCV}(q)$ 를 생성하는 것도 같은 풀이다.
\end{verifybox}

\begin{keybox}
\textbf{전하 보존의 요지.} $V_n$ 은 전하 보존식~\eqref{eq:charge} 의 해다 — 외부에서 주어지는 값이 아니다.
평형이면 그 해가 OCV 이고, 전류가 흐르면 $\xi_j$ 의 지연만큼 $V_n$ 이 OCV 에서 밀린다. 측정은
$V_\app=V_n+s_I|I|R_n$, 속도 구동은 $V_\drive=V_n$ 이다. 이 바탕 위에서 다음 절들이 평형 목표 $\xi_{\eq,j}$,
그 지연, 전위에 의한 배리어 낮춤을 차례로 전개한다.
\end{keybox}

% ====================================================================
\section{평형에서의 peak — 기준선}\label{sec:eqpeak}
% ====================================================================
앞 절에서 내부 전위 $V_n$ 을 전하 보존식의 해로 정의했다. 이제 그 전위축 위에서 첫 질문을 묻는다. 전류가 거의
흐르지 않는 \emph{평형}에서 전이 $j$ 의 peak 는 어디에 서고, 얼마나 넓고, 얼마나 높은가. 그리고 그것이 온도에
따라 어떻게 변하는가. 이 기준선을 먼저 세워야, 전류가 흐를 때의 변화를 그로부터의 \emph{이탈}로 읽을 수 있다.

\subsection{$\dd Q/\dd V_n$ 의 peak 발생}
방전 궤적에서는 외부 전하 $Q\equiv Q_\ext=Q_\cell q$ 와 내부 상태 $\{\xi_j\}$ 가 모두 $V_n$ 을 따라 변하며,
그 관계를 전하 보존식~\eqref{eq:charge} 이 묶는다. 평형에서는 $\xi_j=\xi_{\eq,j}(V_n)$ 이다(아래 소절). 따라서
식~\eqref{eq:charge} 의 양변을 등온 조건에서 $V_n$ 으로 미분하면 관측 ICA $\dd Q/\dd V_n=\dd(Q_\cell q)/\dd V_n$
이 닫힌다:
\begin{equation}
\frac{\dd Q}{\dd V_n}=\frac{\partial Q_\bg}{\partial V_n}+\sum_j Q_j\frac{\dd\xi_j}{\dd V_n}
=C_\bg+\sum_j Q_j\frac{\dd\xi_j}{\dd V_n}.
\label{eq:dQdV}
\end{equation}
배경 $C_\bg$ 는 완만한 기저선이고, 한 전이 $j$ 가 진행되는 좁은 전위 구간에서는 그 항 $Q_j\,\dd\xi_j/\dd V_n$ 이
급증해 \textbf{peak} 를 만든다. \emph{즉 peak 는 $\xi_j$ 가 $V_n$ 에 대해 가장 급히 변하는 곳}이다. 그러므로
peak 의 모양은 곧 $\xi_j(V_n)$ 의 기울기 모양이며, 평형에서는 $\xi_j=\xi_{\eq,j}(V_n)$ 다. 이제 $\xi_{\eq,j}$ 를
격자기체(lattice-gas) 모형으로 유도한다.

\subsection{평형 전환율 $\xi_{\eq,j}(V_n)$ — 격자기체(lattice-gas) 유도}
\textbf{유도의 출발점.} 한 전이를 ``격자 자리(site)에 리튬이 차고 비는'' 격자기체(lattice-gas)로 보면, \emph{점유율}
$\theta_j$(Li 가 자리를 채운 분율)의 평형값은 \emph{자리를 더 채우려는 화학적 경향과 전극 전위의 균형}에서 정해진다.
그 경향을 \emph{점유율의} 화학퍼텐셜 $\mu_\mathrm{Li}(\theta_j)$ 로 적어 두 출처(섞임 엔트로피, 입자 간 상호작용)를 한
줄씩 유도한 뒤, \emph{방전 진행률} $\xi_j=1-\theta_j$ 로 옮긴다 \cite{mckinnon1983,bazant2013}.

\emph{(i) 섞임 엔트로피.} $N$ 개 동등한 자리 중 $n$ 개가 점유된 배열의 수는 $W=\dfrac{N!}{n!\,(N-n)!}$ 이다.
Boltzmann 관계 $S_{\mathrm{mix}}=k_B\ln W$ 에 Stirling 근사 $\ln m!\simeq m\ln m-m$ 를 세 계승에 적용하면
\begin{align}
\ln W&=\ln N!-\ln n!-\ln(N-n)!\notag\\
&\simeq\big(N\ln N-N\big)-\big(n\ln n-n\big)-\big((N-n)\ln(N-n)-(N-n)\big)\notag\\
&=N\ln N-n\ln n-(N-n)\ln(N-n)\qquad\big[\text{선형항}\ -N+n+(N-n)=0\big].
\label{eq:lnW}
\end{align}
점유율 $\theta_j\equiv n/N$ 을 넣으면($n=N\theta_j$, $N-n=N(1-\theta_j)$, 그리고 $N\ln N$ 을 두 항에 분배하면
$N\ln N=n\ln N+(N-n)\ln N$ 이라 $\ln(n/N),\ln((N-n)/N)$ 로 묶인다)
\begin{equation}
S_{\mathrm{mix}}=-k_B N\big[\theta_j\ln\theta_j+(1-\theta_j)\ln(1-\theta_j)\big].
\label{eq:smix}
\end{equation}
혼합 자유에너지는 $G_{\mathrm{mix}}=-T\,S_{\mathrm{mix}}$ 다. 이를 \emph{격자 자리(전체 site) 1몰당}으로 환산하면(전체
자리 수 $N$ 으로 나누고 $N_A$ 를 곱해 $R=k_B N_A$), 1몰당 혼합 자유에너지 $\bar G_{\mathrm{mix}}=RT[\theta_j\ln\theta_j+
(1-\theta_j)\ln(1-\theta_j)]$ 이고, 화학퍼텐셜의 섞임 몫은 이를 $\theta_j$ 로 미분한 것이다:
\begin{equation}
\frac{\partial\bar G_{\mathrm{mix}}}{\partial\theta_j}=RT\big[(\ln\theta_j+1)-(\ln(1-\theta_j)+1)\big]
=RT\ln\frac{\theta_j}{1-\theta_j}
\label{eq:mumix}
\end{equation}
($\theta_j\ln\theta_j$ 미분 $=\ln\theta_j+1$, $(1-\theta_j)\ln(1-\theta_j)$ 미분 $=-(\ln(1-\theta_j)+1)$, 두 $+1$ 이 상쇄).

\emph{(ii) 상호작용.} 정규용액(regular-solution) 근사는 \emph{평균장}에서 점유-인접 상호작용을 본다 — 무작위
배치에서 ``점유$\cdot$빈자리'' 쌍이 생길 확률이 $\propto\theta_j(1-\theta_j)$ 이므로 혼합 엔탈피가
$\bar G_{\mathrm{int}}=\Omega_j\,\theta_j(1-\theta_j)$ 이다 \cite{mckinnon1983}. 미분하면(곱미분
$\dd[\theta_j-\theta_j^2]/\dd\theta_j=1-2\theta_j$)
\begin{equation}
\frac{\partial\bar G_{\mathrm{int}}}{\partial\theta_j}=\Omega_j(1-2\theta_j).
\label{eq:muint}
\end{equation}
표준 상태 $\mu^0$ 까지 더해 삽입 리튬의 \emph{점유율} 화학퍼텐셜은
\begin{equation}
\mu_\mathrm{Li}(\theta_j)=\mu^0+RT\ln\frac{\theta_j}{1-\theta_j}+\Omega_j(1-2\theta_j).
\label{eq:mu}
\end{equation}

\emph{(iii) 전기화학 평형 조건 — 점유율 균형에서.} 식~\eqref{eq:mu} 의 $\mu_\mathrm{Li}$ 는 \emph{점유율
$\theta_j$ 의} 화학퍼텐셜이며, 점유율이 높을수록 크다. 전극에서의 삽입 반응 $\mathrm{Li}^++e^-+[\,]\rightleftharpoons
\mathrm{Li}_{(\text{graphite})}$ ($[\,]$=흑연 빈 자리; 각 stage 전이 $j$ 를 이 반응의 \emph{유효}(coarse-grained)
형으로 본다)의 평형은 양변의 전기화학퍼텐셜이 같은 조건이다: $\mu_\mathrm{Li}(\theta_{\eq,j})=\tilde\mu_{\mathrm{Li}^+}
+\tilde\mu_{e^-}$. 전자의 전기화학퍼텐셜은 $\tilde\mu_{e^-}=\mu_{e^-}^0-FV_n$ 으로 \emph{전기적 일} 항 $-FV_n$ 을
갖는다. 이 항은 전자 1몰을 전위 $V_n$ 의 전극에 놓는 일이며, 전하 $-F$ 와 전위 $V_n$ 의 곱이다. 전해질 활동도가
일정한 조건에서
$\mathrm{Li}^+\cdot$표준 항을 한 상수로 묶어 $\tilde\mu_{\mathrm{Li}^+}+\mu_{e^-}^0\equiv sF\,U_j^0$ 로 전위 $U_j^0$ 를
정의하면, 평형은 \emph{점유율} 화학퍼텐셜이 전위에 묶인 형이 된다. 전기적 일은 $-FV_n$(방향 무관)이고, 진행 방향
부호 $s$ 는 \emph{표기 규약}으로 $sF$ 에 넣어 둔다 — $s=+1$ 에서 아래 식은 $-F(V_n-U_j)$ 와 동일하며, $s$ 는 평형
관계 자체를 바꾸지 않고 진행률 정의($\xi=1-\theta$)와 짝지어 방향만 표시한다:
\begin{equation}
\mu_\mathrm{Li}(\theta_{\eq,j})=\mu^0-sF\big(V_n-U_j\big),\qquad U_j\equiv U_j^0-\frac{\mu^0}{sF}
\label{eq:eqcond}
\end{equation}
\textbf{부호의 물리적 의미.} 좌변 $\mu_\mathrm{Li}$ 는 점유율이 높을수록 크다. 이는 이상 또는 단상
($\Omega_j<2RT$) 영역의 성질이며, $\Omega_j>2RT$ 의 spinodal 구간에서는 깨진다(\S\ref{sec:regsol}). 전자 일이
$-FV_n$ 이므로 $V_n$ 이 오르면 평형 점유율은 낮아져야 하고, 따라서 우변은 $V_n$ 에 대해 감소한다. 이 음의 부호가
곧 ``$V_n\uparrow\Rightarrow$ 탈리튬화 진행''의 물리다. 평형을 진행률 $\xi_j$ 가 아니라 점유율 $\theta_j$ 로
적었기 때문에, 이 부호는 $s$ 규약에 가려지지 않고 명시적으로 드러난다. 부호 $s$ 는 유도되는 양이 아니라 진행률
정의에 묶인 규약이다. 방전 기준 $\xi_j=1-\theta_j$ 와 함께 $s=+1$ 로 두면 $\dd\xi_{\eq,j}/\dd V_n>0$
(식~\eqref{eq:dxidV})이 되어 ``방전에서 $V_n\uparrow\Rightarrow$ 진행$\uparrow$''이라는 관측과 일치한다 — 관측이
규약을 고정하는 것이 아니라, 규약의 자연스러움을 관측이 확인해 주는 것이다. 충방전 양방향에서도 $s$ 는 $+1$ 로
고정되고, 방향은 $\sigma_d$ 가 맡는다(\S\ref{sec:notation}$\cdot$\S\ref{sec:master}).

\emph{(iv) 진행률 $\xi_{\eq,j}$ 로 옮겨 풀기.} 본 장 본문은 먼저 \textbf{이상 극한 $\Omega_j=0$ 을 채택한다}.
이는 ``흑연이 이상 용액이다''라는 물리 주장이 아니다 — 실제로는 \S\ref{sec:regsol} 에서 보듯 $\Omega_j\ne0$ 이
흑연 staging 의 본질이다. 이상 결과를 기준선으로 세워 두면 모든 상호작용 보정 — 폭 $w_j$ 의 좁아짐, 상분리 — 을
그 기준선으로부터의 이탈로 측정할 수 있다. 즉 이상 극한의 채택은 물리 가정이 아니라 전개 순서의 선택이다.
$\Omega_j$ 항은 \S\ref{sec:regsol} 에서 임계까지 포함해 도입하며, 작은 $\Omega_j$ 는 폭 보정의 섭동으로,
$\Omega_j>2RT$ 는 비섭동 상분리로 들어온다. 폭을 $w_j\equiv RT/F$ 로
두면 식~\eqref{eq:eqcond}($\Omega_j=0$)은 $RT\ln[\theta_{\eq,j}/(1-\theta_{\eq,j})]=-sF(V_n-U_j)$ 다. 이제
\emph{점유율을 방전 진행률로} 옮긴다 — $\theta_{\eq,j}=1-\xi_{\eq,j}$ 이므로 $\ln\dfrac{\theta_{\eq,j}}{1-\theta_{\eq,j}}
=\ln\dfrac{1-\xi_{\eq,j}}{\xi_{\eq,j}}=-\ln\dfrac{\xi_{\eq,j}}{1-\xi_{\eq,j}}$ 로 \emph{부호가 한 번 명시적으로
뒤집힌다}(이 반전이 ``점유율$\downarrow$ = 진행률$\uparrow$''의 대수적 표현이다):
\begin{align}
-RT\ln\frac{\xi_{\eq,j}}{1-\xi_{\eq,j}}&=-sF(V_n-U_j)\notag\\
\ln\frac{\xi_{\eq,j}}{1-\xi_{\eq,j}}&=\frac{s(V_n-U_j)}{w_j}\qquad[\div(-RT),\ w_j=RT/F]\notag\\
\frac{\xi_{\eq,j}}{1-\xi_{\eq,j}}&=E,\qquad E\equiv\exp\!\Big[\frac{s(V_n-U_j)}{w_j}\Big]\qquad[\text{양변 지수}]\notag\\
\xi_{\eq,j}&=E(1-\xi_{\eq,j})\ \Rightarrow\ \xi_{\eq,j}(1+E)=E\ \Rightarrow\ \xi_{\eq,j}=\frac{E}{1+E}.\notag
\end{align}
분자$\cdot$분모를 $E$ 로 나누면 $1/E=\exp[-s(V_n-U_j)/w_j]$ 이므로
\begin{equation}
\boxed{\;\xi_{\eq,j}(V_n,T)=\frac{1}{1+\exp\!\big[-s\,(V_n-U_j(T))/w_j(T)\big]}\;,\qquad w_j=\frac{RT}{F}\;.}
\label{eq:logistic}
\end{equation}
이것은 매끄러운 logistic 곡선이지 $0\!\to\!1$ 급점프가 아니다. 본 장은 어떤 step-function 도 쓰지 않는다.

\begin{flagbox}
\textbf{평형 \emph{형태}는 데이터가 정한다.} 위 logistic 은 이상 lattice-gas($\Omega_j=0$)의 닫힌 결과다.
$\Omega_j\ne0$ 이면 식~\eqref{eq:eqcond} 에 상호작용 항 $\Omega_j(1-2\theta_{\eq,j})$ — 진행률로 옮기면
$-\Omega_j(1-2\xi_{\eq,j})$ — 가 살아남아 평형 조건이 음함수가 되고, 등온선은 더 이상 단순 logistic 이 아니다.
그때 $w_j$ 를 effective width 로 두는 것은 닫힌 유도가 아니라 coarse-grained 근사다. 특히 $\Omega_j>2RT$ 면
isotherm 이 비단조가 되어 상분리(2-상 공존, miscibility gap)가 나타나며, 그 구간의 peak 는 단일상 logistic 으로
기술되지 않는다. 이것이 흑연 staging 의 핵심 물리이며 \S\ref{sec:regsol} 과 \S\ref{sec:hys_peak} 가 다룬다.
``Gaussian(erf) 형''으로 보는 것도 유추일 뿐 확정이 아니다. 그럼에도 본 장의 꼬리 결론 — 원거리 감쇠 길이와 부호 —
은 평형 형태의 세부에 둔감하다. 포화 후 꼬리에서는 $\xi_{\eq,j}$ 가 이미 포화해 $\dd\xi_{\eq,j}/\dd V_n\simeq0$
이기 때문이다. 엄밀히 말하면 이 둔감성은 동역학 꼬리가 평형 wing 보다 긴 꼬리-우세 조건 $L_{V,j}\gtrsim RT/F$
에서 성립한다. 그 미만이면 원거리는 진폭 $e^{-\Omega_j/RT}$ 로 억제된 평형 wing 이 지배하는데(점근 유도는
\S\ref{sec:regsol} 의 $w^\eff$ 단락), 그 경우에는 꼬리 자체가 관측되지 않으므로 결론이 위협받지는 않는다. 한편
꼬리의 시작 위치와 진폭은 peak 근처의 평형 형태에 일부 의존한다. 형태 자체의 판별은 저율 OCV(S1, \S\ref{sec:synth})
의 분리 peak 형상을 직접 읽어서 하며, \S\ref{sec:falsify} 가 판별하는 것은 형태가 아니라 율속 기전이다.
\end{flagbox}

\subsection{평형 peak 의 위치$\cdot$너비$\cdot$높이}
\textbf{도함수.} $z\equiv s(V_n-U_j)/w_j$ 로 두면 $\xi_{\eq,j}=(1+e^{-z})^{-1}$ 이다. 직접 미분하면
\begin{equation}
\frac{\dd\xi_{\eq,j}}{\dd z}=-\,(1+e^{-z})^{-2}\cdot(-e^{-z})=\frac{e^{-z}}{(1+e^{-z})^2}.
\label{eq:dxidz}
\end{equation}
한편 $1-\xi_{\eq,j}=\dfrac{e^{-z}}{1+e^{-z}}$ 이므로 $\xi_{\eq,j}(1-\xi_{\eq,j})=\dfrac{1}{1+e^{-z}}\cdot
\dfrac{e^{-z}}{1+e^{-z}}=\dfrac{e^{-z}}{(1+e^{-z})^2}$ 이고, 식~\eqref{eq:dxidz} 와 같다. 즉
$\dd\xi_{\eq,j}/\dd z=\xi_{\eq,j}(1-\xi_{\eq,j})$. 연쇄율 $\dd z/\dd V_n=s/w_j$ 를 곱하면
\begin{equation}
\frac{\dd\xi_{\eq,j}}{\dd V_n}=\frac{s\,\xi_{\eq,j}(1-\xi_{\eq,j})}{w_j}.
\label{eq:dxidV}
\end{equation}

이제 평형 도함수 식~\eqref{eq:dxidV} 를 관측식 식~\eqref{eq:dQdV} 의 peak 항에 대입한다.
\textbf{세 결과(각각 유도).} 식~\eqref{eq:dQdV} 의 peak 항 $Q_j\,\dd\xi_{\eq,j}/\dd V_n$ 으로부터:
\begin{itemize}[leftmargin=1.6em,itemsep=2pt]
\item \emph{위치}: $\dd\xi_{\eq,j}/\dd V_n\propto\xi_{\eq,j}(1-\xi_{\eq,j})$ 는 $\xi_{\eq,j}=\tfrac12$ 에서 최대다
  ($\xi(1-\xi)$ 가 $\xi=\tfrac12$ 서 최대). $\xi_{\eq,j}=\tfrac12\Leftrightarrow z=0\Leftrightarrow V_n=U_j$. 즉
  $V_\peak=U_j(T)$.
\item \emph{높이}: 그 점에서 $\xi_{\eq,j}(1-\xi_{\eq,j})=\tfrac14$ 이므로
  $\big(\dd Q/\dd V_n\big)_{\max}=C_\bg+Q_j/(4w_j)$($C_\bg$ 가 peak 폭서 거의 상수$\cdot$인접 전이 비중첩 가정), peak 순높이 $=Q_j/(4w_j)$.
\item \emph{너비}(FWHM): 아래 검산에서 $\Delta V_{1/2}=2w_j\ln(3+2\sqrt2)\approx3.53\,w_j$.
\end{itemize}
\begin{equation}
\boxed{\;V_\peak=U_j(T),\qquad \Big(\tfrac{\dd Q}{\dd V_n}\Big)^{\text{net}}_{\max}=\frac{Q_j}{4w_j},\qquad
\Delta V_{1/2}=2w_j\ln(3+2\sqrt2)\approx3.53\,w_j\;.}
\label{eq:eqpeak}
\end{equation}
\emph{면적 보존.} 배경을 뺀 peak 면적은 $\int(\dd Q/\dd V_n-C_\bg)\,\dd V_n=Q_j\int_0^1\dd\xi_{\eq,j}=Q_j$ 다.
면적이 $Q_j$ 로 고정되므로 너비가 줄면 높이가 오른다(높이$\times$너비 $\sim Q_j$). 단 $\int_0^1\dd\xi=1$ 은
$\xi_{\eq}(V_n)$ 가 단조 단일값일 때의 항등식이다. 상분리($\Omega_j>2RT$, \S\ref{sec:regsol})에서 등온선이
비단조가 되면 $\xi_{\eq}(V_n)$ 가 다가가 되어 이 적분 형태는 깨진다. 그래도 총 전하 $Q_j$ 자체는 물질량 보존으로
유지되며, 평형 기여는 plateau(Maxwell 공통접선)로 대체된다. \emph{최종식과의 연결.} 이 평형 peak 의 세 양
$\{U_j,w_j,Q_j\}$ 가 최종 피팅식(\S\ref{sec:master})의 평형 상승부 항을 이루고, 꼬리가 사라진 저율 OCV 에서 직접
닫힌다(\S\ref{sec:synth} S1). 단 높이 $Q_j/(4w_j)$ 와 FWHM $3.53\,w_j$ 의 logistic 공식은 단상($\Omega_j<2RT$)
한정이다. 상분리 staging peak($\Omega_j>2RT$)에서 $w_j$ 는 평형 폭이 아니라 broadening 이 정하는 현상학적 폭이므로
(\S\ref{sec:hys_peak}$\cdot$\S\ref{sec:master}), 그 logistic 높이$\cdot$면적 공식을 그대로 쓰지 않고 중심 $U_j$ 와
용량 $Q_j$ 만 직접 운반한다.

\begin{verifybox}
\textbf{FWHM 유도.} $g(z)=\xi_{\eq,j}(1-\xi_{\eq,j})$ 는 $z=0$ 서 최대 $\tfrac14$. 절반 $\tfrac18$ 되는 점:
$\xi(1-\xi)=\tfrac18$. $\xi=\tfrac12(1\pm\tfrac1{\sqrt2})$(이차방정식 $\xi^2-\xi+\tfrac18=0$ 의 해)이고,
$e^z=\xi/(1-\xi)=\dfrac{1+1/\sqrt2}{1-1/\sqrt2}$. 분자$\cdot$분모에 $\sqrt2$ 곱해 $\dfrac{\sqrt2+1}{\sqrt2-1}
=(\sqrt2+1)^2=3+2\sqrt2$. 따라서 $\Delta z=2\ln(3+2\sqrt2)$, $\Delta V_{1/2}=w_j\Delta z\approx3.53\,w_j$.
\end{verifybox}

\subsection{온도 의존(평형 기준선)}
\textbf{너비$\cdot$높이.} 이상 폭은 $w_j=RT/F$ 이므로 너비 $\propto T$ — \emph{저온서 좁아진다}($25^\circ$C 에서
$w=RT/F=8.314\times298/96485=25.7$ mV, $-15^\circ$C 에서 $8.314\times258/96485=22.2$ mV). 높이 $\propto1/w_j\propto
1/T$ — 저온서 높아진다.

\textbf{위치의 온도계수(Gibbs--Helmholtz 관계).} 전이 $j$ 의 \emph{삽입(리튬화) 반응} 자유에너지는 $\Delta G_j=-sF\,U_j$
다. 이는 평형 전위의 정의 그 자체이며, 부호는 삽입 방향 기준이다. 온도로 미분하면 $\partial\Delta G_j/\partial T=-\Delta S_j$(Gibbs 관계)이므로
\begin{equation}
\frac{\partial U_j}{\partial T}=-\frac{1}{sF}\frac{\partial\Delta G_j}{\partial T}=\frac{\Delta S_j}{sF}.
\label{eq:dUdT}
\end{equation}
즉 peak \emph{위치}는 삽입 반응 엔트로피 $\Delta S_j$ 에 따라 온도와 함께 이동한다 \cite{reynier2004}.
탈리튬화 기준으로 보면 $\Delta S_j\to-\Delta S_j$ 가 되지만 물리적으로는 동치다. 또한 이 $\Delta S_j$ 는 활성화
엔트로피 $\Delta S_{a,j}$ 와 다른 양임에 주의한다 — 후자는 \S\ref{sec:lag} 의 \emph{속도}에 들어간다.
수치로 가늠해 두면, 흑연 staging 의 전형적인 $|\Delta S_j|$ 는 수 J/(mol\,K) 수준이라 \cite{reynier2004}
$|\partial U_j/\partial T|=|\Delta S_j|/F\sim0.01$--$0.1$ mV/K 다. $15/45^\circ$C 의 $30$ K 창에서 peak
위치가 $0.3$--$3$ mV 움직이는 규모이며, 이것이 S1 을 각 온도에서 반복해 $U_j(T)$ 를 따로 읽는 까닭이다
(\S\ref{sec:master} walkthrough 1).

\begin{keybox}
\textbf{평형 기준선의 한계.} 평형 peak 는 대칭이고, 저온에서 더 좁고 높다. 그러나 관찰은 ``저온 $\to$ 긴 꼬리
(넓고 비대칭)''로 정반대다. 또 평형 peak 는 전류(C-rate)와 무관한데, 실제 peak 는 C-rate 에 민감하다.
\textbf{따라서 꼬리의 주된 원인은 평형 열역학이 아니라, 전류가 흐를 때 평형을 따라가지 못하는 동역학이다.}
단 그 동역학으로 가기 전에 한 가지가 남아 있다 — 평형 형태 자체가 이상 극한으로 끝나지 않는다는 점이다. 다음 절이
정규용액 상호작용($\Omega_j\ne0$)과 상분리까지 평형 기준선을 넓히고, 그 위에서 \S\ref{sec:lag} 가 동역학 지연을
다룬다.
\end{keybox}

% ====================================================================
\section{정규용액 상호작용과 상분리}\label{sec:regsol}
% ====================================================================
앞 절의 평형 기준선은 이상 극한($\Omega_j=0$)의 매끄러운 logistic 이었다. 그러나 흑연 staging 은 정규용액
상호작용 $\Omega_j\ne0$ 이 만드는 \emph{일차 상전이}여서, 동역학을 따지기도 전에 평형 형태 자체가 비단조가 될
수 있다. 본 절은 그 상분리 열역학을 세운다. 이 열역학은 두 군데서 쓰인다. 첫째, $\Omega_j>2RT$ 구간에서 평형
등온선을 비단조로 만들어 peak 의 성격을 바꾼다(본 절). 둘째, 뒤의 확장에서 \emph{충방전 히스테리시스}
(\S\ref{sec:hys_branch})의 분기를 만드는 뿌리가 된다.

\subsection{정규용액 자유에너지 — 앞 절 $\mu$ 와 같은 $g_j$}
앞 절 화학퍼텐셜 식~\eqref{eq:mu} 의 상호작용 항 $\Omega_j(1-2\theta_j)$ 는 \emph{정규용액}(regular solution)
모형에서 온다. 전이 $j$ 의 진행률 $\xi$($=1-\theta_j$)에 대한 몰당 자유에너지는 혼합 entropy 와 상호작용 enthalpy
의 합이다:
\begin{equation}
g_j(\xi)=g_j^0+\underbrace{RT[\xi\ln\xi+(1-\xi)\ln(1-\xi)]}_{-T S_\mathrm{mix}}
+\underbrace{\Omega_j\,\xi(1-\xi)}_{\text{상호작용}}.
\label{eq:regsol_g}
\end{equation}
혼합 entropy 항은 앞 절 식~\eqref{eq:smix} 의 $S_\mathrm{mix}$ 그대로이고($\xi\leftrightarrow\theta$ 대칭), 상호작용
항은 \emph{평균장}에서 ``전환$\cdot$미전환'' 이웃 쌍의 분율 $\propto\xi(1-\xi)$ 에 쌍당 초과 에너지 $\Omega_j$
($\sim z[\epsilon_{AB}-\tfrac12(\epsilon_{AA}+\epsilon_{BB})]$, $z$=이웃 수)를 곱한 것이다 — $\Omega_j>0$ 이면
\emph{동종 자리 선호}(이종 쌍 불리)라 분리를 부추긴다. 이 최근접쌍 분해는 평균장 어림이다. 실제 흑연 staging 의 상호작용은 층간 탄성 변형 같은 \emph{장거리} 성분 —
Daumas--H\'erold 도메인의 기원 — 을 포함하므로, $\Omega_j$ 는 그것까지 한 쌍-에너지로 뭉뚱그린 \emph{유효}
척도로 읽는다. 즉 $\Omega_j$ 는 부호와 규모의 안내이지 미시 정량 예측이 아니다. $g_j$ 의 점유율 미분이 곧 앞 절
식~\eqref{eq:mu} 의 $\mu_\mathrm{Li}(\theta)$ 다. 정확히는 $\partial g_j/\partial\theta=\mu_\mathrm{Li}-\mu^0$
으로 상수 $\mu^0$ 만큼 차이 난다. 요컨대 \emph{같은 $g_j$ 를 한 번은 peak(1계 미분)로, 한 번은 상분리(2계 곡률)로
본다} — 이 한 함수의 두 사영이 본 절 전체의 구도다.

\subsection{곡률과 임계 온도 — 이중 우물}
상분리는 $g_j$ 의 \emph{곡률}이 정한다:
\begin{equation}
g_j''(\xi)=\frac{RT}{\xi(1-\xi)}-2\Omega_j.
\label{eq:regsol_gpp}
\end{equation}
entropy 항 $+RT/[\xi(1-\xi)]>0$ 은 항상 볼록(섞이려 함), 상호작용 항 $-2\Omega_j$ 는 $\Omega_j>0$ 라 항상 오목
(뭉치려 함). 양 끝($\xi\to0,1$)서는 entropy 가 발산해 $g_j''>0$(볼록)이지만, 가운데서 $\Omega_j$ 가 크면 $-2\Omega_j$ 가
이겨 $g_j''<0$(오목)이다. \emph{가운데 오목 구간이 $g_j$ 를 위로 밀어 그 \textbf{양쪽} 볼록 구간에 극소가 하나씩} 생기므로
— 두 골, 곧 \emph{이중 우물}이다. 국소 $g''<0$ 한 점이 아니라 ``양끝 볼록 $+$ 가운데 오목''의 \emph{전역} 조합이 두 최소를
만든다. \emph{정확히 둘}인 까닭은 entropy 곡률 $RT/[\xi(1-\xi)]$ 가 $\xi=\tfrac12$ 서 최소인 U-형이라 오목 구간
($g''<0$)이 \emph{단 하나}뿐이기 때문이다 — 변곡점이 정확히 둘이고 $g'$ 가 증가–감소–증가의 한 굽이만 가져, 공통접선
기울기 $\mu_c$ 에 대한 $g_j'=\mu_c$ 의 해가 많아야 셋(극소 둘$+$극대 하나)이다. 두 변곡점(spinodal, $g_j''=0$,
$\xi_{s,j}^\pm$)은 극소가 \emph{아니라} 오목 구간의 경계이고, 두 극소는 그 \emph{바깥}에 놓인다. $\xi\!\leftrightarrow\!
1-\xi$ \emph{대칭}이 주는 것은 둘의 \emph{대칭 배치}($\mu_c$ 가 $\xi=\tfrac12$ 기울기)다 — 비대칭계에선 배치만
비대칭해질 뿐 ``오목 구간 하나 $\Rightarrow$ 극소 둘''은 유지된다. 한 우물과 이중 우물의 경계는 $g_j''$ 이 가장 작은 $\xi=\tfrac12$ 에서 $4RT-2\Omega_j=0$ 이 되는 곳, 곧
\textbf{임계 온도 $T_{c,j}=\Omega_j/2R$} 다. $T>T_{c,j}$(즉 $\Omega_j<2RT$)면 우물이 하나로 고용체 단상이다.
다만 단상이라 해도 $\Omega_j\ne0$ 이면 등온선은 순수 logistic 이 아니라 좁아진 peak 가 된다 — 아래 $w^\eff$ 가
그 정량이며, ``logistic 그대로''는 $\Omega_j=0$ 에 한한다. 반대로 $T<T_{c,j}$($\Omega_j>2RT$)면 이중 우물, 곧
상분리다.

\subsection{binodal $\cdot$ spinodal $\cdot$ 준안정}
이중 우물에서 한 균질 상보다 두 상($\xi$ 작은$\cdot$큰)으로 갈라지는 편이 자유에너지가 낮다 — miscibility gap.
갈림은 두 경계로 규정된다:
\begin{itemize}[leftmargin=1.4em,itemsep=2pt]
\item \textbf{binodal} $\xi_{b,j}^\pm$ — 두 상이 \emph{평형}으로 공존하는 조성. 공존 평형은 두 조건의 \emph{연립}이다:
  (i) 화학퍼텐셜 일치 $g_j'(\xi_b^-)=g_j'(\xi_b^+)$($=\mu$, 같은 기울기), (ii) 대평형(grand potential) 일치
  $g_j(\xi_b^-)-\mu\,\xi_b^-=g_j(\xi_b^+)-\mu\,\xi_b^+$(상 교환이 자유에너지를 안 바꿈 $\Leftrightarrow$ 같은 절편).
  (i)$+$(ii)가 곧 \emph{한 직선이 $g_j$ 를 두 점서 접한다}는 \emph{공통 접선} 기하이며, 그 접선 기울기가 평탄
  전위(plateau)다 — 이 binodal plateau 전위가 곧 \emph{가역 평형 중심} $U_j$ 에 대응하고, 히스 분기의 $\gamma_j\to0$
  극한(\S\ref{sec:hys_branch})이 이 binodal(가역)로 수렴한다. 본 장 정량은 spinodal$\cdot$$\Omega_j$ 로 닫고 binodal 은
  이 \emph{가역 기준}으로 쓴다.
\item \textbf{spinodal} $\xi_{s,j}^\pm$ — $g_j''=0$ 의 \emph{역학적 불안정} 경계. 식~\eqref{eq:regsol_gpp}$=0$ 을 풀면
  닫힌 꼴 $\xi_{s,j}^\pm=\tfrac12\pm\tfrac12\sqrt{1-2RT/\Omega_j}$ 다(상세 유도는 식~\eqref{eq:hys_spinodal}).
  안쪽($g_j''<0$)은 무한소 요동에도 즉시 분리(불안정), binodal--spinodal 사이($g_j''>0$ 이나 공통 접선 위)는 \emph{준안정}.
\end{itemize}
조성축이 세 영역(안정 단일 상 / 준안정 / 불안정)으로 나뉘며 \emph{준안정}이 뒤 히스테리시스의 무대다. 준안정
영역에선 새 상의 \emph{핵}이 임계 크기를 넘어야(핵생성 장벽) 전이가 시작된다 — 장벽은 구동력이 작은 binodal
근처서 크고 안쪽으로 갈수록 준다. 장벽이 크면 계가 균질 준안정 분기를 따라 \emph{spinodal 까지 과주행}한 뒤에야
(거기선 준안정 최솟값 자체가 사라져 장벽 없이) 분리한다. \textbf{충$\cdot$방전이 서로 다른 끝에서 과주행하는 것이
충방전 히스테리시스의 근원}이다(\S\ref{sec:hys_branch}).

\textbf{전위로 보면 van der Waals 등온선.} 이중 우물이면 평형 전위 $V_{\eq,j}(\xi)$ — 평형 조건~\eqref{eq:eqcond}
를 $V_n$ 에 대해 푼 것으로, 닫힌 꼴은 식~\eqref{eq:hys_Veq} 다 — 가 \emph{비단조} loop 를 그린다. 이는 압력--부피
van der Waals 등온선과 동형이다. 대응은 $g_j\leftrightarrow$ Helmholtz $A$, $\dd g_j/\dd\xi\leftrightarrow-P$,
$\xi\leftrightarrow$ 부피로 읽는다. 한 가지 부호에 주의한다 — 여기 대응에 쓰인 $\dd g_j/\dd\xi$ 는 점유율
화학퍼텐셜 $\mu_\mathrm{Li}$ 와 부호가 반대다. $\xi=1-\theta$ 이므로
$\dd g_j/\dd\xi=-\dd g_j/\dd\theta=-(\mu_\mathrm{Li}-\mu^0)$ 이고, 같은 물리를 변수만 바꿔 본 것이다.
음의 기울기 구간(spinodal 안쪽)은 vdW 에서 압축률이 음($\partial P/\partial V>0$, $\kappa_T<0$)인 불안정 가지에
대응하고, 공통 접선은 Maxwell 등면적 작도, 곧 평탄 전위에 대응한다. \emph{이 loop 의 두 어깨가 곧 충$\cdot$방전
분기가 된다}(\S\ref{sec:hys_branch}). 본 절의 상분리가 뒤의 히스테리시스 확장으로 이어지는 다리가 바로 이 loop 다.
\begin{flagbox}
\textbf{단일 입자 vs 다입자.} 위 spinodal 과주행은 \emph{단일 입자} 그림이다. 실제 전극은 입자 집단이라 입자마다
임계 핵생성이 조금씩 다른 전위에서 일어나 \emph{하나씩 차례로} 전환된다(Dreyer 다입자~\cite{dreyer2010}) —
같은 압력의 풍선 다발이 동시에 부풀지 않고 하나씩 부풀듯. 그래서 집단 거동은 단일 입자 과주행보다 분기 갈림을
\emph{매끄럽게$\cdot$좁게} 만들고, 단일 입자 spinodal 이 \emph{상한}, 실측은 그 아래(축소 인자 $\gamma_j$,
\S\ref{sec:hys_branch})에 놓인다.
\end{flagbox}

\subsection{peak 에의 영향 — 유효 폭의 좁아짐}
정규용액 상호작용은 \emph{상분리 이전에 이미} 평형 peak 의 \emph{폭}을 바꾼다. 평형 조건 식~\eqref{eq:eqcond}
$\mu_\mathrm{Li}(\theta)=\mu^0-sF(V_n-U_j)$ 를 $V_n$ 으로 미분하면 $[RT/(\theta(1-\theta))-2\Omega_j]\,\dd\theta/\dd V_n=-sF$
이고, $\xi=1-\theta$ 라 전환율의 전위 민감도가
\begin{equation}
\frac{\dd\xi_{\eq,j}}{\dd V_n}=\frac{sF}{\,RT/[\xi(1-\xi)]-2\Omega_j\,}
\label{eq:regsol_slope}
\end{equation}
이다. peak 중심($\xi=\tfrac12$)에서
\begin{equation*}
\frac{\dd\xi_{\eq,j}}{\dd V_n}\Big|_{\xi=1/2}=\frac{sF}{4RT-2\Omega_j}
\end{equation*}
이고, 이상 logistic 의 중심 기울기
$s/(4w_j)$(식~\eqref{eq:dxidV} 의 $\xi=\tfrac12$ 값)와 맞추면 \emph{유효 폭}
\begin{equation}
\boxed{\;w_j^\eff=\frac{RT}{F}\Big(1-\frac{\Omega_j}{2RT}\Big)=\frac{RT}{F}\Big(1-\frac{T_{c,j}}{T}\Big)\qquad(\Omega_j<2RT).\;}
\label{eq:weff}
\end{equation}
$w_j^\eff$ 는 peak \emph{중심 곡률}만 이상 logistic 과 맞춘 \emph{국소} 정의다. 정규용액 peak 는 logistic 이
아니므로, 중심 기울기가 같아도 전역 모양은 다르다. 두 군데서 다르다. 첫째, \emph{반치 영역}은 matched logistic
보다 좁아서 $3.53\,w_j^\eff$ 가 FWHM 을 과대 추정하며, 그 편향은 $\Omega_j\to2RT$ 로 갈수록 커진다. 둘째,
\emph{원거리 wing} 은 감쇠 길이가 $\Omega_j$ 와 무관하게 $RT/F$($>w_j^\eff$)로 고정된 채 진폭만
$e^{-\Omega_j/RT}$ 로 억제되어, matched logistic 보다 오히려 무겁다. 이 점근의 출처를 적어 두면 — wing
($\xi\to1$)에서 상호작용 항이 $\Omega_j(1-2\xi)\to-\Omega_j$ 로 포화하므로 평형 조건~\eqref{eq:eqcond} 이
$RT\ln[\xi/(1-\xi)]\simeq sF(V_n-U_j)+\Omega_j$ 로 닫히고, 따라서
$1-\xi\simeq e^{-\Omega_j/RT}e^{-sF(V_n-U_j)/RT}$($s{=}+1$ 방전 규약)이다. 감쇠 길이는 $RT/F$ 로 고정되고
진폭만 $\Omega_j$ 가 억제한다. 요컨대 $w^\eff$ 는 중심 곡률만 정확하고 전역 모양은 근사다. 즉 상호작용이 클수록($\Omega_j\to2RT$) peak 가 \emph{좁아지고},
$\Omega_j\to2RT^-$($T\to T_{c,j}^+$) 극한에서 폭이 $0$ 으로 — 그 지점이 곧 상분리(delta 형 plateau)의 시작이다(식 단서
$\Omega_j<2RT$ 와 정합: 경계는 극한값). \textbf{이것이
staging peak 가 이상 $RT/F$($25^\circ$C 서 $25.7$ mV)보다 \emph{날카로운} 까닭}이며, 동시에 \S\ref{sec:master} 의 폭
파라미터 $w_j$ 에 물리적 근거와 \emph{온도 의존}을
준다: $w_j(T)=(RT-\Omega_j/2)/F$ 는 이상의 $RT/F$ 와 달리 $T$ 에 단순 비례가 아니라 절편 $-\Omega_j/2F$ 를 가져,
\emph{한 온도의 $w_j$ 가 $\Omega_j$ 를 정하고 그것이 다른 온도의 $w_j$ 를 예측}한다 — 그리고 그 $\Omega_j$ 는
충방전 gap(\S\ref{sec:hys_branch})에서 얻는 $\Omega_j$ 와 \emph{같아야} 한다. 다만 둘은 같은 곡률 $g_j''$ 의
두 사영이므로 독립 측정이 아니라 내적 일관성 점검이다(\S\ref{sec:falsify}).

$\Omega_j>2RT$ 면 식~\eqref{eq:weff} 가 음이 되어 폭의 의미를 잃는다. 그 구간의 isotherm 은 비단조이고, peak 는
단일상 logistic 이 아니라 상분리 plateau — 매우 좁고 큰 peak — 가 된다. 이것이 staging 의 가장 날카로운 peak 다.
\textbf{이로써 평형 기준선이 이상($\Omega=0$) $\to$ 폭 좁아짐($0<\Omega<2RT$) $\to$ 임계$\cdot$상분리
($\Omega\ge2RT$)의 한 모수축 $\Omega_j/2RT$ 위에 닫혔다.} 단 $\Omega_j=2RT$ 는 매끄러운 연장이 아니라 단상과
2상이 갈리는 질적 경계, 곧 임계점임에 유의한다. 평형 기준선이 완성되었으니, 다음 절부터는 전류가 흐를 때 평형에서
\emph{이탈}하는 동역학으로 간다.
\begin{verifybox}
\textbf{Worked example — 폭에서 $\Omega_j$ 읽기.} $25^\circ$C 저율 OCV 에서 어느 단상 전이의 폭이
$w_j=18$ mV 로 읽혔다고 하자. 이상 폭은 $RT/F=25.7$ mV 이므로 식~\eqref{eq:weff} 에서
$\Omega_j=2RT(1-Fw_j/RT)=2RT\times0.30\approx1.5$ kJ/mol 이고, 임계 온도는
$T_{c,j}=\Omega_j/2R\approx90$ K 로 측정창보다 한참 아래다. 이 전이는 전 온도창에서 단상이며, gap 경로가
아니라 폭 경로 $w_j(T)=(RT-\Omega_j/2)/F$ 의 온도 예측으로 검증된다. 반대로 상온에서 상분리가 보이는 전이는
$\Omega_j>2RT\approx5.0$ kJ/mol 의 하한을 가지며, 그 $\Omega_j$ 는 폭이 아니라 충방전 gap 절편의 온도
의존(S5)에서 읽는다. 같은 $\Omega_j$ 가 두 경로를 내적 일관성으로 묶는다(\S\ref{sec:falsify}(iii)).
\end{verifybox}

% ====================================================================
\section{동역학 지연과 비대칭 꼬리}\label{sec:lag}
% ====================================================================
앞 두 절이 평형 기준선을 완성했다. 이제 전류를 흘린다. 전류가 흐를 때 전환율 $\xi_j$ 는 평형 목표
$\xi_{\eq,j}$ 를 즉시 따라가지 못하며, 그 \emph{지연}이 peak 를 비대칭으로 늘린다. 이 절은 지연의 운동방정식을
세우고 그 \emph{일반해}를 먼저 구한 뒤, 상승부와 꼬리를 그 일반해의 두 극한으로 읽는다. 그 과정에서 ``왜
저온$\cdot$고율에서 꼬리가 길어지는가''가 닫힌 식으로 답해진다.

\subsection{완화 방정식의 유도}
\textbf{1차 완화의 근거.} 전이 $j$ 를 정/역 두 방향 속도 $r_j^+,r_j^-$ 를 갖는 2-상태 과정으로 보면, \emph{진행률}의
순 변화는 \emph{전향}(전이 완성=탈리튬화, $\xi_j$ 증가; $\propto$ 아직 전이 안 된 분율 $1-\xi_j=\theta_j$, 즉 잔여 Li
점유율)과 \emph{역향}(재리튬화, $\xi_j$ 감소; $\propto$ 이미 전이된(빈자리) 분율 $\xi_j$)의 차다:
\begin{align}
\frac{\dd\xi_j}{\dd t}&=r_j^+(1-\xi_j)-r_j^-\xi_j
=r_j^+-(r_j^++r_j^-)\xi_j\notag\\
&=(r_j^++r_j^-)\Big[\frac{r_j^+}{r_j^++r_j^-}-\xi_j\Big]
\equiv k_j\,(\xi_{\sstat,j}-\xi_j),
\label{eq:relax}
\end{align}
$k_j\equiv r_j^++r_j^-$(완화 속도), $\xi_{\sstat,j}\equiv r_j^+/(r_j^++r_j^-)$(정상점). 유일한 가정은 ``$r_j^\pm$ 를
$\xi_j$ 의 작은 변화에 대해 국소 상수로 본다''이며, 그 위에서 위 재배열은 \emph{근사가 아니라 대수적 항등}이다.
\emph{단, 전이 전체를 \textbf{단일 완화 모드}(1차 ODE)로 보는 것 자체가 ansatz 다} — 핵생성$\cdot$상경계 이동 등
여러 시간척도를 하나의 유효 모드로 coarse-grain 한 것으로, 한 완화 시간이 지배할 때 타당하다(다중 모드$\cdot$입자
분산 효과는 \S\ref{sec:dist}).
\textbf{정상점과 평형값의 일치}($\xi_{\sstat,j}=\xi_{\eq,j}$)는 detailed balance 의 귀결인데, 그 검증
($r_j^+/r_j^-=e^{\mathcal A_j/RT}$ 가 식~\eqref{eq:logistic} 을 재생함)은 \S\ref{sec:barrier} 에서 정/역 속도를
구체화할 때 이루어진다 — 여기서는 $\xi_{\sstat,j}=\xi_{\eq,j}$ 로 둔다(전방 참조).

\textbf{속도상수와 활성화 배리어 (Eyring).} 전이상태 이론에서 속도상수는, 전이상태를 통과하는 보편 빈도
$k_BT/h$ 에 장벽을 넘을 Boltzmann 인자를 곱한 것이다 \cite{eyring1935}:
\begin{equation}
k_j\simeq k_0\,\exp\!\Big[-\frac{\Delta G_{a,j}}{RT}\Big]\ \ (k_0=k_BT/h,\ \text{현상론적 prefactor}),\qquad \Delta G_{a,j}=\Delta H_{a,j}-T\Delta S_{a,j}.
\label{eq:eyring}
\end{equation}
여기 $k_j$ 는 \emph{완화 속도}($=r_j^++r_j^-$, 식~\eqref{eq:relax}) \emph{그대로}이며, 그 \emph{척도}가 활성화 배리어의
Arrhenius/Eyring 형임을 적은 것이다. 그래서 등호가 아니라 ``$\simeq$''를 썼으며, $k_0$ 는 $O(1)$ 수
인자$\cdot$교환전류$\cdot$활동도를 흡수하는 현상론 prefactor 다(\S\ref{sec:barrier}). 정/역 분해는
\S\ref{sec:barrier} 의 $r_j^\pm$ 로 구체화되며, forward 우세 꼬리에서 $k_j\simeq r_j^+$ 가 된다. peak 중심부
($\mathcal A_j\approx0$)에서는 $r_j^+\approx r_j^-$ 라 합이 인자 $2$ 만큼 크지만, 그 수 인자도 $k_0$ 에
흡수된다. $k_j$ 를 forward 속도로 재정의하지는 않는다.
\textbf{저온에서 $k_j$ 가 작아진다} — $e^{-\Delta H_{a,j}/RT}$ 가 급감하기 때문이며, 이것이 곧 ``저온 긴
꼬리''의 근원이다. 꼬리 길이의 절대값은 $L_{q,j}\propto1/k_0$ 라 $k_0$ 에 의존한다. 그러나 본 장의 결론 — 꼬리의
온도$\cdot$전위 \emph{의존성}과 그 부호, Arrhenius 기울기로 뽑는 $\Delta H_a$ — 은 $k_0$ 를 피팅 prefactor 로
흡수하므로 그 값에 무관하다. 이 흡수에는 전제가 하나 있다. 흡수된 인자들이 피팅 온도창에서 $T$-무관이어야 하며,
강한 $T$ 의존이 있으면 Arrhenius 절편과 기울기에 계통오차로 들어간다.

\subsection{지연이 만드는 비대칭 꼬리}
\textbf{측정축으로의 변환(시간 $\to$ 용량).} 정전류 방전에서 무차원 용량 좌표 $q=\int|I|\,\dd t/Q_\cell$ 는
$\dd q/\dd t=|I|/Q_\cell$(단위시간당 흘린 용량분율)이다. 연쇄율 $\dd\xi_j/\dd q=(\dd\xi_j/\dd t)/(\dd q/\dd t)$ 로
식~\eqref{eq:relax} 을 $q$ 미분으로 바꾸면
\begin{equation}
\frac{\dd\xi_j}{\dd q}=\frac{k_j(\xi_{\eq,j}-\xi_j)}{|I|/Q_\cell}=\frac{Q_\cell\,k_j}{|I|}\,(\xi_{\eq,j}-\xi_j).
\label{eq:dxidq}
\end{equation}

\textbf{지연의 운동방정식 — 보편 형태를 먼저 세운다.} 지연을 $r_j\equiv\xi_{\eq,j}-\xi_j$ 라 하자. 첨자 없는
$r_j$ 는 지연이며, 정/역 \emph{속도} $r_j^\pm$ 와 구별한다. 지연의 변화율은 목표의 변화율에서 실제 진행률의
변화율을 뺀 것이므로, $\dd r_j/\dd q=\dd\xi_{\eq,j}/\dd q-\dd\xi_j/\dd q$ 에 식~\eqref{eq:dxidq} 의
$\dd\xi_j/\dd q=(Q_\cell k_j/|I|)\,r_j$ 를 넣으면, 영역 제한 없이 성립하는 보편 방정식을 얻는다:
\begin{equation}
\frac{\dd r_j}{\dd q}=\frac{\dd\xi_{\eq,j}}{\dd q}-\frac{1}{L_{q,j}}\,r_j,\qquad
\boxed{\;L_{q,j}=\frac{|I|}{Q_\cell\,k_j}\;}.
\label{eq:Lq}
\end{equation}
첫째 항은 평형 목표의 이동이 지연을 \emph{공급}하는 원천이고, 둘째 항은 완화가 지연을 \emph{소거}하는 감쇠다.
감쇠의 특성 길이 $L_{q,j}$ 가 본 장 꼬리의 핵심 양이다. 여기 $k_j=r_j^++r_j^-$ 인데, 꼬리 영역
($\mathcal A_j\gtrsim RT$)에서는 forward 우세라 $k_j\simeq r_j^+\equiv k_j^{\mathrm{fwd}}$
(식~\eqref{eq:keff})로 읽는다. 즉 본 식과 아래 모든 꼬리 표현의 $k_j$ 는 이 forward 속도 $k_j^{\mathrm{fwd}}$ 다.
이는 \S\ref{sec:barrier} 의 \emph{직렬 율속} $k_j^\eff$ 와 다른 양으로, $k_j^\eff$ 는
$k_{j,\mathrm{act}}\simeq k_j^{\mathrm{fwd}}$ 를 공율속과 직렬 합성한 것이다. peak 중심에서는
$k_j\simeq2r_j^+$ 라 인자 $2$ 의 차이가 있고, 이 차이는 $L_q$ 의 절대값을 정확히 $2$배 바꾼다. 그러나 그 인자는
피팅 prefactor $k_0$ 에 흡수되므로, 본 장이 데이터에서 뽑으려는 양 — 꼬리의 온도$\cdot$전위 의존성과 Arrhenius
기울기 — 에는 영향이 없다. ``무영향''은 절대값이 아니라 의존성에 대한 말이다.

\textbf{일반해 — 지수 기억 커널.} 식~\eqref{eq:Lq} 는 1계 선형 ODE 이므로 적분인자 $e^{q/L_{q,j}}$ 로 닫힌
일반해를 갖는다. $L_{q,j}$ 가 적분 구간에서 천천히 변한다는 가정(slowly-varying — 정량 조건은 아래 검산 박스)
아래, 임의 시작점 $q_0$ 에 대해
\begin{equation*}
r_j(q)=r_j(q_0)\,e^{-(q-q_0)/L_{q,j}}
\;+\;\int_{q_0}^{q}e^{-(q-q')/L_{q,j}}\,\frac{\dd\xi_{\eq,j}}{\dd q'}\,\dd q'.
\end{equation*}
이 일반해가 본 절의 중심 결과다. 지연은 두 몫의 합이다. 첫째 몫은 초기 지연의 지수 감쇠(동차해)이고, 둘째 몫은
과거의 평형 목표 변화율을 지수 커널 $e^{-\Delta q/L_{q,j}}$ 로 가중해 누적한 \emph{기억}(특수해)이다. 즉 계는
용량축에서 길이 $L_{q,j}$ 만큼의 ``최근 과거''만 기억한다. 이제 이 보편식에서 두 극한을 뽑는다 — 상승부와 꼬리는
서로 다른 가정이 아니라 \emph{같은 일반해의 두 코너 케이스}다. 기억 길이의 수치 감각도 들어 두자.
$0.1$C($|I|/Q_\cell=0.1\ \mathrm h^{-1}$)에서 $k_j=10^{-3}\ \mathrm s^{-1}$ 이면
$L_{q,j}=|I|/(Q_\cell k_j)\approx0.028$ — 전체 방전의 약 $3\%$ 거리만 기억한다. 같은 율에서 저온으로
$k_j$ 가 $12$배 느려지면(아래 keybox 수치) $L_{q,j}\approx0.33$ 으로, 기억이 전이 폭을 넘어 다음 peak
까지 뻗는다. 꼬리 융합(\S\ref{sec:overlap})의 정량이 이 한 수에서 나온다.

\emph{극한 (i) — 상승부(quasi-static).} 평형 목표가 기억 길이에 비해 천천히 변하는 영역, 곧
$\dd\xi_{\eq,j}/\dd q$ 가 $L_{q,j}$ 척도에서 거의 일정한 영역에서는, 적분 안의 $\dd\xi_{\eq,j}/\dd q'$ 를
현재값으로 근사해 적분 밖으로 꺼낼 수 있다. 초기 감쇠가 끝난 뒤($q-q_0\gg L_{q,j}$) 남는 것은
\begin{equation*}
r_j(q)\;\simeq\;L_{q,j}\,\frac{\dd\xi_{\eq,j}}{\dd q}\Big|_{q}\;+\;O\!\big(L_{q,j}^2\big),
\end{equation*}
즉 지연은 목표 변화율에 비례하는 \emph{소량}이다. 빠른 동역학 극한($k_j\to\infty$, $L_{q,j}\to0$)에서 지연이
$0$ 으로 사라져 $\xi_j\to\xi_{\eq,j}$ — \S\ref{sec:eqpeak} 의 평형 추종이 회복된다. 상승부의 관측이 평형
peak 로 지배되는 까닭이 이것이다.

\emph{극한 (ii) — post-peak 꼬리.} peak 를 지나면 평형 목표가 logistic 양극단으로 포화해
$\xi_{\eq,j}(1-\xi_{\eq,j})\to0$, 따라서 원천항 $\dd\xi_{\eq,j}/\dd q\simeq0$ 이 된다. 일반해에서 적분
몫이 죽고 동차해만 남는다. 적분 하한이 될 \emph{꼬리 시작점} $q_a$ 를 정의하자: 평형 항이 포화해
$\dd\xi_{\eq,j}/\dd q$ 가 정점값의 일정 분율(예 $\sim5\%$) 이하로 떨어지는 \emph{단일} 좌표 — 상승부에서
꼬리로 넘어가는 곳이다(정확한 컷은 피팅에서 선택, \S\ref{sec:synth} S1). 그러면
\begin{equation}
r_j(q)=r_j(q_a)\,\exp\!\Big[-\frac{q-q_a}{L_{q,j}}\Big],
\label{eq:rsol}
\end{equation}
즉 지연이 지수로 감쇠한다. \textbf{$L_{q,j}$ 는 ``지연이 $1/e$ 로 줄어드는 방전 거리''}, 곧 꼬리의 특성 길이다.
관측 peak 항은 $\dd\xi_j/\dd q=-\dd r_j/\dd q=(r_j(q_a)/L_{q,j})\exp[-(q-q_a)/L_{q,j}]$ 로 같은 지수형이다.
요약하면 — 보편 방정식~\eqref{eq:Lq} 하나가 상승부의 ``평형 $+$ 소량 지연''(극한 i)과 꼬리의 ``순수 지수
감쇠''(극한 ii)를 모두 담으며, 단일지수 꼬리식~\eqref{eq:rsol} 은 전 구간 가정이 아니라 포화 영역의 극한임이
유도 순서에서 분명해진다. \S\ref{sec:synth} 의 닫힌식이 이 두 극한을 한 곡선으로 합칠 때 이 구분을 그대로
물려받는다.

\textbf{용량축 $\to$ 전압축.} post-peak 에서 $V_n$ 이 $q$ 에 거의 선형이면($\dd V_n/\dd q$ 가 $q_a$ 둘레 상수;
$V_a\equiv V_n(q_a)$), $q-q_a\simeq(V_n-V_a)/(\dd V_n/\dd q)$ 이고, 꼬리 진행 방향에서 $(V_n-V_a)$ 와 $\dd V_n/\dd q$
는 \emph{같은 부호}라 그 비 $=|V_n-V_a|/|\dd V_n/\dd q|\ge0$ 이다. (staging plateau 처럼 $\dd V_n/\dd q\to0$ 인 2-상
공존 구간은 $L_{V,j}\to0$ 이라 단일-phase 가정 밖이다 — 비고$\cdot$\S\ref{sec:falsify}; 본 꼬리식은 plateau 를
벗어난 단상 영역에 쓴다.) 따라서 꼬리는 전압축에서 감쇠한다(\emph{단방향}
— post-peak 전향 영역 $s(V_n-V_a)\ge0$ 에서만 정의되며, 양방향 대칭 감쇠가 아니다):
\begin{equation}
\frac{\dd Q}{\dd V_n}\Big|_{\text{tail}}\propto\exp\!\Big[-\frac{|V_n-V_a|}{L_{V,j}}\Big]\quad\big(s(V_n-V_a)\ge0\big),\qquad
L_{V,j}=\Big|\frac{\dd V_n}{\dd q}\Big|_{q_a}L_{q,j}.
\label{eq:tail}
\end{equation}

\begin{verifybox}
\textbf{$L_{q,j}$ 무차원.} $|I|$=$\mathrm{C/s}$, $Q_\cell$=$\mathrm C$, $k_j$=$\mathrm{s^{-1}}$ 이므로
$|I|/(Q_\cell k_j)=(\mathrm{C/s})/(\mathrm C\cdot\mathrm{s^{-1}})=1$ — $q$ 와 같은 무차원 길이. \textbf{$L_{q,j}$ 상수는
\emph{국소 근사}다}: 정확한 지연 해는
\[ r_j(q)=r_j(q_a)\exp\!\Big[-\int_{q_a}^q\frac{\dd q'}{L_{q,j}(q')}\Big] \]
이고, 꼬리 감쇠 길이 안에서 구동력 $\mathcal A_j$(따라서 $k_j$)가 천천히 변해 $|\dd L_{q,j}/\dd q|\ll1$ 일 때만 $L_{q,j}$ 를
그 구간 상수로 보고 적분 밖으로 뺀다(slowly-varying). 이 조건을 구동력으로 구체화하면 — $V_\drive=V_n$ 이라 $\mathcal A_j$
가 꼬리 좌표 자체를 따라 변하므로, 한 감쇠 길이 $L_{V,j}$ 동안 $\ln L_q$ 의 변화량
$|\partial\ln L_{q,j}/\partial V_n|\,L_{V,j}=(\chi F/RT)L_{V,j}\ll1$(식~\eqref{eq:LqV}) — 이어야 한다. 전형값($\chi=0.5$,
$L_{V,j}\sim5$ mV, $25^\circ$C)으로 $\approx0.1$ 이라 성립하나 \emph{경계적}이다: $\chi$ 가 크거나 꼬리가 길면 위반되고,
그때 꼬리는 위치마다 감쇠율이 변하는 형으로 휘어 \S\ref{sec:dist} 의 중첩으로 넘어간다.
\end{verifybox}

\subsection{온도$\cdot$C-rate 의존 — 관찰과의 일치}
식~\eqref{eq:Lq},~\eqref{eq:eyring} 에서 $1/k_j\simeq(h/k_BT)e^{\Delta G_{a,j}/RT}$ 이므로($k_0=k_BT/h$ 현상론)
\begin{equation}
L_{q,j}\simeq\frac{|I|}{Q_\cell}\frac{h}{k_BT}\exp\!\Big[\frac{\Delta G_{a,j}}{RT}\Big]
\;\Longrightarrow\;
\begin{cases}
\text{저온}\ (T\downarrow): & k_j\downarrow\Rightarrow L_{q,j}\uparrow\Rightarrow\text{꼬리 길어짐}\ (\checkmark),\\[2pt]
\text{고율}\ (|I|\uparrow): & L_{q,j}\uparrow\Rightarrow\text{꼬리 길어짐, peak 낮아짐}.
\end{cases}
\label{eq:tailTC}
\end{equation}
(``peak 낮아짐''은 직접 효과가 아니라 \emph{면적 보존 경유}다 — 꼬리로 미뤄지는 용량만큼 중심이 깎인다;
꼬리 우세 시의 한계는 \S\ref{sec:synth} 3$\times$3 표의 단서 참조.)

\textbf{단일지수에서 늘어진(stretched) 꼬리로 — 후반부의 동기.} 위 꼬리는 \emph{한 입자$\cdot$한 배리어}
$\Delta G_{a,j}$ 가정의 \emph{단일 지수}다. 그러나 실제 전극의 꼬리는 그보다 \emph{늘어져}(stretched) 관측된다 —
입자마다 크기$\cdot$결정성$\cdot$국소 staging 이 달라 활성화 배리어가 \emph{분포}를 이루고, 서로 다른 감쇠 길이
$L_{q,j}$ 가 겹치기 때문이다. 이 늘어짐을 정량화하는 것이 후반부(\S\ref{sec:stattools}$\cdot$\S\ref{sec:dist})의
과제이며, 여기서 얻은 지수 꼬리 항은 최종식(\S\ref{sec:master})의 \emph{동역학 항}이 된다(다율$\cdot$다온도로
닫힌다 — \S\ref{sec:synth} S3$\cdot$S4).

\begin{keybox}
\textbf{동역학 지연의 요지.} 동역학 지연은 평형 peak 의 뒤쪽에 지수 꼬리를 붙여 peak 를 비대칭으로 만들고,
면적 보존에 의해 높이를 낮춘다. 수치로 감을 잡아 두면 — $\Delta H_{a,j}\approx40$ kJ/mol 일 때
$25^\circ$C$\to-15^\circ$C 에서 꼬리 길이 $L_{q,j}$ 는 약 $10$배 이상 길어진다.
$e^{(\Delta H_a/R)(1/258-1/298)}\approx e^{2.5}\approx12$ 에 prefactor 의 $1/T$ 까지 곱해지기 때문이다.
빠른 kinetics 극한($k_j\to\infty$)에서는 $L_{q,j}\to0$ 이라 꼬리가 사라지고, 일반해의 극한 (i) 그대로 평형
peak(\S\ref{sec:eqpeak})로 매끄럽게 복귀한다. 온도와 C-rate 의존성은 위 식들로 설명되었다. 남은 것은 전위가
$k_j$ 를 바꾸는 방식이며, 다음 절이 그것을 다룬다.
\end{keybox}

% ====================================================================
\section{전위에 의한 유효 배리어}\label{sec:barrier}
% ====================================================================
앞 절은 완화 속도 $k_j$ 의 온도 의존을 Eyring 배리어로 닫았다. 그러나 관찰의 핵심은 ``온도에 의한 배리어
\emph{외에}, 극판 전위가 배리어를 추가로 낮춘다''는 것이었다. 그렇다면 전위는 식~\eqref{eq:eyring} 의
$\Delta G_{a,j}$ 에 어떻게 들어가는가. 이 절은 구동력을 정의하고, 정/역 속도의 보편형을 세운 뒤, 꼬리 영역의
forward 극한을 그로부터 뽑는다.

\subsection{구동력}
\S\ref{sec:stage} 에서 구동 전위 $V_\drive$(기본형 $V_\drive=V_n$)를 두었다. 전이 $j$ 의 \textbf{구동력}
(affinity)은 구동 전위가 전이 중심에서 벗어난 정도이며, 식~\eqref{eq:eqcond} 의 전기화학 평형 조건
$\mu_\mathrm{Li}=\mu^0-sF(V_n-U_j)$ 에서 \emph{독립적으로} 오는 몰당 자유에너지다. logistic 에서 빌려 온 양이
아니다. 기본형 $V_\drive=V_n$ 에서 $\mathcal A_j/RT=s(V_n-U_j)/w_j$ 가 식~\eqref{eq:logistic} 의 logit 과
같아지는데, 이는 순환 논법이 아니라 평형 열역학과 동역학 detailed balance 라는 두 독립 경로가 같은 logistic 을
내는 정합 검증이다:
\begin{equation}
\mathcal A_j=sF\,(V_\drive-U_j).
\label{eq:affinity}
\end{equation}
\emph{왜 한 항뿐인가} — $\mu_\mathrm{Li}$ 차 전체가 아니라 전기 항만 들어가는 까닭은, 완화식~\eqref{eq:relax} 의
점유 prefactor $(1-\xi_j)\cdot\xi_j$ 가 혼합 엔트로피의 통계를 \emph{이미} 담당하기 때문이다(친화도에
$RT\ln[\xi/(1-\xi)]$ 를 또 넣으면 이중 계상; $\Omega_j\ne0$ 의 상호작용 몫만 추가 대상 — 식~\eqref{eq:sseq} 아래
일반화 주).

\subsection{유효 배리어 — Butler--Volmer 방향성 장벽 이동}
\textbf{구동력에 의한 forward 장벽 강하.} 정/역 두 방향(\S\ref{sec:lag} 의 $r_j^\pm$)을 보면, 구동력 $\mathcal A_j$ 는
두 방향의 활성화 장벽을 \emph{비대칭}으로 옮긴다 — 표준 Butler--Volmer/Brønsted--Evans--Polanyi 선형 자유에너지
관계다 \cite{bardfaulkner,evanspolanyi1938}. 전달 계수 $\chi\in[0,1]$(구동력 중 forward 장벽을 낮추는 몫; 반응
좌표상 전이상태의 위치로 해석)로 — 정$\cdot$역 반응은 \emph{같은 전이상태}를 공유하므로 전이상태 통과 빈도
prefactor $k_0=k_BT/h$(식~\eqref{eq:eyring})가 두 방향에 \emph{공통}이다 —
\begin{equation}
r_j^+=k_0\,e^{-(\Delta G_{a,j}-\chi\mathcal A_j)/RT},\qquad
r_j^-=k_0\,e^{-(\Delta G_{a,j}+(1-\chi)\mathcal A_j)/RT},\qquad k_0=\frac{k_BT}{h}.
\label{eq:bv}
\end{equation}

\textbf{전달계수 분할과 합 1의 근거.} 전달 계수 $\chi$ 는 전이상태의 \emph{반응 좌표상 분율
위치}다($0$=반응물 쪽, $1$=생성물 쪽, $\tfrac12$=가운데). 구동력 $\mathcal A_j$ 가 두 우물의 상대 높이를
기울이면, \emph{같은 한 점}인 전이상태가 그 분율만큼 따라 움직여 — forward 장벽(반응물$\to$전이상태)은
$\chi\mathcal A_j$ 만큼 내려가고, backward 장벽(생성물$\to$전이상태)은 나머지 $(1-\chi)\mathcal A_j$ 만큼 올라간다
\cite{evanspolanyi1938,bardfaulkner}. \textbf{두 계수의 합이 정확히 $1$ 인 것은 자유 선택이 아니라 detailed
balance 의 강제}다: 아래에서 정/역 비 $r_j^+/r_j^-=e^{[\chi+(1-\chi)]\mathcal A_j/RT}$ 가 자유에너지 차의
Boltzmann 인자 $e^{\mathcal A_j/RT}$ 와 같으려면 지수의 두 몫의 합이 반드시 $1$ 이어야 한다. 합이 $1$ 에서
벗어나면 운동학이 주는 평형 상수가 열역학 구동력과 어긋나 평형이 깨진다. 즉 $\chi$ 자체는 $0\!\sim\!1$ 의 자유
파라미터지만, ``$\chi$ 와 $1-\chi$ 로 갈라져 합이 $1$''은 두 방향이 \emph{같은 전이상태를 공유}함의 대수적
귀결이다.

\textbf{평형 불변 (detailed balance).} 위 합$=1$ 로 두 방향 비는
\begin{equation}
\frac{r_j^+}{r_j^-}=\exp\!\Big[\frac{(\chi\mathcal A_j)-(-(1-\chi)\mathcal A_j)}{RT}\Big]
=\exp\!\Big[\frac{\mathcal A_j}{RT}\Big],
\label{eq:db}
\end{equation}
($\chi$ 가 상쇄). 따라서 정상점은 $\xi_{\sstat,j}=\dfrac{r_j^+}{r_j^++r_j^-}=\dfrac{1}{1+r_j^-/r_j^+}
=\dfrac{1}{1+e^{-\mathcal A_j/RT}}$. 기본형 $V_\drive=V_n$ 에서 $\mathcal A_j=sF(V_n-U_j)$, $w_j=RT/F$ 이므로 이는
\begin{equation}
\xi_{\sstat,j}=\frac{1}{1+\exp[-s(V_n-U_j)/w_j]}=\xi_{\eq,j}
\label{eq:sseq}
\end{equation}
즉 \S\ref{sec:eqpeak} 의 logistic~\eqref{eq:logistic} 과 \emph{정확히 일치}한다 — \S\ref{sec:lag} 에서 전방참조로
둔 $\xi_{\sstat,j}=\xi_{\eq,j}$ 가 여기서 \emph{증명}된다. 단 이 닫힌 증명은 \emph{이상 극한}($\Omega_j=0$,
$w_j=RT/F$) 위의 것이다 — $\Omega_j\ne0$ 이면 친화도를 $\mathcal A_j=sF(V_\drive-U_j)-\Omega_j(1-2\xi_j)$ 로(상호작용
몫만 추가 — 혼합 엔트로피 항을 넣지 않는 까닭은 식~\eqref{eq:affinity} 직후의 이중 계상 논거 그대로) 일반화하면 같은 논법이
$\Omega_j\ne0$ 등온선을 정확히 재생하며, 본 장 꼬리 결론(forward 우세 감쇠)은 이 세부에 둔감하다. \textbf{장벽 이동은 평형 target 을 옮기지 않는다}: 평형은 두
상태의 자유에너지 차($\mathcal A_j$)가 정하고, 장벽(활성화)은 거기 도달하는 속도만 정한다. 이것이 전이상태
이론의 표준 역할 분리이며, $\chi$ 가 식~\eqref{eq:db} 에서 상쇄되는 것이 그 표현이다.

\textbf{완화 속도의 보편형, 그리고 forward 극한.} 완화 속도는 정/역 속도의 합 $k_j=r_j^++r_j^-$ 다.
식~\eqref{eq:bv} 의 두 속도와 detailed balance 식~\eqref{eq:db} 를 쓰면, 구동력의 크기에 제한 없이 성립하는
보편형이 바로 닫힌다:
\begin{equation*}
k_j=r_j^+\big(1+e^{-\mathcal A_j/RT}\big)
=k_0\,e^{-(\Delta G_{a,j}-\chi\mathcal A_j)/RT}\big(1+e^{-\mathcal A_j/RT}\big).
\end{equation*}
괄호 인자 $(1+e^{-\mathcal A_j/RT})$ 가 역방향 몫의 전부다. 이제 이 보편형에서 꼬리 영역의 극한을 뽑는다.
방전(탈리튬화)에서 전이가 완성되는 꼬리 영역은 구동 전위가 전이 중심을 forward 방향으로 지난 곳이다 — $\xi_j$ 가
$1$ 로 갈 때 $V_n>U_j$ 이므로 $\mathcal A_j=sF(V_n-U_j)>0$($s=+1$)이다. 그러면 역방향 몫은 $e^{-\mathcal A_j/RT}$
로 억제되며, 그 크기는 $\mathcal A_j=RT$ 에서 약 $37\%$, $3RT$ 에서 약 $5\%$ 다. 따라서 구동력이 수 $RT$ 를 넘는
곳($\mathcal A_j\gtrsim3RT$, 보정 $\lesssim5\%$)에서는 역방향을 버린 forward 극한이 정확하다:
\begin{equation}
\boxed{\;\Delta G_{\eff,j}=\Delta G_{a,j}-\chi\,\mathcal A_j\;,\qquad
k_j\simeq k_0\exp\!\Big[-\frac{\Delta G_{a,j}-\chi\mathcal A_j}{RT}\Big]\quad(\mathcal A_j\gtrsim 3RT)\;.}
\label{eq:keff}
\end{equation}
구동력이 클수록($\mathcal A_j>0$) forward 유효 배리어가 낮아져 완화 속도 $k_j$ 가 커진다. 평형 target 은
\S\ref{sec:eqpeak} 그대로 남는다. \textbf{이것이 ``전위가 배리어를 낮춘다''의 정확한 자리다 — 전위가 바꾸는
것은 평형이 아니라 속도다.}
\emph{유효 작동 창 — 보편형이 가르는 세 구간.} 보편형의 괄호 인자를 기준으로 작동 창이 셋으로 갈린다.
(i) $\mathcal A_j\approx0$ 인 peak \emph{중심부}에서는 $r_j^+\approx r_j^-$ 라 $k_j\simeq2r_j^+$ 가 되어
단일지수 형이 깨지지만, 그 구간의 관측은 평형 상승부(\S\ref{sec:eqpeak})가 지배하므로 무관하다.
(ii) $\mathcal A_j\gtrsim3RT$ 의 \emph{원거리 꼬리}에서는 단일지수가 정확하다(보정 $\lesssim5\%$,
식~\eqref{eq:keff} box). (iii) 그 사이의 \emph{전이대} $RT\lesssim\mathcal A_j\lesssim3RT$ 에서는 단일지수
($k_j\simeq r_j^+$)가 선도 근사일 뿐, 보편형의 보정 인자 $(1+e^{-\mathcal A_j/RT})$ 가 최대 $\sim37\%$
($\mathcal A_j=RT$)까지 살아 있다. \textbf{이 보정은 $\mathcal A_j$ 를 통해 $V_n$ 만의 함수라
\emph{집단 공통}이므로 \S\ref{sec:dist} 의 $\rho_G$($G$-분산) 중첩으로는 생성되지 않는다}(뒤 \S\ref{sec:dist} 의
분포 적분은 배리어 $G$ 에\emph{만} 걸리므로, $V_n$ 만의 함수인 이 인자는 좁은$\cdot$넓은 분포 어느 쪽에서도 적분
\emph{밖}으로 빠져나가는 공통 인자다) — 즉 \emph{$\rho_G$ 가 흡수하지 못한다}. 정직한 처리는 식~\eqref{eq:Lq} 의 $k_j$ 를 forward 근사 $r_j^+$ 가 아니라 위 \emph{보편형 그대로} 쓰는 것이며,
그러면 보정이 $L_q$ 에 자동 포함된다. 단일지수 $r_j^+$ 형은 그 인자를 피팅 prefactor $k_0$ 에 잠근, $\gtrsim3RT$
한정의 극한 근사다. 따라서 ``꼬리식을 $\mathcal A_j\gtrsim RT$ 에 쓴다''는 말은 선도형이 그렇다는 뜻이고, 정확성은
$\gtrsim3RT$ 에 한한다. 파라미터의 식별은 \S\ref{sec:synth} S1--S4 에서 다룬다 — $\{\Delta G_{a,j},\chi,k_0\}$
는 피팅으로, $V_\drive=V_n$ 은 측정으로, $U_j$ 는 열역학 입력으로 닫힌다.
\begin{boundbox}
\textbf{선형의 한계(Marcus 유추)와 작동 창.} $-\chi\mathcal A_j$ 는 \emph{소구동력 1차} 근사다. Marcus 포물면
$\Delta G^\ddagger(\mathcal A)=\dfrac{(\lambda-\mathcal A)^2}{4\lambda}=\dfrac{\lambda}{4}-\dfrac{\mathcal A}{2}
+\dfrac{\mathcal A^2}{4\lambda}$ 의 $\mathcal A=0$ 1차 전개는 계수 $\partial\Delta G^\ddagger/\partial\mathcal A|_0
=-\tfrac12$, 즉 대칭 극한 $\chi=\tfrac12$ 을 \emph{유추}로 준다 \cite{marcus1956}. $|\mathcal A_j|\sim\lambda$(큰
과전위$\cdot$역전 영역)에서는 2차 곡률로 깨진다. \emph{intercalation 상전이로의 1차 근거는 BV/BEP 경험식이고,
Marcus 는 그 선형 형태의 물리적 유추로만 인용한다.} \textbf{본 장의 유효 작동 창은 $3RT\lesssim\mathcal A_j\ll\lambda$}
— forward 우세로 단일지수 꼬리가 정확하고($\mathcal A_j\gtrsim3RT$, 위), $\chi$ 의 $\mathcal A_j$-상수 근사 상대오차가 Marcus
2차항 $\mathcal A_j^2/4\lambda$ 대 $\chi\mathcal A_j$ 비 $\sim\mathcal A_j/(4\chi\lambda)=\mathcal A_j/2\lambda$($\chi=\tfrac12$)로 작다($\mathcal A_j=3RT,\ \lambda\gg RT$ 면 $\ll1$).
$\lambda\gg RT$ 라 이 창($3RT\sim\lambda$)이 존재한다.
\end{boundbox}
\begin{boundbox}
\textbf{속도 제한(직렬 시간척도; 불연속 절단을 피하는 구현).} 전하전달이 빨라도 수송$\cdot$유한 자리$\cdot$
고체 확산이 전체 진행을 제한할 수 있다. 임의 $\min$/clip 대신 율속 단계의 직렬 합 $1/k_j^{\eff}=1/k_{j,\mathrm{act}}+1/k_{j,\lim}$
으로 둔다($k_{j,\mathrm{act}}$=식~\eqref{eq:keff} 의 활성화 속도; $k_{j,\lim}\to\infty$ 면 $k_j^{\eff}\to k_{j,\mathrm{act}}$ 로 활성화
지배 환원). 이로써 강구동서 $\Delta G_{\eff,j}<0$ 가 되어도 $\max(\cdot,0)$ 같은 절단을 쓰지 않는다. 이하 꼬리식($L_{q,j}$)의
속도는 이 $k_j^{\eff}$ 다. 활성화 지배면 $k_j^{\eff}\simeq k_{j,\mathrm{act}}$ 이고, $k_{j,\lim}$ 의
$V_\drive$/SOC 의존 가능성은 \S\ref{sec:falsify}(i) 수송 진단으로 점검한다. 첨자에 관한 주의 하나 — 이
$\eff$ 는 직렬 율속의 합성을 뜻하고, 식~\eqref{eq:keff} 의 $\Delta G_{\eff,j}$ 에 붙은 $\eff$ 는 전위에 의한
배리어 낮춤을 뜻한다. 같은 첨자지만 다른 물리다.
\end{boundbox}

\subsection{전위$\cdot$C-rate 가 peak 를 바꾸는 통로}
식~\eqref{eq:keff} 를 식~\eqref{eq:Lq} 에 넣으면 $L_{q,j}\propto1/k_j^{\eff}$ 라 $\partial\ln L_{q,j}/\partial V_\drive=-\partial\ln k_j^{\eff}/\partial V_\drive$ 이고($1/k_j^{\eff}=1/k_{j,\mathrm{act}}+1/k_{j,\lim}$,
$k_{j,\lim}$ 은 $V_\drive$ 무관 가정), $\partial\mathcal A_j/\partial V_\drive=sF$(식~\eqref{eq:affinity})이고 전위 의존은
$k_{j,\mathrm{act}}$ 만 통하므로(연쇄율 $\partial\ln k_j^{\eff}/\partial V_\drive=\tfrac{k_{j,\lim}}{k_{j,\mathrm{act}}+k_{j,\lim}}\,\partial\ln k_{j,\mathrm{act}}/\partial V_\drive$)
꼬리 길이의 \emph{전위 의존}은
\begin{equation}
\frac{\partial\ln L_{q,j}}{\partial V_\drive}=-\frac{k_{j,\lim}}{k_{j,\mathrm{act}}+k_{j,\lim}}\,\frac{\chi sF}{RT}\le0\quad(s=+1).
\label{eq:LqV}
\end{equation}
활성화 지배 $k_{j,\mathrm{act}}\ll k_{j,\lim}$ 극한에서 인자 $k_{j,\lim}/(k_{j,\mathrm{act}}+k_{j,\lim})\to1$ 이라 $-\chi sF/RT$ 로
환원된다; 유한 공율속이면 그만큼 \emph{감쇠}하나 인자 $\in(0,1]$ 이라 \emph{부호}($\chi>0$ 이면 $<0$ — 전위 의존)는
그대로다(\S\ref{sec:falsify} 반증 지문이 율속 비율에 견고한 까닭).
\begin{keybox}
\textbf{세 인자의 peak 진입.} 구동 전위가 커지면 forward 유효 배리어가 낮아져 꼬리가 짧아진다. 측정축에서는
분극 $s_I|I|R_n$(식~\eqref{eq:vapp})이 apparent peak 위치를 밀므로, C-rate 는 위치 이동에도 기여한다. 이로써
세 인자 — 온도$\cdot$C-rate$\cdot$전위 — 가 모두 $k_j$ 와 $L_{q,j}$ 를 통해 peak 에 들어왔다. 남은 물리는
하나다. 전극 입자가 동일하지 않다는 것이다. 그 비균질이 \S\ref{sec:lag} 에서 예고한 늘어진(stretched) 꼬리를
만든다. 이를 다루려면 분포$\cdot$가중 평균$\cdot$분산 전파의 통계 도구가 필요하므로, 다음 절(\S\ref{sec:stattools})
에서 잠시 그 도구를 정비하고, 그 위에서 \S\ref{sec:dist} 가 입자 분포를 정량화한다.
\end{keybox}

% ====================================================================
\section{입자 분포를 다루는 통계적 기술}\label{sec:stattools}
% ====================================================================
\textbf{도구 정비 막간.} 앞 절들이 한 입자의 물리를 닫았다. 여기서 잠시 멈춰, 입자 \emph{집단}을 다룰 통계
도구를 준비한다. \S\ref{sec:lag} 가 예고하고 \S\ref{sec:barrier} 가 가리킨 늘어진 꼬리가 바로 이 도구로
정량화된다. \S\ref{sec:barrier} 까지는 한 전이를 활성화 배리어 한 값 $\Delta G_{a,j}$ 로 다뤘다. 그러나 실제
전극에서는 입자 집단마다 배리어가 다르므로(\S\ref{sec:dist}), 관측은 배리어 \emph{분포}에 대한 합이 된다. 이를
기술하는 데 필요한 본문 도구는 넷이다 — 분포, 가중 평균, 분산 전파, 연속 중첩. 다섯째 도구인 변수 변환(Jacobian)
은 본문 spine 에는 들어가지 않으며, \S\ref{sec:dist} rigorbox 의 길이-분포 형식 점검 전용이다(아래에 별도로
표시). 본 절은 이 도구들을 확률과 미적분만으로, 본 문제에 맞게 자족적으로 확립한다. 이후 \S\ref{sec:dist} 가
그대로 인용하므로, 본 절의 결과식은 \S\ref{sec:dist} 의 입력이다.

\subsection{활성화 배리어의 분포와 밀도}
전극 입자 집단은 크기$\cdot$결정성$\cdot$국소 staging 이 제각각이라 활성화 배리어 $G$ 가 단일 값이 아니라 분포를
이룬다. 이를 밀도(density) $\rho_G(G)$ 로 기술하며, $\rho_G(G)\,\dd G$ 는 배리어가 구간 $[G,\,G+\dd G]$ 에 드는
집단의 분율이다:
\begin{equation}
\rho_G(G)\,\dd G=\big(\text{배리어가 }[G,\,G+\dd G]\text{ 에 드는 집단의 분율}\big),\qquad [\rho_G]=1/[\text{에너지}].
\label{eq:density}
\end{equation}
전 구간의 분율 합이 1 이므로 정규화(normalization) 조건은
\begin{equation}
\int_0^\infty\rho_G(G)\,\dd G=1.
\label{eq:norm}
\end{equation}
본 장의 $\rho_G$ 는 \emph{용량-분율} 밀도다 — 배리어 $G$ 집단이 전이 용량 $Q_j$ 의 \emph{어느 분율}을 내는가의 \emph{단위
정규화}($\int\rho_G\,\dd G=1$) 분포이며, 총 용량 $Q_j$ 는 식~\eqref{eq:superpose} 에서 \emph{별도}로 곱한다(이중 계상
아님: $Q_j\!\cdot\!\rho_G$ 가 $G$별 용량). 따라서 $\rho_G$ 가중 합이 곧 관측 신호다. \emph{식별 주의}: 이 용량-분율 가중은
입도$\cdot$구조 분포($G\!\sim\!\ln a$)에 \emph{용량-반경 매핑}을 곱해 \emph{모형}하는 것이지 꼬리에서 역산하지 않는다
(forward 예측; 순수 ex-situ 입도만으로는 용량 가중이 안 나오므로 매핑이 그 다리다, \S\ref{sec:falsify}).

\subsection{분포에 대한 가중 평균(ensemble average)}
관측량은 한 집단이 아니라 모든 집단의 기여를 분율로 가중해 합한 값이다. 집단 양 $f(G)$ 의 가중 평균은 합을
적분으로 옮겨
\begin{equation}
\langle f\rangle=\int_0^\infty f(G)\,\rho_G(G)\,\dd G,
\label{eq:wavg}
\end{equation}
즉 각 집단 값 $f(G)$ 를 분율 $\rho_G\,\dd G$ 로 가중한 적분이다. \S\ref{sec:dist} 의 중첩식(식~\eqref{eq:superpose})
은 각 배리어 집단의 꼬리 기여를 $f(G)$ 로 두고 이 평균을 취한 것으로, 임의의 결합이 아니라 분율 가중 합이다.

\subsection{분산의 전파와 $1/RT$ 선형 증폭}
분포는 평균뿐 아니라 그 폭으로도 특징지어진다. 분산(variance) $\sigma_G^2=\langle(G-\langle G\rangle)^2\rangle$,
표준편차(standard deviation) $\sigma_G$ 가 그 척도이며, 선형 척도 변환의 분산은 $\mathrm{Var}(aG)=a^2\sigma_G^2$ 다.
꼬리 길이는 배리어의 지수이므로 — 동일 $T\cdot|I|\cdot$구동력에서 집단 간 차이가 배리어 $G$ 뿐이라
$L_q\propto e^{(G-\chi\mathcal A_j)/RT}$(\S\ref{sec:barrier}$\cdot$\S\ref{sec:dist}; 구동항 $\chi\mathcal A_j$ 는 집단
\emph{공통}이라 아래 const 에 흡수) — 로그를 취하면 $\ln L_q=\text{const}+G/RT$ 로 선형이 되어($\chi\mathcal A_j/RT$
는 그 const 에 들어가 \emph{분산에 무관}), $a=1/RT$ 로
\begin{equation}
\mathrm{Var}(\ln L_q)=\Big(\tfrac{1}{RT}\Big)^2\sigma_G^2,\qquad \sigma_{\ln L_q}=\frac{\sigma_G}{RT}.
\label{eq:varprop}
\end{equation}
배리어 \emph{표준편차}가 로그-길이 표준편차로 $1/RT$ 배 전파되므로($\sigma_{\ln L}=\sigma_G/RT$),
저온($RT\downarrow$)일수록 길이 산포가 커진다. 수치로 보면 $\sigma_G=5$ kJ/mol 일 때 $25^\circ$C 에서
$\sigma_{\ln L}\approx2.0$, $-15^\circ$C 에서 $\approx2.3$ 이다. 로그 폭 $2$ 란 길이가 $e^{\pm2}$, 곧
$7$배 범위로 퍼진다는 뜻이라, 같은 전극 안에서도 집단 간 꼬리 길이가 한 자릿수 가깝게 다를 수 있다. 고정된
입자 분포가 온도를 통해 거동을 바꾸는 메커니즘이며, \S\ref{sec:dist} 의 ``저온에서 꼬리가 더 늘어진다''의 근원이다.

\subsection{변수 변환과 밀도 보존}
분포의 가로축을 배리어 $G$ 에서 길이 $L$ 로 옮길 때 한 집단의 분율은 좌표 선택에 무관하게 보존되므로, 두 밀도는
\begin{equation}
A_L(L)\,\dd L=\rho_G(G)\,\dd G\quad\Longrightarrow\quad A_L(L)=\rho_G\big(G(L)\big)\,\Big|\frac{\dd G}{\dd L}\Big|
\label{eq:jacobian}
\end{equation}
로 연결된다($|\dd G/\dd L|$=Jacobian). 본 문제에 대한 구체적 전개는 \S\ref{sec:dist} 의 심화에서 수행하며,
본문에는 분율 보존만으로 충분하다.

\subsection{지수 감쇠의 연속 중첩}
빠른 집단(작은 $L$)은 peak 근처에서 감쇠를 마치고, 느린 집단(큰 $L$)은 꼬리를 멀리까지 연장한다. 단일 지수
$e^{-x/L}$ 를 여러 $L$ 에 대해 분율로 가중 중첩하면(식~\eqref{eq:wavg}), 결과는 더 이상 단일 지수가 아니라 더
무거운(heavier-tailed) 꼬리 — 초기에 빠르고 후미에서 완만한 형 — 가 된다. 그 점근형이 \emph{정확히} stretched
exponential $e^{-(x/L_{\mathrm{KWW}})^\beta}$($0<\beta<1$)이 되는 것은 $\rho_G$ 가 특정 두꺼운-꼬리 분포, 곧
KWW 의 역-Laplace 핵일 때이며, 지수 $\beta$ 는 그 분포가 정한다 \cite{lindsey1980,johnston2006,svare2000}.
여기서 $L_{\mathrm{KWW}}$ 는 KWW 척도로, \S\ref{sec:dist} rigorbox 의 $L_0$ 와는 별개의 양이다. 요컨대
``stretched''는 $\rho_G$ 형태에 대한 조건부 결론이다. 임의의 $\rho_G$ 가 stretched 를 주지는 않는다 — 좁은
$\rho_G$ 면 거의 단일 지수가 되고, 두꺼운 꼬리면 KWW 가 된다.

\begin{keybox}
\textbf{요약.} 입자 분포는 밀도 $\rho_G$ 로 기술하며 그 정규화는 식~\eqref{eq:norm} 이다. 관측은 그에 대한
가중 평균 $\int f\rho_G\,\dd G$ (식~\eqref{eq:wavg})로 기술한다. 길이가 배리어의 지수이므로, 배리어 표준편차는 로그-길이
표준편차 $\sigma_G/RT$ 로 전파되고(식~\eqref{eq:varprop}) 저온에서 꼬리가 더 늘어진다. 배리어 분포는 분율 보존에
의해 길이 분포로 변환되며(식~\eqref{eq:jacobian}), 지수 감쇠의 연속 중첩은 더 무거운 꼬리를 준다 — 특정 두꺼운-꼬리
$\rho_G$ 에서는 stretched exponential 이 된다. 다음 절(\S\ref{sec:dist})은 이 네 본문 도구(변수 변환은 rigorbox
전용)로 입자 분포를 정량화한다. 거기서 일어나는 일은 새 물리가 아니라 분율 가중 합산이다.
\end{keybox}

% ====================================================================
\section{배리어 분포와 중첩}\label{sec:dist}
% ====================================================================
지금까지 한 전이의 배리어 $\Delta G_{a,j}$ 를 한 값으로 다뤘다. 그러나 실제 전극의 입자는 크기$\cdot$국소
staging$\cdot$결정성$\cdot$결정립계가 제각각이다. 이 차이는 무시할 수 없다. 관측되는 peak 와 꼬리는 여러 입자
집단의 합이기 때문이다. 그렇다면 어떻게 합치는가 — 앞 절에서 준비한 통계 도구가 여기서 쓰인다.

\textbf{입자 분포 $\to$ 배리어 분포 $\to$ 중첩(네 단계).}
\begin{enumerate}[label=\textbf{(\alph*)},leftmargin=2.2em,itemsep=2pt]
\item \textbf{입자별 배리어 차이.} 큰 입자/작은 입자, 결정성이 좋은/나쁜 입자, staging 환경이 다른 도메인은
  같은 전이라도 넘어야 할 활성화 장벽이 다르다.
\item \textbf{배리어 분포.} 한 값 대신 분포 $\rho_G(G)$ 로 기술한다. 이는 \S\ref{sec:stattools} 의 밀도이며
  정규화는 식~\eqref{eq:norm}, $\int_0^\infty\rho_G(G)\,\dd G=1$ 이다. $\rho_G$ 는 꼬리에서 역산하지 않는다 —
  입자 크기$\cdot$국소 상태 분포에 용량-반경 매핑을 곱해 독립적으로 모형하는 용량-분율 분포다
  (\S\ref{sec:stattools} 식별 주의$\cdot$아래 boundbox). 분포를 배리어 $G$ 에만 두는 까닭은 구동력
  $\mathcal A_j$ 가 집단 공통 — 모든 입자가 같은 평균장 $V_n$ 을 본다 — 이기 때문이다. 이 공통성은 이상형
  $\mathcal A_j=sF(V_\drive-U_j)$ 기준의 말이며, $\Omega\ne0$ 일반화에서 생기는 $\xi_j$ 의존 몫은 꼬리
  포화창에서 $\Delta\xi$ 가 작아 1차에서 무시한다.
\item \textbf{집단별 꼬리 길이.} 장벽 $G$ 인 집단은 식~\eqref{eq:Lq},~\eqref{eq:keff} 로 자신의 꼬리 길이
  \begin{equation*}
  L_{q}(G)\simeq\frac{|I|}{Q_\cell}\frac{h}{k_BT}\,e^{(G-\chi\mathcal A_j)/RT}
  \end{equation*}
  를 갖는다(\emph{활성화 지배}
  $k_j^{\eff}\simeq k_{j,\mathrm{act}}$ 형 — \emph{이하 \S\ref{sec:dist}$\cdot$\S\ref{sec:synth} 닫힌식 공통 scope};
  유한 공율속은 식~\eqref{eq:LqV} 의 감쇠인자와 \S\ref{sec:falsify} 경쟁원 배제로 다룬다). \emph{장벽이 길이로
  지수 변환}됨에 유의 — 길이는 배리어의 \emph{지수}이나, \emph{로그 길이}의 표준편차는 $\sigma_G/RT$ 로 ($1/RT$ 배)
  \emph{선형} 전파다. ``지수적 증폭''이라는 말은 변환 $L\!\sim\!e^{G/RT}$ 자체를 가리키는 것이지, 로그축 분산
전파가 지수라는 뜻이 아니다.
\item \textbf{선형 중첩.} 관측되는 $\dd Q$ 는 각 집단의 기여 $\dd Q_G$ 의 \emph{선형 합}이다
  (독립 완화 가정에서 신호가 가산적). 각 집단 기여(식~\eqref{eq:tail}; 그 집단의 잔류 지연 $r_a(G)=r_j(q_a;G)$ 가
  진폭)를 $\rho_G$ 로 가중 적분하면(\S\ref{sec:stattools} 의 가중 평균, 식~\eqref{eq:wavg})
\end{enumerate}
\begin{equation}
\boxed{\;\begin{aligned}
\frac{\dd Q}{\dd V_n}\Big|_{\text{tail}}&=Q_j\!\int_0^\infty\!\rho_G(G)\,\frac{r_a(G)}{L_V(G)}\,
\exp\!\Big[-\frac{V_n-V_a}{L_V(G)}\Big]\,\dd G\quad\big(s(V_n-V_a)\ge0\big),\\[2pt]
&\qquad L_V(G)=\Big|\frac{\dd V_n}{\dd q}\Big|_{q_a}\!\!L_q(G),\qquad r_a(G)\equiv r_j(q_a;G).
\end{aligned}\;}
\label{eq:superpose}
\end{equation}
\textbf{중첩 결과의 해석.} 빠른 집단(작은 $G$, 작은 $L_q$)은 peak 근처에서 감쇠를 마치고, 느린 집단(큰 $G$)이
꼬리를 멀리까지 연장한다. 분포가 좁아지는 극한 $\rho_G\to\delta(G-\bar G)$ 에서는 적분이 sifting 으로 닫혀
식~\eqref{eq:simplefit} 의 단일 지수($L_V=L_V(\bar G)$)로 정확히 환원된다 — \S\ref{sec:synth} 와 정합한다.
\emph{용량 보존}도 성립한다. 일반 분포에서 꼬리 면적은 $Q_j\!\int\!\rho_G(G)\,r_a(G)\,\dd G=Q_j\overline{r_a}$
이고, 포화 근사에서 $\overline{r_a}=1-\overline{\xi_j(q_a)}$ 는 미완 용량과 같으므로 전이 총 용량 $Q_j$ 를 넘지
않는다. 단일 지수의 연속 중첩은 단일 지수보다 무거운 꼬리를 주며, $\rho_G$ 가 충분히 두꺼운 꼬리를 가질 때 그
점근형이 stretched(늘어진) 꼬리, 곧 KWW 형이다 \cite{lindsey1980,johnston2006,svare2000}.
\textbf{\S\ref{sec:lag} 가 예고한 ``실측 꼬리의 늘어짐''이 바로 이 분포 중첩의 귀결}이며, 단일 입자 단일지수
(\S\ref{sec:lag})는 분포가 좁은($\rho_G\to\delta$) 극한일 뿐이다. 어떤 $\rho_G$ 가 두꺼운 꼬리를 주는지는
물리에서 온다 — 입자 반경 $a$ 의 분포가 확산$\cdot$staging 시간을 거쳐 배리어로 가는 사상 $G\sim\ln a$ 가
$\rho_G$ 의 꼬리를 정한다.

\textbf{저온 꼬리 확장(분산 전파 — \S\ref{sec:stattools}).} 식~\eqref{eq:superpose} 의 $L_V(G)$ 는
$\ln(L_V/L_{V,0})=G/RT$ 를 만족한다. 기준 길이 $L_{V,0}$ 는 구동항 $\chi\mathcal A_j$ 등 집단 공통 상수를 흡수해
로그를 무차원화하는 역할이다 — $L_V\propto L_q\propto e^{(G-\chi\mathcal A_j)/RT}$ 에서 $\chi\mathcal A_j$ 는
상수 오프셋이라 비에서 탈락하기 때문이다. 따라서 배리어 분산 $\sigma_G^2$ 는
$\mathrm{Var}[\ln(L_V/L_{V,0})]=\sigma_G^2/(RT)^2$ 로, 곧 표준편차 $\sigma_G/RT$ 로 전파된다
(식~\eqref{eq:varprop} 을 $L_V\propto L_q$ 에 적용한 것). \textbf{저온일수록($RT\downarrow$) 길이 분산이 커져 꼬리가 더 늘어진다} — 고정된
배리어 분포가 온도를 통해 늘어짐을 만드는 메커니즘이며, 이것이 관측의 ``저온 긴 꼬리''를 \emph{입자 분포까지
포함해} 설명한다. \emph{단 이 $\sigma_G/RT$ 척도는 $\sigma_G$ 유한한 $\rho_G$(좁은$\cdot$중간 폭)에서 정량적}이다 —
점근형이 정확한 KWW stretched 가 되는 두꺼운-꼬리 $\rho_G$ 에선 $\sigma_G$ 가 발산할 수 있어, 그 영역은 stretched
지수 $\beta$ 를 온도별 \emph{현상학 지수}로 두고 그 온도 추이로 본다. 둘은 \emph{배타 regime} 이 아니라 \emph{같은 중첩 적분의 두 점근}이다 —
코어가 유한분산이면서 꼬리만 무거운 흔한 $\rho_G$ 에서는 $\sigma_G/RT$(코어 폭)와 $\beta(T)$(무거운 꼬리)가 \emph{동시}
적용된다.

\begin{rigorbox}
\textbf{길이 분포로의 좌표 변환(밀도 보존, 세 단계).} 식~\eqref{eq:superpose} 의 적분을 배리어 대신 길이 $L_q$
(전압축은 $L_V=|\dd V_n/\dd q|\,L_q$ 로 비례)로 옮기는 변환(\S\ref{sec:stattools} 의 변수 변환, 식~\eqref{eq:jacobian} 의 본 문제 형식 전개)을 명시한다. \emph{단계 1(밀도 보존)}: 집단의 분율은 좌표를 바꿔도 보존되므로 $A_L(L_q)\,\dd L_q=\rho_G(G)\,\dd G$.
\emph{단계 2(Jacobian)}: 역함수 $G(L_q)=\chi\mathcal A_j+RT\ln(L_q/L_0)$($L_0=|I|h/(Q_\cell k_BT)$)에서
$\dd G/\dd L_q=RT/L_q>0$ 이라 일대일이므로 $A_L(L_q)=\rho_G(G(L_q))\,RT/L_q$. \emph{단계 3(지지 범위)}: 배리어는
음수가 못 되므로($G\ge0$) 길이에 최단 하한 $L_{\min}=L_0\,e^{-\chi\mathcal A_j/RT}$ 가 생긴다($L_q<L_{\min}$ 에서
$A_L=0$). 본문에는 이 형식이 필요 없다 — 식~\eqref{eq:superpose} 의 ``$\rho_G$ 가중 합''만으로 충분하다.
\end{rigorbox}
\begin{boundbox}
\textbf{독립 완화 가정·비유일성.} 식~\eqref{eq:superpose} 는 입자 집단이 \emph{서로 독립}으로 완화한다고 본다
(평균장) \cite{funabiki1999}. \emph{협동성과의 화해(스케일 분리)}: \S\ref{sec:regsol} 의 협동(상분리$\cdot$핵생성)은 한
입자 \emph{내부}의 일차 전이로 이미 그 입자의 평형 $\xi_{\eq,j}$ 형태에 흡수돼 있고, 본 절의 독립은 \emph{입자 \emph{간}}
(서로 다른 입자)에만 적용된다 — 내부 협동 $\ne$ 집단 간 독립이라 모순이 아니다. 또 $\rho_G$ 가 배리어 $G$ \emph{하나}로 모든 집단
차이를 흡수한다는 것($r_a(G)\cdot L_V(G)$ 와의 상관이 평균장$\cdot$단일전위 $V_n$ 근사로 제한)이 본 절의 최강 가정이다.
같은 맥락에서 $q_a$ 는 평형 형태가 정하는 \emph{집단 공통} 좌표라 — 빠른 집단(작은 $L_q$)은 $q_a$ 이전에 이미 거의
완화하고 느린 집단은 상승부 잔여가 남을 수 있어 — $r_a(G)$ 를 공통 $q_a$ 에서 평가하는 것은 분포 폭이 상승부 폭보다
좁을 때의 1차 근사다. 무시한 그 상관은 꼬리를 더 무겁게 미는 방향이므로 stretched 결론을 약화시키지 않는다.
또 stretched 꼬리에서 $\rho_G$ 로의 역산은
\emph{비유일}하므로($\rho_G$ 외 다른 메커니즘도 같은 꼬리를 낼 수 있다), $\rho_G$ 를 꼬리에서 역산하지 않고
독립 측정/모형한 $\rho_G$ 로 \emph{forward 예측}만 한다(\S\ref{sec:falsify}).
\end{boundbox}

% ====================================================================
\section{종합: 피팅 가능한 식과 3$\times$3 의존성}\label{sec:synth}
% ====================================================================
여기까지 여섯 개의 물리 절 — 전하 보존, 평형, 정규용액, 지연, 유효 배리어, 분포(통계 도구절 \S\ref{sec:stattools}
는 도구이므로 제외) — 이 각각 한 조각씩을 닫았다. 이 절은 그것들을 묶는다. 묶으면 $\dd Q/\dd V$ 가
$(V;\,T,|I|,V_\drive)$ 의 닫힌 함수가 된다. 이 중 정규용액의 몫은 둘로 갈린다. 본 절의 방전 단일 전이 닫힌식에는
$w_j$ 의 물리적 근거(식~\eqref{eq:weff})로 들어가고, 충방전 분기 몫은 \S\ref{sec:hys_branch} 에서 합류한다.
본 절이 답할 질문은 둘이다. peak 의 위치$\cdot$너비$\cdot$높이는 세 인자에 각각 어떻게 반응하는가. 그리고
실무에 바로 쓸 가장 단순한 식은 무엇인가.

\subsection{닫힌 형태 — ``평형에서 지연을 뺀 것의 미분''}
\textbf{\S\ref{sec:dist} 의 꼬리를 상승부와 한 식으로 묶기.} \S\ref{sec:dist} 는 \emph{한 전이의 꼬리}(post-peak)만
배리어 분포 적분 식~\eqref{eq:superpose} 으로 줬다. 그러나 관측 peak 는 상승부(rise)와 꼬리가 \emph{한 곡선}이다.
이제 둘을 \emph{하나의 닫힌 식}으로 묶으면, \S\ref{sec:eqpeak} 의 평형 peak 가 rise 를, 식~\eqref{eq:superpose} 가
그 뒤 꼬리를 담당함이 한눈에 드러난다.

관측 peak 항은 식~\eqref{eq:dQdV} 의 $Q_j\,\dd\xi_j/\dd V_n$ 이다. 실제 전환율은 어느 영역에서나
$\xi_j=\xi_{\eq,j}-r_j$ 로 적을 수 있으며, 지연 $r_j$ 의 거동은 \S\ref{sec:lag} 의 일반해가 전 구간에서 정해
준다. 상승부에서는 극한 (i)에 따라 $r_j\simeq L_q\,\dd\xi_{\eq,j}/\dd q$ 의 소량이고, post-peak 에서는 극한
(ii)에 따라 지수 감쇠다. 입자 집단을 들이면 이 지연을 집단 평균해야 하므로, \emph{가중 평균 지연}
\[
\bar r_j(q)\equiv\int_0^\infty\rho_G(G)\,r_j(q;G)\,\dd G,\qquad
r_j(q;G)=r_a(G)\,e^{-(q-q_a)/L_q(G)}\qquad(q\ge q_a)
\]
을 빼서 $\xi_j=\xi_{\eq,j}-\bar r_j$ 로 쓴다. 여기서 $r_j(q;G)$ 는 집단 $G$ 의 지연(식~\eqref{eq:rsol})로
$q\ge q_a$ 에서 정의되며, $q<q_a$ 의 지연은 극한 (i)의 quasi-static 값이다. 즉 $r_j(q;G)$ 는 \S\ref{sec:lag} 가 \emph{한 집단}에 대해 구한 지연이고,
$\bar r_j$ 는 그것을 \S\ref{sec:dist} 의 분포 $\rho_G$ 로 평균(\S\ref{sec:stattools} 의 가중 평균)한 양이다.

\textbf{이 지연의 미분이 곧 식~\eqref{eq:superpose} 의 꼬리다(연결 확인).} $\xi_j=\xi_{\eq,j}-\bar r_j$ 를
식~\eqref{eq:dQdV} 에 넣으면 peak 항이 \emph{평형 미분} 빼기 \emph{지연 미분} 두 조각으로 갈린다. 꼬리를 만드는
지연 조각 $-Q_j\,\dd\bar r_j/\dd V_n$ 을 풀면, post-peak 에서 한 집단의 지연이 $\dd r_j/\dd q=
-(r_a(G)/L_q(G))\,e^{-(q-q_a)/L_q(G)}$ 로 감쇠하고, 연쇄율 $\dd/\dd V_n=(\dd q/\dd V_n)\,\dd/\dd q$ 와
$L_V(G)=|\dd V_n/\dd q|\,L_q(G)$ 를 넣으면
\[
-Q_j\frac{\dd\bar r_j}{\dd V_n}
=Q_j\!\int_0^\infty\!\rho_G(G)\,\frac{r_a(G)}{L_V(G)}\,e^{-(V_n-V_a)/L_V(G)}\,\dd G,
\]
즉 식~\eqref{eq:superpose} 의 꼬리가 \emph{정확히 재현}된다 — \S\ref{sec:dist} 의 분포 적분은 새 가정이 아니라
$\bar r_j$ 미분의 다른 표기였던 것이다. 그러므로 peak 전체는
\begin{equation}
\boxed{\;\frac{\dd Q}{\dd V_n}=C_\bg+Q_j\frac{\dd}{\dd V_n}\big(\xi_{\eq,j}-\bar r_j\big)
=\underbrace{C_\bg+Q_j\frac{\dd\xi_{\eq,j}}{\dd V_n}}_{\text{평형 peak (rise)}}
\;-\;\underbrace{Q_j\frac{\dd\bar r_j}{\dd V_n}}_{\text{지연이 깎는 꼬리}}\;.}
\label{eq:closed}
\end{equation}
이것이 종합식이다. 모호한 결합 연산이 아니라 \emph{평형 전환율에서 지연을 뺀 것의 미분}이라는, 정의된 형태다.
두 영역의 분담은 \S\ref{sec:lag} 일반해의 두 극한 그대로다. peak 앞쪽(상승부)에서는 극한 (i)대로 지연이
$O(L_q)$ 의 소량이라 평형 항 $Q_j\dd\xi_{\eq,j}/\dd V_n$ 이 지배하고, peak 뒤쪽에서는
$\dd\xi_{\eq,j}/\dd V_n\to0$ 이라 극한 (ii)의 지수 꼬리 $-Q_j\dd\bar r_j/\dd V_n$(식~\eqref{eq:superpose})이
지배한다. 빠른 kinetics 극한($\bar r_j\to0$)에서는 평형 peak(\S\ref{sec:eqpeak})로 복귀하므로 식이 무결하다.
측정축으로는 $V_n\to V_\app=V_n+s_I|I|R_n$ 으로 환산한다.

\subsection{3$\times$3 의존성 표 (실무 산출물)}
표의 \emph{행}은 peak 의 위치$\cdot$너비$\cdot$높이, \emph{열}은 세 인자($T\downarrow$$\cdot$$|I|\uparrow$$\cdot$$V_\drive\uparrow$)다.
각 칸은 그 인자가 그 특성을 바꾸는 \emph{경로} — \emph{어느 식의 어느 항}이 지배하고 부호가 어느 쪽인지 — 를 위
도출에서 직접 읽은 것이며($s=+1$ 방전 기준), 평형 항(\S\ref{sec:eqpeak})과 동역학 꼬리 항(\S\ref{sec:lag}$\cdot$\S\ref{sec:dist})이
\emph{반대로} 미는 칸(특히 너비$\cdot$높이의 $T$ 열)이 곧 관찰 해소의 핵심이다.
\begin{center}
\renewcommand{\arraystretch}{1.4}
\begin{tabular}{p{0.12\textwidth} p{0.26\textwidth} p{0.24\textwidth} p{0.25\textwidth}}
\toprule
 & \textbf{온도 $T\downarrow$} & \textbf{C-rate $|I|\uparrow$} & \textbf{구동 전위 $V_\drive\uparrow$} \\
\midrule
\textbf{위치} & $U_j(T)$ 이동, $\partial U_j/\partial T=\Delta S_j/(sF)$ \eqref{eq:dUdT} & 분극 $+s_I|I|R_n$ 이동 \eqref{eq:vapp} & 위치 직접 이동 없음(\emph{단상 이상 한정}) — $V_\peak=U_j$ 는 $V_\drive$ 무관 \eqref{eq:eqpeak}; \eqref{eq:LqV} 로 너비$\cdot$높이에만($\Omega_j>2RT$ 분기서는 중심이 방향 의존 이동, \S\ref{sec:hys_branch}) \\
\textbf{너비} & 평형 $w_j=RT/F\downarrow$ \emph{좁힘} \eqref{eq:logistic}; 지연$\cdot$분포 꼬리 $\uparrow$ \emph{넓힘}($L_{V,j}\!\gtrsim\!w_j$ 면 우세) \eqref{eq:tailTC}\eqref{eq:superpose} & 꼬리 $L_{q,j}\propto|I|\uparrow$ 넓힘 \eqref{eq:tailTC} & 유효 배리어$\downarrow\Rightarrow L_{q,j}\downarrow$ 좁힘 \eqref{eq:LqV} \\
\textbf{높이} & 평형 $1/w_j\uparrow$, 지연$\cdot$분포가 낮춤(면적 보존) \eqref{eq:tailTC} & $|I|\uparrow\Rightarrow$ 낮춤 \eqref{eq:tailTC} & 배리어$\downarrow\Rightarrow$ 꼬리 짧아져 높임 \eqref{eq:LqV} \\
\bottomrule
\end{tabular}
\end{center}
\textbf{표 읽는 법(칸별 기전).} 표는 외울 것이 아니라 \emph{어느 식의 어느 항}에서 읽히는지를 보는 것이다.
\emph{위치 행} — 온도는 평형 중심 $U_j(T)$ 를 반응 엔트로피만큼 옮기고(식~\eqref{eq:dUdT}), C-rate 는 분극
$s_I|I|R_n$ 으로 \emph{apparent} 위치를 밀며(식~\eqref{eq:vapp}), 구동 전위는 평형 중심을 \emph{옮기지 못한다}
($V_\peak=U_j$ 는 $V_\drive$ 무관 — 전위는 평형 target 이 아니라 도달 \emph{속도}만 바꾸므로). \emph{너비 행} — 평형
폭 $w_j=RT/F$ 는 저온서 \emph{좁아}지지만(식~\eqref{eq:logistic}), 동역학 꼬리 $L_{V,j}$ 는 저온$\cdot$고율서
\emph{넓어}지고(식~\eqref{eq:tailTC}), 구동 전위는 유효 배리어를 낮춰 꼬리를 \emph{좁힌다}(식~\eqref{eq:LqV}).
\emph{높이 행} — 면적이 $\sim Q_j$ 로 보존되므로 너비와 반대로 움직여, 평형은 저온에서 높이고 꼬리는 낮춘다.
이 면적-반비례는 꼬리가 면적의 일부일 때의 1차 어림이다. 꼬리가 우세해지면 면적이 상승부와 꼬리로 갈려 단순
반비례가 약화된다(식~\eqref{eq:simplefit}). 또 $V_\drive$ 열의 ``좁힘$\cdot$높임''(식~\eqref{eq:LqV})은
$\mathcal A_j=sF(V_\drive-U_j)>0$ 인 꼬리 영역에 한하며, 상승부($\mathcal A_j\lesssim0$)에서는 무효이거나
반대다. \textbf{이 표의 핵심은 너비$\cdot$높이의
$T$ 열에서 평형 항과 꼬리 항이 \emph{반대 부호}로 경쟁}한다는 점이다 — 평형만 보면 ``저온서 좁고 높다''인데, 꼬리가
이기면 관측은 ``저온서 넓고 낮다''로 뒤집힌다.

\textbf{핵심 관찰 해소.} 평형은 너비를 ``저온서 좁힘''으로 예측하나, 동역학 꼬리가 평형 폭보다 크면($L_{V,j}
\gtrsim\Delta V_{1/2}\sim w_j$; 저온$\cdot$유한 전류에서 일반적) 꼬리 넓힘이 우세해 \emph{저온서 넓고 비대칭}이
된다 — 관찰과 일치. (두 항의 \emph{크기} 우열은 데이터로 확인할 정량 문제이며, 본 장은 각 항의 \emph{부호}와
\emph{우세 조건} $L_{V,j}\gtrsim w_j$ 를 도출한다.) C-rate$\cdot$전위 방향도 가정이 아니라
위 식들에서 부호로 도출됐다.

\subsection{실무 피팅 근사식}
\textbf{우세 집단 극한.} $\rho_G$ 가 한 배리어 $\bar G$ 둘레에 좁게 몰리면($\rho_G\to\delta(G-\bar G)$),
식~\eqref{eq:superpose} 의 적분이 델타 sifting 으로 닫힌다($\int\delta(G-\bar G)f(G)\,\dd G=f(\bar G)$):
\begin{equation}
\boxed{\;L_{q,j}\simeq\frac{|I|}{Q_\cell}\frac{h}{k_BT}\,e^{(\bar G-\chi\mathcal A_j)/RT},\qquad
\frac{\dd Q}{\dd V_n}\Big|_{\text{tail}}\approx\frac{Q_j\,r_{a,j}}{L_{V,j}}\,e^{-(V_n-V_a)/L_{V,j}}\;,\;\;
L_{V,j}=\Big|\frac{\dd V_n}{\dd q}\Big|_{q_a}L_{q,j}\;.}
\label{eq:simplefit}
\end{equation}
즉 peak 는 \emph{평형 logistic 상승부 + 단일 지수 꼬리}이고, 꼬리는 특성 길이 $L_{V,j}$(전압 단위)의 지수 감쇠다
(꼬리는 식~\eqref{eq:tail} 처럼 post-peak \emph{단방향} $s(V_n-V_a)\ge0$ 에서만).
\textbf{단일지수 vs stretched 선택.} 분포가 좁으면($\rho_G\to\delta$) 이 단일지수가 정확하고(\S\ref{sec:lag} 환원),
넓어 꼬리가 \emph{휘면}(stretched) 식~\eqref{eq:superpose} 의 분포 적분으로 돌아가 $\rho_G$ 폭을 독립 입력으로
더한다 — \S\ref{sec:stattools}$\cdot$\S\ref{sec:dist} 의 분포는 \emph{버린 우회가 아니라} 단일지수가 깨지는
데이터에서 켜지는 항이다.
여기서 $r_{a,j}\equiv r_j(q_a)=\xi_{\eq,j}(q_a)-\xi_j(q_a)\in(0,\xi_{\eq,j}(q_a)]$(포화창 $\xi_{\eq,j}(q_a)\simeq1$ 서
상한 $1$)는 꼬리 시작점의 잔류 지연(전이가 꼬리로 넘기는
미완 분율)이며, 꼬리 면적은 $\int(\dd Q/\dd V_n)_{\text{tail}}\,\dd V_n=Q_j\,r_{a,j}$ 다. 정확히는 상승부가
$Q_j\,\xi_j(q_a)=Q_j[\xi_{\eq,j}(q_a)-r_{a,j}]$ 를, 꼬리가 $Q_j\,r_{a,j}$ 를 담아 합이 $Q_j\,\xi_{\eq,j}(q_a)$ 인데,
꼬리 영역의 \emph{포화 근사} $\xi_{\eq,j}(q_a)\simeq1$(식~\eqref{eq:Lq} 유도 전제) 하에서 상승부 $\simeq Q_j(1-r_{a,j})$,
총합 $\simeq Q_j$ 로 보존된다. \emph{이 인자 $r_{a,j}$ 없이 $Q_j/L_{V,j}$ 로 두면 꼬리 항만으로 $Q_j$ 가 소진되어 용량 보존이 깨진다.}
\begin{verifybox}
\textbf{Worked example — 한 전이의 peak 치수($Q_j=0.5\,Q_\cell$ 가정).} 평형 폭은 $25^\circ$C 에서
$w_j=RT/F=25.7$ mV 라 FWHM $=3.53\,w_j\approx91$ mV, 순 peak 높이 $=Q_j/(4w_j)=0.5\,Q_\cell/(4\times0.0257)
\approx4.9\,Q_\cell/\mathrm V$. 같은 전이가 $-15^\circ$C 에선 $w_j=22.2$ mV 로 좁아져 FWHM $\approx78$ mV, 높이
$\approx5.6\,Q_\cell/\mathrm V$ — \emph{평형만 보면 저온서 더 좁고 높다}. 그러나 동역학 꼬리 길이 $L_{q,j}\propto
e^{\Delta H_a/RT}$ 는 $\Delta H_a=40$ kJ/mol 이면 $25^\circ$C$\to-15^\circ$C 에서 약 $12$배 길어져(\S\ref{sec:lag}),
꼬리가 평형 폭을 넘어서면($L_{V,j}\gtrsim w_j$) 관측 peak 는 거꾸로 \emph{저온서 넓고$\cdot$낮고$\cdot$비대칭}이
된다 — 표의 너비$\cdot$높이 $T$ 열에서 평형과 꼬리가 반대로 미는 바로 그 칸이다. 이 한 예가 본 장 전체의
관찰(저온 긴 꼬리)을 수치로 압축한다. \emph{이 수치는 뒤에서 쓰인다}: $91$ mV 는 \emph{이상 극한($\Omega_j=0$) FWHM
상한}이다 — 실측 staging peak 는 더 좁다. 단상($\Omega_j<2RT$)이면 $w^\eff$(식~\eqref{eq:weff})만큼 좁아지고,
상분리($\Omega_j>2RT$)면 plateau 라 더욱 날카롭다(\S\ref{sec:hys_peak}). 따라서 이 수치는 S1 의 OCV peak
판독에서 상한 sanity-check 로만 쓴다 — 읽은 FWHM 이 $91$ mV 를 넘으면 이상 폭보다 넓다는 뜻이라 꼬리 오염
신호다. 등호로 닫지 않는다; 다음 절 융합 판정
(\S\ref{sec:overlap})도 이 $91$ mV 를 \emph{이상 상한 규모}로(동일 폭 가정 하) 전이 간격과 비교하는 입력으로 받는다.
\end{verifybox}

\textbf{식별 순서(비순환).} 평형항과 꼬리항을 한 데이터에서 동시에 자유 피팅하면, 둘이 같은 전위창에서 섞여
공선형이 되어 분리할 수 없다. 이 공선형은 평형$\leftrightarrow$꼬리 축의 문제다. 해법으로 두 종류의 식별을
구분한다. (a) \textbf{독립 측정으로 애초에 공선형이 없는 경우.} 저율 OCV(S1)에서는 꼬리가 물리적으로 부재
($|I|\to0$)하므로, 한 곡선이어도 $\{U_j,w_j,Q_j\}$ 가 서로 직교한 peak 모멘트 — 위치는 중심(1차 모멘트),
너비는 중심 곡률, 면적은 0차 모멘트(적분) — 라서 동시에 분리된다. 전류 차단(S2)은 전압 도약이라는 별개의
관측량이다. (b) \textbf{공선형을 해소해야 하는 경우.} $\chi$(S3)와 $\Delta H_a$(S4)는 둘 다 꼬리 지수에
섞이므로, 직교 측정축 — 등온 다전위 대 다온도 — 으로 한 겹씩 벗긴다. 요컨대 독립 실험으로 먼저 고정해 뒤 단계의
기지값으로 만들면, 각 단계에 미지수가 하나씩만 남아 순환이 끊긴다. 아래의 순차 제약이 그 벗김의 순서이며, 각
단계는 ``무엇을 고정하고, 무엇이 분리되며, 공선형이 어떻게 해소되는가''로 읽는다.
\begin{enumerate}[label=\textbf{S\arabic*.},leftmargin=2.6em,itemsep=2pt]
\item \textbf{저율 OCV}($|I|\to0$, 꼬리가 사라진 평형). 고정$\to\{U_j,w_j,Q_j\}$. 분리 근거: 동역학 꼬리가
  없어 평형 peak(식~\eqref{eq:eqpeak})만 남으므로 위치$\cdot$너비$\cdot$면적이 곧 $U_j,w_j,Q_j$.
\item \textbf{펄스/전류 차단}. 고정$\to R_n$. 분리 근거: 전류 차단 직후 전압의 즉각 도약(IR-drop) $\Delta V=
  s_I|I|R_n$ 에서 $R_n$ 을 읽는다(GITT) \cite{weppner1977}. $R_n$ 과 $\chi$ 가 둘 다 꼬리를 늘려 공선형이므로
  $R_n$ 을 \emph{먼저} 묶는다.
\item \textbf{다전류/전위} — $\chi$ 를 Arrhenius 앞에서 먼저 확정한다. 고정$\to\chi$. 분리 근거:
  식~\eqref{eq:LqV} 의 $\partial\ln L_{q,j}/\partial V_\drive$ 기울기에서 $\chi$ 가 나온다. 활성화 지배
  극한에서 그 기울기는 $-\chi sF/RT$ 이고, 공율속이 유한하면 식~\eqref{eq:LqV} 의 감쇠인자
  $k_{j,\lim}/(k_{j,\mathrm{act}}+k_{j,\lim})$ 로 보정한다. 이때 $\mathcal A_j$ 는 S1$\cdot$S2 의 기지값으로
  계산된다. $\chi$ 를 다음 단계 S4 보다 먼저 확정하는 까닭은, S4 의 구동항 보정 $\chi\mathcal A_j$ 를 기지값으로
  만들기 위해서다. 이 순서가 ``$\chi$ 와 $\Delta H_a$ 가 둘 다 꼬리에 들어가 공선형''인 순환 의존을 끊는다.
  실측으로는 고정 $T$ 에서, S1$\cdot$S2 로 환산한 동일 SOC$\cdot$tail window 의 $\ln L_{q,j}$ 를 $V_\drive$ 로
  회귀한다(상세 절차는 \S\ref{sec:master} 의 피팅 알고리즘).
\item \textbf{다온도 Arrhenius}. 고정$\to\Delta H_a$(순수, 기울기에서; $\Delta S_a$ 는 절편 $\Delta S_a/R+\text{const}$
  \emph{조합}으로만 — 절 끝에서 상술). \emph{보정의 대수 유도}: 식~\eqref{eq:simplefit}
  에 $\Delta G_a=\Delta H_a-T\Delta S_a$ 를 넣으면
  \begin{equation*}
  \begin{gathered}
  \ln L_{q,j}\simeq\ln\frac{T_*}{T}+\frac{\Delta H_a-T\Delta S_a-\chi\mathcal A_j}{RT},\\[2pt]
  T_*\equiv\frac{|I|h}{Q_\cell k_B}\ (\text{특성 온도, K — 로그 인자 무차원화}).
  \end{gathered}
  \end{equation*}
  $\ln(T_*/T)$ 의 $-\ln T$ 부분은 \emph{약한 온도 의존 보정항}, 엔트로피 항 $-\Delta S_a/R$(=$-T\Delta S_a/RT$)는
  \emph{상수 절편}($1/T$ 무관)이며, 구동항 $-\chi\mathcal A_j/RT$ 는 $\mathcal A_j(T)$ 때문에 $1/T$ 에 비선형으로
  섞인다. 좌변에서 $\ln(T_*/T)$ 와 구동항을 이항해
  \begin{equation}
  \boxed{\;\begin{aligned}
  &y(T)\equiv\ln\frac{T_*}{L_{q,j}\,T}-\frac{\chi\mathcal A_j(T)}{RT}\ \text{를}\ \tfrac1T\ \text{로 회귀}\\[2pt]
  &\quad\Rightarrow\ \text{기울기}=-\frac{\Delta H_a}{R}\ (\text{순수})
  \end{aligned}\;}
  \label{eq:arrhenius}
  \end{equation}
  (이항 후 남는 $1/T$ 계수는 $-\Delta H_a/R$ 뿐, 절편은 $\Delta S_a/R+\text{const}$). 분리 근거: prefactor$\cdot$
  엔트로피$\cdot$구동항을 보정해 순수 엔탈피만 기울기에 남긴다. 보정에 쓰는 $\chi\mathcal A_j$ 는 앞 단계
  S3 의 기지 $\chi$ 와 S1$\cdot$S2 로 계산되므로 순환이 없다.
\end{enumerate}
이로써 $\{U_j,w_j,Q_j,R_n,\Delta H_a,\chi\}$ 가 \emph{위 도출식만으로}(\emph{활성화 지배} 가정 하 — 유한 공율속이면
$k_{j,\lim}$ 이 추가 미지수라 \S\ref{sec:falsify} (i) 수송 배제로 닫는다) 데이터에서 분리 추출된다. \emph{단 $\Delta S_a$ 는}
Arrhenius 절편 $\Delta S_a/R+\text{const}$ 에서 \emph{현상론적 prefactor}(흡수된 $k_0\cdot$활성면적 등; \S\ref{sec:lag} 의
$k_0$ 주)와 \emph{결합}으로만 나오므로, prefactor 를 독립 고정하지 않으면 $\Delta S_a$ 단독값은 식별되지 않는다(절편
조합으로만 결정). $\rho_G$(분포 폭)는 꼬리가 단일 지수에서 벗어나(stretched) 휠 때 독립 입력으로 더한다(역산 금지).
\begin{boundbox}
\textbf{본 장 범위와 식별 교란 요인.} (i) 본 \emph{절}의 닫힌식은 \emph{단일 방향}(방전)$\cdot$단일 전이 한정이다 —
충방전 OCV hysteresis(분기)는 본 장 \emph{안}의 \S\ref{sec:hys_branch} 가 유도하고 \S\ref{sec:master} 통합형이
확장한다(충전 \emph{꼬리} 동역학만 방향별 별도). (ii) 여러 전이가 겹치는 영역(관측의 ``다음 peak 와 겹침'')은 \S\ref{sec:overlap} 에서 전이 $j$ 별
합산$\cdot$OCV-anchored 분해로 본격 다룬다. (iii) 배경 용량 $C_\bg$(식~\eqref{eq:dQdV})은 꼬리와 공선형이
될 수 있어, S1 에서 저율 OCV 로 함께 고정한 뒤 그 위에서 꼬리를 회귀한다. (iv) 본 장 전위는 \emph{반쪽전지}(상대전극
기준 내부 $V_n$) 척도이며 ICA 도 $\dd Q/\dd V_n$ 이다 — full-cell $V_\cell=V_p-V_n$ 의 $\dd Q/\dd V_\cell$ 는 양극
기여가 섞이므로, 본 장 결과를 full-cell 데이터에 적용하려면 양극 분리(반쪽전지 기준 환원)가 선행돼야 한다.
\end{boundbox}

% ====================================================================
\section{다중 전이 겹침과 forward 합산}\label{sec:overlap}
% ====================================================================
방전 본론은 이제 두 조각을 남겨 두고 있다. 앞 절(\S\ref{sec:synth})이 \emph{단일 전이}의 peak 를 닫힌식과
식별 절차로 닫았으니, 남은 것은 ``전이가 여럿일 때''(본 절)와 ``같은 곡선을 역미분 좌표로 볼 때''
(\S\ref{sec:dva})다. 이 둘까지 닫으면 방전 $\dd Q/\dd V$ 의 이론이 완결된다. 본 절의 주제는 겹침이다. 관측의 한
축은 ``C-rate 가 커지면 peak 가 겹친다''였고, 흑연은 전이가 여럿이다($N_p\approx3\sim4$). 그 단독 곡선들이
겹칠 때 관측 ICA 는 어떻게 합쳐지는가. 미리 방침부터 밝혀 둔다. \textbf{본 절은 융합 곡선을 전이별로 역산하는
``분해''를 하지 않는다.} 그 역산은 비유일이어서 금지된다(아래$\cdot$\S\ref{sec:dist}). 대신 저율 anchor 로 단독
곡선을 고정한 뒤, 고율로 \textbf{forward 합산}해 관측과 대조한다. 식별의 방향은 ``저율로 고정 $\to$ 고율로
예측$\cdot$검증'' 하나뿐이다.

\textbf{총 ICA 의 선형 중첩.} 전하 보존식~\eqref{eq:charge} 은 처음부터 $\sum_j Q_j\xi_j$ 형이었고, 그 미분
~\eqref{eq:dQdV} 이 곧 총 ICA 다:
\begin{equation}
\boxed{\;\frac{\dd Q}{\dd V_n}=C_\bg(V_n,T)+\sum_{j=1}^{N_p}Q_j\,\frac{\dd\xi_j}{\dd V_n}\;}\qquad(\text{식~\eqref{eq:dQdV} 의 다전이 형})
\label{eq:total}
\end{equation}
각 항 $Q_j\,\dd\xi_j/\dd V_n$ 은 식~\eqref{eq:closed} 의 (평형 rise $-$ 지연 꼬리) 구조 그대로다. 여기서
두 층위를 구분해 둔다. \emph{합산 자체}, 곧 선형 가산은 무조건 성립한다 — 식~\eqref{eq:total} 은 전하
보존~\eqref{eq:charge} 의 직접 미분이므로 가정이 아니다. 반면 ``전이 \emph{간} 독립''은 별개의 가정이다. 이는
각 $\xi_j$ 의 동역학($L_{q,j}$ 와 $\xi_{\eq,j}$)이 이웃 전이의 진행에 불변이라는 가정으로, 전이 간 $\Omega$-결합과
순차 staging 의 상호작용을 무시한다. \S\ref{sec:dist} 의 ``전이 \emph{내} 입자 독립''과는 다른 층위다. 요컨대
합은 보존에서 오고, 각 항의 독립 진화는 이 별도 가정에서 온다.

\textbf{겹침의 정량 기준.} 인접 전이 $j,j{+}1$ 의 단독 곡선이 겹치는 조건은 둘이다: (i) \emph{평형} peak
폭으로도 가까우면 — 두 peak 의 \emph{반치반폭}($\Delta V_{1/2}/2$ 씩)이 서로 닿는 조건, 곧 그 합
$|U_{j+1}-U_j|\lesssim\tfrac12\Delta V_{1/2,j}+\tfrac12\Delta V_{1/2,j+1}=(\Delta V_{1/2,j}+\Delta V_{1/2,j+1})/2$
(각 $\Delta V_{1/2}\simeq3.53\,w$; 식~\eqref{eq:eqpeak}); (ii) \emph{동역학} 꼬리가 옆 peak 까지 닿으면 — 꼬리 길이가 중심 차 이상
$L_{V,j}\gtrsim|U_{j+1}-U_j|$(식~\eqref{eq:tail}; 꼬리는 단방향이라 $j$ 의 꼬리는 방전 방향 \emph{다음} 전이 쪽으로
뻗는다). \textbf{(ii)가 C-rate$\cdot$저온의 겹침이다}: 식~
\eqref{eq:tailTC} 로 $|I|\uparrow$ 또는 $T\downarrow$ 면 $L_{V,j}\uparrow$ 라, 저율서 분리됐던 peak 가 고율서 조건
(ii)를 넘겨 융합한다. 융합 시작 전류$\cdot$온도는 위 부등식의 등호로 추정된다.
\begin{verifybox}
\textbf{Worked example — 두 peak 의 융합 개시.} 인접 stage 전이가 $\Delta U=|U_{j+1}-U_j|\approx100$ mV 떨어져
있고 평형 폭이 \S\ref{sec:synth} worked example 의 값($w_j=25.7$ mV, FWHM $\approx91$ mV)이면, 저율 평형서는 그
FWHM($\approx91$ mV)이 간격($100$ mV)보다
작아 \emph{겨우 분리}된다. 전류를 키워 꼬리 $L_{V,j}$ 가 $\sim100$ mV 로 자라면(조건 (ii)) $j$ 의 꼬리가 $j{+}1$
peak 에 닿아 \emph{융합이 시작}된다. $L_{V,j}\propto|I|\,e^{\Delta H_a/RT}$ 라 저온일수록 더 \emph{낮은} 전류에서
융합이 개시된다 — 관측의 ``저온$\cdot$고율서 peak 가 겹친다''가 이 부등식의 등호로 정량화된다. 일단 융합하면 두
전이의 면적이 한 봉우리에 섞여, 저율 anchor 없이는 $Q_j$ 를 되돌릴 수 없다(아래).
\end{verifybox}

\textbf{OCV-anchored forward 합산$\cdot$예측(절차).} 융합한 고율 데이터에서 $j$별 기여를 \emph{역산}하면 비유일이다(\S\ref{sec:dist}). 대신 저율 정보를 anchor 로 \emph{forward} 예측한다:
\begin{enumerate}[label=\textbf{S$_\arabic*'$.},leftmargin=2.8em,itemsep=2pt]
\item \textbf{저율 OCV 로 분리 peak 에서 $\{U_j,w_j,Q_j\}_{j=1}^{N_p}$ 고정}(식~\eqref{eq:eqpeak}; \S\ref{sec:synth}
  의 S1 과 동일). 저율서는 꼬리가 짧아 peak 가 분리되므로 전이별 평형 파라미터가 식별된다.
\item \textbf{각 전이의 고율 단독 곡선 생성}: 고정 $\{U_j,w_j,Q_j\}$ 에 동역학 $\{\Delta H_{a,j},\chi\}$ 와 Arrhenius
  \emph{절편 조합값}(\S\ref{sec:synth} S3$\cdot$S4 — $\Delta S_a$ 단독이 아니라 흡수 prefactor 와 결합된 절편이
  $L_{q,j}$ 절대값을 정함)을 넣어 $j$별 평형 rise $+$ 지수 꼬리(식~\eqref{eq:simplefit})를 그린다.
\item \textbf{합산} 식~\eqref{eq:total} 로 \emph{융합 형상을 예측}한다($\sum_j$).
\item \textbf{데이터 대조}: 예측 융합 곡선을 고율 측정과 맞춰 동역학 파라미터를 정련한다 — $\{U_j,w_j,Q_j\}$ 는
  S$_1'$ 고정값 유지(자유화하면 공선형).
\end{enumerate}
\textbf{분리 가능성의 한계.} peak 가 완전히 융합하면 $j$별 면적 $Q_j$ 의 분리는 고율 데이터만으로는 불가능하다.
이는 \S\ref{sec:synth} 공선형 논리의 겹침판이다. 저율 OCV anchor 가 있을 때만 융합 peak 의 합산 정보를 전이별로
되돌릴 수 있다. 그래서 본 절은 ``저율로 고정, 고율로 예측''의 forward 방향만 허용하며, 이는 \S\ref{sec:dist}
의 비유일성 결론과 정합한다. 한 가지 확장을 미리 적어 두면 — 본 절의 합산은 단상$\cdot$임계($\Omega_j\le2RT$,
전이당 peak 하나)를 전제로 적었으나, $\Omega_j>2RT$ 전이는 충$\cdot$방전 두 중심 $U_j^{\dis/\chg}$ 의 peak 쌍
(\S\ref{sec:hys_peak})이 되므로 $\sum_j$ 가 방향별 중심 각각으로 확장된다. 통합식~\eqref{eq:hys_master} 의
$\sum_j$ 가 그 겹침까지 포함한다. 이로써 여러 전이의 합산이 닫혔다. 방전 본론의 마지막 조각으로, 다음 절은 같은
측정 곡선을 역미분 좌표 $\dd V/\dd Q$ 로 읽는 상보적 관점을 정리한다.

% ====================================================================
\section{보완 관점: $\dd V/\dd Q$ (DVA)}\label{sec:dva}
% ====================================================================
지금까지 $\dd Q/\dd V$(ICA)를 세웠다. 같은 측정 곡선 $Q(V)$ 의 역미분 $\dd V/\dd Q$ — differential voltage
analysis, DVA — 는 같은 물리를 상보적 좌표로 드러낸다. 새 파라미터는 없으며, 같은 $\{U_j,w_j,Q_j\}$ 가 다른
형태로 나타날 뿐이다. 본 절은 둘의 대응 관계와 둘을 함께 쓰는 이유를 정리한다. 새 모델이 아니다.

\subsection{peak $\leftrightarrow$ valley 대응}
정의상 점별로 $\dd V/\dd Q=(\dd Q/\dd V)^{-1}$ 다(다전이 합 식~\eqref{eq:total} 의 점별 역수). 따라서
\begin{equation}
\frac{\dd V}{\dd Q}(V)=\Big[\,C_\bg(V_n,T)+\sum_{j}Q_j\big(\text{평형 상승부}+\text{꼬리}\big)\Big]^{-1},
\label{eq:dva}
\end{equation}
이고 두 표현의 극점이 \emph{서로 뒤집힌다}: $\dd Q/\dd V$ 가 \emph{peak}(전이 plateau — 전하가 많고 전위가 평탄)인
$V=U_j$ 에서 $\dd V/\dd Q$ 는 \emph{최소(valley)}가 되고, 전이 \emph{사이}(stage 경계 — $\dd Q/\dd V$ 가 최소,
전위가 급변)에서 $\dd V/\dd Q$ 가 \emph{최대(peak)}가 된다. \textbf{즉 DVA peak 는 stage \emph{경계}를, ICA peak 는
stage \emph{plateau}를 표시}한다 — 같은 전이의 양쪽 정보다.

\subsection{왜 둘을 함께 쓰나}
\begin{enumerate}[label=(\roman*),leftmargin=2.0em,itemsep=2pt]
\item \textbf{기저선 분리 — 더 많은 정보가 아니다}: $\dd V/\dd Q$ 와 $\dd Q/\dd V$ 는 전단사 관계라
  정보량이 같고, 어느 좌표도 본질적으로 우월하지 않다. 역수는 오히려 $\dd Q/\dd V$ 가 작은 영역에서 잡음을
  $1/y^2$ 로 증폭한다. 실무 이점은 분산 감소가 아니라 표현에 있다. plateau 에서 ICA peak 는 배경 $C_\bg$ 와
  겹쳐 기저선 분리가 어려운 반면, DVA 는 stage 경계의 급변을 형상으로 도드라지게 해 경계 SOC 의 판독을 쉽게
  한다. 분산의 정량 비교 자체는 잡음 모형에 의존하는 문제라 본 장 범위 밖이다.
\item \textbf{용량의 직접성}: ICA peak \emph{면적}은 곧 $Q_j$(전이 용량)다 — DVA 의 vs-$V$ 좌표에선 valley 폭으로
  간접적이다. 다만 $\dd V/\dd Q$ 를 vs-$Q$ 로 그리면 경계 간 $\Delta Q$ 에서 $Q_j$ 를 직독할 수 있으므로,
  이는 좌표 선택의 문제다.
\item \textbf{교차검증}: \emph{DVA valley} 수 $=$ ICA \emph{peak} 수 $=N_p$(둘 다 $U_j$ 표시; DVA \emph{peak}=경계는 전이
  \emph{사이}라 별개로 셈), \emph{DVA valley 위치 $=$ ICA peak 위치 $=U_j$}. 두 좌표는 전단사라 \emph{독립 정보가 아니므로
  잡음을 통계적으로 줄이진 못한다} — 다만 같은 데이터를 두 표현에서 읽어 \emph{판독$\cdot$기저선 처리의 \textbf{절차 오류}}
  (peak 누락$\cdot$경계 오지정 등)를 교차 점검하는 데 그친다(비통계적).
\end{enumerate}
노화 진단에서 \emph{DVA peak 간격}의 변화가 전극별 용량 손실(LLI$\cdot$LAM)을 분리하는 표준 도구인 것도 같은
이유다 \cite{dubarry2012,bloom2005} — 경계의 급변이 \emph{형상}으로 도드라져 간격 변화의 판독이 쉽기 때문이다(통계적
우월이 아니라 표현상 이점, 위 (i)).

\subsection{같은 파라미터, 같은 의존}
$\dd V/\dd Q$ 의 valley 폭$\cdot$깊이도 \emph{새 물리가 아니라} 같은 모델에서 온다: valley 깊이는
$1/[C_\bg+Q_j/(4w_j)]$ ($C_\bg$ 포함), 폭은 $w_j$(단상이면 식~\eqref{eq:weff}, 분기 영역이면 \S\ref{sec:hys_peak} 의 현상학적 폭), 비대칭은
식~\eqref{eq:Lq} 동역학 꼬리, 충$\cdot$방전 갈림은 \S\ref{sec:hys_branch} 분기. \emph{단 valley 의 단일-$w_j$ 귀속은 전이
\textbf{비겹침} 영역에 한한다} — 겹치면 합의 역수라 전이별로 안 쪼개진다(\S\ref{sec:overlap} forward-only 와 같은 제약).
또 폭(\S\ref{sec:regsol})과 gap 의 $\Omega_j$ 는 같은 $g_j''$ 의 사영이라 \emph{독립이 아닌 내적 일관성}이다(\S\ref{sec:falsify}).
그러므로 본 장의 피팅은 $\dd Q/\dd V$ 로 하고, $\dd V/\dd Q$ 는 경계 SOC 의 판독을 쉽게 하는 데 병용한다.
새 정보가 아니라 같은 모델식의 두 그림이다.

\medskip
\textbf{방전 본론의 종결.} 이로써 방전 $\dd Q/\dd V$ 의 이론이 완결되었다. 내부 전위의 정의(\S\ref{sec:stage})
에서 출발해 평형 기준선(\S\ref{sec:eqpeak}$\cdot$\S\ref{sec:regsol})과 동역학 꼬리(\S\ref{sec:lag}$\cdot$
\S\ref{sec:barrier}), 입자 분포(\S\ref{sec:dist})를 지나, 닫힌식과 식별 절차(\S\ref{sec:synth}), 다전이 합산
(\S\ref{sec:overlap}), 상보 좌표(본 절)까지 닫았다. 이후의 절들은 이 본론 위에 세우는 \emph{확장}이다.
\S\ref{sec:regsol} 에서 이미 등장한 상분리가 충전과 방전을 가르는 분기를 만들기 때문에, 다음 절부터 그
히스테리시스를 같은 골격 위에 얹는다(\S\ref{sec:hys_branch}--\S\ref{sec:hys_partial}). 마지막으로
\S\ref{sec:master} 가 방전 한 줄 피팅식과 그 양방향 통합형을, 측정 파일에서 파라미터까지 가는 피팅 알고리즘과
함께 닫는다.

% ====================================================================
\section{충방전 분기와 히스테리시스 gap $\Delta U_j^\hys$}\label{sec:hys_branch}
% ====================================================================
여기서부터 확장이 시작된다. 방전 본론(\S\ref{sec:stage}--\S\ref{sec:dva})은 한 방향의 $\dd Q/\dd V$ 를
닫았다. 이제 같은 전극을 충전과 방전으로 오갈 때, 같은 SOC 에서 전위가 갈리는 \emph{충방전 히스테리시스}로
넓힌다. 새 물리는 거의 필요하지 않다. 그 뿌리는 \S\ref{sec:regsol} 에서 이미 세운 상분리이기 때문이다 —
충전과 방전이 서로 다른 준안정 분기를 따라 spinodal 까지 과주행하기에 두 방향의 전위가 갈린다. 본 절은 그
갈림(gap)을 정규용액 $\Omega_j$ 와 온도의 닫힌 식으로 유도한다.

\subsection{비단조 평형 전위 $V_{\eq,j}(\xi)$}
\S\ref{sec:regsol} 의 상분리를 \emph{전위}의 언어로 옮긴다. 평형 조건 식~\eqref{eq:eqcond}
$\mu_\mathrm{Li}(\theta_{\eq,j})=\mu^0-sF(V_n-U_j)$ 에 화학퍼텐셜 식~\eqref{eq:mu}(이제 $\Omega_j\ne0$)를 넣고
$\theta=1-\xi$ 로 바꾸면 — $\ln[\theta/(1-\theta)]=-\ln[\xi/(1-\xi)]$, $(1-2\theta)=-(1-2\xi)$ 라 두 항 부호가 함께
뒤집혀 —
\begin{equation}
V_{\eq,j}(\xi)=U_j+\frac{RT}{sF}\ln\frac{\xi}{1-\xi}+\frac{\Omega_j}{sF}(1-2\xi).
\label{eq:hys_Veq}
\end{equation}
$\Omega_j=0$ 이면 둘째 항만 남아 앞 절 logistic 식~\eqref{eq:logistic} 의 역함수(단조)다. $\Omega_j$ 항이
\emph{되돌림}을 더한다. 기울기는
\begin{equation}
\frac{\dd V_{\eq,j}}{\dd\xi}=\frac{1}{sF}\Big[\frac{RT}{\xi(1-\xi)}-2\Omega_j\Big]=\frac{1}{sF}\,g_j''(\xi),
\label{eq:hys_slope}
\end{equation}
즉 \emph{전위 곡선의 기울기가 곧 \S\ref{sec:regsol} 의 자유에너지 곡률}(식~\eqref{eq:regsol_gpp})이다 —
$V_{\eq,j}$ 기울기 $0$ 인 곳이 spinodal($g_j''=0$). $\Omega_j>2RT$ 면 $\xi=\tfrac12$($RT/[\xi(1-\xi)]=4RT$) 둘레서
음의 기울기(역학적 불안정)다. 기울기 $0$ 인 spinodal 진행률은 식~\eqref{eq:hys_slope}$=0$ 에서
$\xi(1-\xi)=RT/(2\Omega_j)$, 곧
\begin{equation}
\boxed{\;\xi_{s,j}^\pm=\tfrac12\pm\tfrac12\sqrt{1-\tfrac{2RT}{\Omega_j}}\;,\qquad(\Omega_j>2RT).\;}
\label{eq:hys_spinodal}
\end{equation}
\begin{verifybox}
\textbf{임계$\cdot$환원.} 실수 spinodal 은 $\Omega_j>2RT$($T<T_{c,j}=\Omega_j/2R$, \S\ref{sec:regsol})에서만. $\Omega_j
\to2RT^+$ 면 $\xi_{s,j}^\pm\to\tfrac12$, $\Omega_j\to\infty$ 면 $\to0,1$. $\Omega_j\le2RT$ 면 $V_{\eq,j}$ 단조라
spinodal$\cdot$히스 없음 — 앞 절 단일 peak.
\end{verifybox}

\subsection{충방전 분기와 gap}
비단조 $V_{\eq,j}$ 는 $\xi_{s,j}^-$ 서 \emph{극대}, 불안정 구간서 하강, $\xi_{s,j}^+$ 서 \emph{극소}, 다시 상승한다.
한 전이는 불안정 구간을 균질하게 못 지나 \S\ref{sec:regsol} 의 준안정 분기를 따른다. \emph{방향이 출발 어깨를
고정}한다: \textbf{방전}(탈리튬화)은 $\xi\!\approx\!0$ 의 안정 단상에서 출발해 $\xi$ 를 키우며 \emph{상승} 준안정 가지를
타고, 핵생성 장벽이 사라지는 \emph{첫} 극값 $\xi_{s,j}^-$(상승 가지의 spinodal, \S\ref{sec:regsol})에서 비로소 분리한다.
\textbf{충전}(리튬화)은 반대로 $\xi\!\approx\!1$ 에서 내려와 $\xi_{s,j}^+$ 에서 분리한다 — 같은 핵생성 물리라도 \emph{진행
방향}이 도달하는 spinodal 어깨를 가르므로 두 분기가 갈린다. 분기의 \emph{상한} 중심 전위는 그 극값이고(아래
$\gamma_j$ 로 축소된 실측 중심은 식~\eqref{eq:hys_center})
\begin{equation}
U_j^{\dis,\sup}=V_{\eq,j}(\xi_{s,j}^-),\qquad U_j^{\chg,\sup}=V_{\eq,j}(\xi_{s,j}^+),
\label{eq:hys_branchU}
\end{equation}
극대$>$극소라 $U_j^{\dis,\sup}>U_j^{\chg,\sup}$ — 방전 전위가 충전보다 높다는 관찰과 부호가 맞는다(부호는 $s=+1$
규약 하 양수, $s$ 는 $\Delta U_j^\hys$ 의 $2/sF$ 에 흡수). gap 은 두 극값 차다.
$\xi_{s,j}^\pm=\tfrac12\pm\tfrac12 u_j$($u_j\equiv\sqrt{1-2RT/\Omega_j}$)를 식~\eqref{eq:hys_Veq} 에 넣어 빼면 —
방전 끝점서 $\tfrac{\xi}{1-\xi}=\tfrac{1-u_j}{1+u_j}$, 충전서 $\tfrac{1+u_j}{1-u_j}$, $(1-2\xi)=\pm u_j$ — 로그 항
차 $2\ln\tfrac{1-u_j}{1+u_j}=-4\,\mathrm{artanh}\,u_j$, $\Omega_j$ 항 차 $2u_j$ 라
\begin{equation}
\boxed{\;\Delta U_j^\hys=U_j^{\dis,\sup}-U_j^{\chg,\sup}=\frac{2}{sF}\big[\Omega_j u_j-2RT\,\mathrm{artanh}\,u_j\big],\quad
u_j=\sqrt{1-\tfrac{2RT}{\Omega_j}}\;.}
\label{eq:hys_dU}
\end{equation}
이로써 히스테리시스 gap 이 \emph{정규용액 $\Omega_j$ 와 온도에서 유도}됐다 — 임의 피팅 상수가 아니다.
\begin{verifybox}
\textbf{극한$\cdot$수치.} $\Omega_j\to2RT^+$ 면 $\Omega_j u_j-2RT\,\mathrm{artanh}\,u_j\to\tfrac43RT u_j^3$ 라
$\Delta U_j^\hys\to(8RT/3sF)u_j^3\to0^+$(임계서 매끄럽게 소멸). $\Omega_j=4RT$ 면 $u_j=\sqrt{0.5}\approx0.707$,
$\mathrm{artanh}(0.707)\approx0.881$, $\Delta U_j^\hys\approx(2/sF)(4{\cdot}0.707-2{\cdot}0.881)RT\approx2.1\,RT/(sF)
\approx54$ mV($25^\circ$C). 차원: $\Omega,RT$ J/mol, $sF$ C/mol $\Rightarrow$ V.
\end{verifybox}
\begin{verifybox}
\textbf{Worked example — gap 절편의 온도 의존이 보이는 크기.} S5 가 회귀하는 절편$(T)=\gamma_j\Delta U_j^\hys(T)$
의 감각을 위해 $\gamma_j=0.6$ 으로 두고, $\Omega_j$ 를 $25^\circ$C 의 $RT_0$ 배수로 잡아 측정창 세 온도에서
절편을 계산하면 다음과 같다(mV).
\begin{center}
\renewcommand{\arraystretch}{1.25}
\begin{tabular}{lcccc}
\toprule
$\Omega_j$ & $T_{c,j}$ & $15^\circ$C & $23^\circ$C & $45^\circ$C \\
\midrule
$3RT_0$ & $447$ K & $14.2$ & $13.1$ & $10.2$ \\
$4RT_0$ & $596$ K & $34.7$ & $33.2$ & $29.3$ \\
$8RT_0$ & $1193$ K & $135.2$ & $132.8$ & $127.0$ \\
\bottomrule
\end{tabular}
\end{center}
창 안($15\to45^\circ$C) 상대 감소는 각각 $-28\%$, $-16\%$, $-6\%$ 다. $T_{c,j}$ 가 측정창에 가까울수록
절편의 온도 의존이 가파르고, 깊은 2상측($8RT_0$)은 $-6\%$ 라 측정 불확도에 묻히기 쉽다 — 피팅 알고리즘
(6$'$)의 퇴화 분기(lumped $\Delta_j$ 강등)가 필요한 까닭이 이 수치에서 그대로 보인다.
\end{verifybox}
\begin{boundbox}
\textbf{$U_j$-대칭과 축소 $\gamma_j$.} 두 spinodal 극값은 $U_j$ 에 \emph{대칭}이다 — 식~\eqref{eq:hys_Veq} 의 로그항과
$\Omega$ 항이 모두 $\xi=\tfrac12$ 대칭이라 $\tfrac12[V_{\eq}(\xi_s^-)+V_{\eq}(\xi_s^+)]=U_j$(끝점 $\xi_s^\pm=\tfrac12\pm
\tfrac12 u_j$ 에서 로그항 합 $\ln\tfrac{1-u}{1+u}+\ln\tfrac{1+u}{1-u}=0$, $\Omega$ 항 합 $(+u)+(-u)=0$). 이 대칭은
\emph{가정이 아니라 결과}이며(비대칭 site energy 면 깨짐), 그래서 분기 중심을 $U_j$ 둘레 $\pm\tfrac12$ 로 쪼갤 수 있다.
식~\eqref{eq:hys_branchU} 의 spinodal gap 은 각 분기가 끝까지 과주행하는 \emph{상한}이고, 실제로는 핵생성이 spinodal
전에 시작되고 다입자 평활화(\S\ref{sec:regsol} Dreyer)가 갈림을 좁혀 실측은 그 아래다. 이를 \emph{현상학적} 한 자유도
$\gamma_j\in[0,1]$ 로 흡수해
\begin{equation}
U_j^{\dis/\chg}=U_j\pm\tfrac12\gamma_j\Delta U_j^\hys
\label{eq:hys_center}
\end{equation}
로 둔다. $\gamma_j$ 의 선형 스케일링은 1차 근사다. $\gamma_j=1$ 은 spinodal 상한이고, $\gamma_j\to0$ 은
binodal 잔류 plateau 가 아니라 히스 소멸을 가리키는 규약적 하한이다. 다중 전환점의 엄밀한 취급은 이력을 전환점
분포의 중첩으로 쓰는 Preisach 류 모형을 요구하므로 본 장 범위 밖이다(\S\ref{sec:hys_partial}). 정량에는
$\Omega_j$ 와 $\gamma_j$ 한 양씩이면 충분하다. \textbf{식별 주의.} 관측 gap 의 \emph{절편} — 여러 율의 충방전 gap $\Delta V_{\mathrm{obs},j}(|I|)$ 를
$|I|\to0$ 으로 외삽한 값(\S\ref{sec:hys_pol}) — 은 $\gamma_j\Delta U_j^\hys$ 의 곱이므로, 한 온도의 절편만으로는
$\gamma_j$ 와 $\Omega_j$ 가 분리되지 않는다. 1차 식별 경로는 절편의 온도 의존이다. 식~\eqref{eq:hys_dU} 의
$\Delta U_j^\hys(T)$ 형이 $\Omega_j$ 를 정하고, 크기가 $\gamma_j$ 를 정한다. 이는 \S\ref{sec:master} 의
S5$\cdot$파라미터 표와 동일한 절차이며, $\gamma_j$ 가 측정창에서 $T$-무관이라는 전제 위에 선다. 단상 온도
구간이 있으면 평형 폭 $w^\eff(T)$(식~\eqref{eq:weff})가 검증 경로가 된다. 다만 창 전체에서 히스가 보이는
전이는 그 창이 전부 2상측($\Omega_j>2RT$, $T<T_{c,j}$)이라 폭-경로가 무효임에 유의한다 — 폭과 gap 은 같은
$g_j''$ 의 사영이라 어차피 내적 일관성 점검용이다(\S\ref{sec:falsify}(iii)).
\end{boundbox}

% ====================================================================
\section{분기별 평형 $\dd Q/\dd V$ 와 peak 쌍}\label{sec:hys_peak}
% ====================================================================
앞 절이 분기 중심 $U_j^{\dis/\chg}$(식~\eqref{eq:hys_center})를 열역학으로 닫았다. 이 절은 그것을 관측량
$\dd Q/\dd V$ 로 옮긴다. \S\ref{sec:eqpeak} 의 평형 peak 는 중심 $U_j$ 의 단상 logistic
식~\eqref{eq:logistic} 이었고, 그 형상은 중심값만 바꿔 그대로 쓴다. 충$\cdot$방전 각 방향에서 전이 $j$ 의
상승부는 자기 분기 중심에 선다:
\begin{equation}
\frac{\dd Q}{\dd V_n}\Big|^{\dis}=C_\bg+\sum_j Q_j\,\frac{\dd\xi_{\eq,j}}{\dd V_n}\Big|_{U_j^\dis},\qquad
\frac{\dd Q}{\dd V_n}\Big|^{\chg}=C_\bg+\sum_j Q_j\,\frac{\dd\xi_{\eq,j}}{\dd V_n}\Big|_{U_j^\chg},
\label{eq:hys_dQdV}
\end{equation}
여기서 $\dd\xi_{\eq,j}/\dd V_n$ 는 폭 $w_j$ 의 종 모양이고, 중심만 $U_j\to U_j^{\dis/\chg}$ 로 옮겼다.

\textbf{이 종 모양의 지위 — 왜 plateau 가 유한 폭으로 보이는가.} 형상의 출처를 단계별로 분명히 해 둘 필요가
있다. 분기는 $\Omega_j>2RT$ 에서만 존재하고, 바로 그 영역에서 \emph{평형} 등온선은 logistic 이 아니라 상분리
plateau 다(\S\ref{sec:regsol}; 평형 유효 폭 $w^\eff$ 는 음이 되어 무의미하다). 단일 입자의 평형 그림이라면
plateau 의 $\dd Q/\dd V$ 는 무한히 날카로운 선이어야 한다. 그런데 관측 peak 는 유한 폭의 종 모양이다. 그 폭은
두 군데서 온다. 첫째, 실제 전극은 입자 집단이고 입자마다 분기 전위가 조금씩 달라(\S\ref{sec:regsol} 의 Dreyer
다입자 그림) 날카로운 선이 입자 분포의 폭만큼 번진다. 둘째, 동역학 지연(\S\ref{sec:lag})이 비대칭 꼬리를
더한다. 따라서 식~\eqref{eq:hys_dQdV} 의 종 모양은 평형형이 아니라 관측되는 broadened peak 의 \emph{현상학적
형상 ansatz}이며, 그 유한 폭 $w_j$ 는 평형 예측이 아니라 \emph{독립 피팅 폭}이다(\S\ref{sec:master} 통합형의
주와 동일). 같은 기호 $w_j$ 의 지위가 두 영역에서 다르다는 점을 명심해야 한다 — 단상($\Omega_j\le2RT$)에서는
식~\eqref{eq:weff} 의 평형 예측(검증 가능)이고, 분기 영역($\Omega_j>2RT$)에서는 broadening 이 정하는 피팅
자유도다. 이 단서 위에서 한 전이는 ICA 에 \textbf{peak 쌍}으로 나타나며, 그 쌍의 성질은 다음 셋으로 정리된다.
\begin{itemize}[leftmargin=1.4em,itemsep=2pt]
\item \textbf{중심 갈림} $\Delta V_{\mathrm{peak},j}=U_j^\dis-U_j^\chg=\gamma_j\,\Delta U_j^\hys$ (식~\eqref{eq:hys_center}) —
  \emph{평형 히스테리시스}의 직접 관측량,
\item \textbf{면적} $Q_j$ 는 양방향 동일(같은 전이가 지나는 \emph{총 전하}는 물질량 보존으로 유지 — \S\ref{sec:eqpeak}
  대로 상분리서 $\int\dd\xi=1$ 적분형은 plateau 로 대체되나 $Q_j$ 자체는 불변),
\item \textbf{모양$\cdot$폭}은 같은 (현상학적) $w_j$ (양방향 같은 broadening).
\end{itemize}
정리하면, 충방전 $\dd Q/\dd V$ 를 겹쳐 그리면 면적과 폭이 같고 중심만 $\gamma_j\Delta U_j^\hys$ 어긋난 거울
쌍이 전이마다 보인다. 다만 이는 아직 평형 그림이다 — 실측 gap 에는 전류가 만드는 분극이 더해져 있으므로, 다음
절이 그 동역학 몫을 분리해 열역학 gap 만 남기는 절차를 세운다.

% ====================================================================
\section{동역학 분극과 열역학 gap 의 분리}\label{sec:hys_pol}
% ====================================================================
식~\eqref{eq:hys_center} 의 분기 갈림은 \emph{평형}($|I|\to0$)에서도 남는 \textbf{열역학} 성분이다. 실측 충방전 gap 은
거기에 전류가 만드는 \textbf{동역학 분극}이 더해진 것이다. 전압 분해 식~\eqref{eq:vapp}
$V_\app=V_n+s_I|I|R_n$ 에서, 방전($s_I=+1$)은 음극 전위를 $|I|R_n$ 만큼 올리고 충전은 내리므로 관측 peak 중심은
\begin{equation}
V_{\mathrm{peak},j}^{\dis}=U_j^\dis+|I|R_n,\qquad V_{\mathrm{peak},j}^{\chg}=U_j^\chg-|I|R_n,
\label{eq:hys_obspeak}
\end{equation}
\begin{equation}
\boxed{\;\Delta V_{\mathrm{obs},j}(|I|)=\underbrace{\gamma_j\,\Delta U_j^\hys}_{\text{열역학}(|I|\to0\text{ 잔존})}
+\underbrace{2|I|R_n}_{\text{동역학}(|I|\to0\text{ 소멸})}.\;}
\label{eq:hys_obsgap}
\end{equation}
\textbf{이 분리가 본 장의 실용 핵심이다.} 여러 율에서 gap 을 재어 $|I|\to0$ 으로 외삽하면, 절편이 열역학 gap
$\gamma_j\Delta U_j^\hys$(\S\ref{sec:regsol} 의 $\Omega_j$ 와 $\gamma_j$)이고 기울기가 $2R_n$ 이다. 이 가산
분해 — 절편 더하기 기울기 — 는 $\gamma_j$ 가 전류에 무관할 때 성립한다. 고율에서 과주행 정도가 바뀌어
$\gamma_j=\gamma_j(|I|)$ 가 되면 절편과 기울기가 섞여 분리할 수 없다. $\gamma_j$ 평탄 가정의 물리 조건은
``핵생성$\cdot$전환 시간상수가 구동 시간 $1/$율보다 훨씬 짧다''이다. 그 조건에서 과주행은 핵생성 장벽이 정하는
율-무관 수준($\gamma_j<1$, spinodal 아래)에 고정된다. 빠른 핵생성은 작은 과주행으로의 포화를 뜻하지 spinodal
도달을 뜻하지 않는다 — 앞 절 boundbox 와 정합한다. 가정이 위배되는 고율에서는 과주행이 율을 따라 더 밀리므로,
위배는 gap--$|I|$ 곡선의 곡률로 검출된다. 따라서 외삽은 저$\sim$중율 구간에 한정하고, 분극이 peak 위치만 옮기며
형상은 \S\ref{sec:lag} 꼬리로 따로 처리된다는 전제를 둔다.
$R_n$ 은 ohmic$\cdot$전하이동$\cdot$확산을 한 양으로 묶은 lumped 저항이다(식~\eqref{eq:vapp} 의 공통 $R_n$).
공통 $R_n$ 의 채택은 검증 가능한 가정이다. 그 내용은 ohmic 과 SEI 성분이 전이 간에 공유되고, 전하이동과 확산의
전이별$\cdot$SOC별 차이가 측정 창 안에서 2차라는 것이다. 이 가정은 도피처가 아니라 falsifiable 예측을 낳는다 —
모든 전이$\cdot$SOC 에서 $2|I|R_n$ 기울기가 같아야 하므로, 전이별 gap 기울기가 유의하게 갈리거나 같은 전이에서
IR-drop 이 SOC 에 의존하면 기각되고, 그때는 전이별$\cdot$SOC별 $R_{n,j}(\xi)$ 로 간다. staging 별
$D_{\mathrm{Li}}$ 차가 크면 위배가 나타난다. 또 하나, $2|I|R_n$ 의 선형성은 ohmic 근사다. charge-transfer
과전위는 $\eta_\mathrm{ct}\propto\ln|I|$ 로 비선형이라(\S\ref{sec:falsify}) 저$\sim$중율에서는 $R_n$ 에
흡수되지만, 고율에서 gap--$|I|$ 곡선의 기울기가 줄며 아래로(오목하게) 휘면 선형 근사가 기각된다. $\ln$ 포화
곡률은 charge-transfer 우세의 신호이고, 반대로 위로 휘는 것은 확산 한계 또는 $\gamma_j(|I|)$ 쪽 신호다. 이
율 의존은 \S\ref{sec:lag} 의 지연 꼬리와 같은 뿌리에서 나온다.
\begin{flagbox}
\textbf{왜 분극과 히스를 구분하나.} 분극만 있는 계(상분리 없음, $\Omega_j\le2RT$)는 $|I|\to0$ 에서 gap 이 \emph{닫힌다}.
흑연 staging 은 닫히지 않는 잔류 gap 을 보이며, 그 절편이 곧 상분리($\Omega_j>2RT$)의 증거다. gap 의 \emph{온도}
의존($\Delta U_j^\hys(T)$ 식~\eqref{eq:hys_dU}, $T\to T_{c,j}$ 서 소멸)도 같은 결론을 교차 검증한다.
\end{flagbox}

% ====================================================================
\section{부분 cycle 과 return-point 기억}\label{sec:hys_partial}
% ====================================================================
지금까지는 완전 충방전을 가정했다. 중간 SOC 에서 방향을 바꾸는 부분 cycle 에서는 한 전이가 분기를 끝까지 돌지
못해 \emph{minor loop} 를 그린다. 이때의 핵심 관측이 \textbf{return-point memory} 다 — 방향을 되돌리면 계가
직전 전환점으로 되돌아와 큰 loop 에 합류한다. 이는 상분리 1차 전이의 일반적 성질이다~\cite{sethna1993}. 본 장은
이를 전이별 보간 인자 하나로 최소 기술한다. 전환점 $\xi_\mathrm{turn}$ 의 분기 내 정규화 위치를
$\eta\equiv(\xi_\mathrm{turn}-\xi_{s,j}^-)/(\xi_{s,j}^+-\xi_{s,j}^-)\in[0,1]$ 로 두면, $\eta$ 만큼만 진행한
부분 cycle 의 유효 gap 은
\begin{equation}
\Delta U_{j}^{\hys,\,\mathrm{part}}=h_j(\eta)\,\Delta U_j^\hys,\qquad h_j(0)=0,\;h_j(1)=1,
\label{eq:hys_partial}
\end{equation}
로 줄어든다. $h_j$ 는 단조 증가 함수이고, 데이터가 없으면 무정보 기본형 $h_j(\eta)=\eta$ 를 쓴다 — 이는 물리
유도가 아니라 최대-무지 prior 이며, 곡률은 부분 cycle 데이터로 보정한다. 완전 cycle 만 다루면 $h_j=1$ 이라 앞
절들이 그대로 성립하고, 부분 cycle 데이터가 있을 때 $h_j$ 가 major-loop 사이의 보간을 준다. 한계도 정직하게
명시해 둔다. 단일 스칼라 $h_j(\eta)$ 는 minor loop 의 gap \emph{크기}만 근사한다. return-point memory 의
본질인 경로 의존 — 같은 $\xi$ 라도 어느 전환점에서 왔는가, 중첩된 nested loop — 은 담지 못하며, 그것은 전환점
분포를 쓰는 Preisach 류 모형을 요구하므로 본 장 범위 밖이다. 즉 ``return-point 와 $h_j$ 면 충분하다''는 말은
gap 크기의 보간에 한한 것이고, 다중 전환점 경로가 중요해지면 $h_j$ 단일 인자의 한계가 드러난다. 그 위배 자체가
진단 신호가 된다. 이로써 확장의 마지막 조각인 $h_j$ 까지 닫혔다. 다음 절이 방전식과 분기, 부분 cycle 을 하나의
통합 피팅식으로 모은다.

% ====================================================================
\section{종합 모델식과 피팅$\cdot$예측 — 한 줄로}\label{sec:master}
% ====================================================================
앞 절들이 \emph{각 절은 그 식의 한 항}이라 했다. 이제 그 항들을 하나로 모은다. 먼저 \emph{방전} $\dd Q/\dd V$ 를
\emph{(배경) $+$ 전이별 (평형 상승부 $+$ 동역학 꼬리)}의 합으로 닫고, 이어 \emph{분기 중심과 분극 부호만} 방향별로
바꿔 \emph{충방전 양방향}으로 넓힌다 — 이것이 본 장이 도달하는 \emph{단일 통합 피팅식}이다.

\subsection{최종 모델식}
측정 전위 $V_\app$ 를 내부 전위 $V_n=V_\app-s_I|I|R_n$(식~\eqref{eq:vapp})으로 환산한 뒤, 전이 $j=1,\dots,N_p$ 에 대해
\begin{equation}
\boxed{\;\frac{\dd Q}{\dd V_\app}=C_\bg(V_n,T)+\sum_{j=1}^{N_p}Q_j\bigg[\underbrace{\frac{\xi_{\eq,j}(1-\xi_{\eq,j})}{w_j}}_{\text{평형 상승부}}
\;+\;\underbrace{\frac{r_{a,j}}{L_{V,j}}\,e^{-(V_n-V_{a,j})/L_{V,j}}}_{\substack{\text{동역학 꼬리}\\ s(V_n-V_{a,j})\ge0}}\bigg]\;.}
\label{eq:master}
\end{equation}
각 항은 앞 절의 \emph{닫힌} 결과를 그대로 부른다:
\begin{align*}
\xi_{\eq,j}&=\big[1+e^{-s(V_n-U_j)/w_j}\big]^{-1} &&\text{(식~\eqref{eq:logistic}; 평형 상승부 → \S\ref{sec:eqpeak})},\\
L_{V,j}&=\Big|\tfrac{\dd V_n}{\dd q}\Big|_{q_a}L_{q,j},\quad
L_{q,j}=\tfrac{|I|}{Q_\cell}\tfrac{h}{k_BT}\,e^{(\Delta H_{a,j}-T\Delta S_{a,j}-\chi\mathcal A_j)/RT}
&&\text{(식~\eqref{eq:Lq}\eqref{eq:keff}; \S\ref{sec:lag}\S\ref{sec:barrier})},\\
\mathcal A_j&=sF(V_n-U_j),\quad r_{a,j}=\xi_{\eq,j}(q_a)-\xi_j(q_a)&&\text{(식~\eqref{eq:affinity}; 구동력$\cdot$잔류 지연)}.
\end{align*}
$V_{a,j}=V_n(q_{a,j})$=전이 $j$ 의 꼬리 시작 전위. 분포가 좁으면($\rho_G\to\delta$) 이 단일지수 꼬리가 정확하고,
넓어 꼬리가 휘면(stretched) 그 꼬리 항을 식~\eqref{eq:superpose} 의 $\rho_G$ 적분으로 바꾼다(\S\ref{sec:dist}).
\emph{이 한 식이 peak 의 위치$\cdot$너비$\cdot$높이와 그 온도$\cdot$C-rate$\cdot$전위 의존(3$\times$3, \S\ref{sec:synth})을 모두 담는다.}

\subsection{충방전 통합형 — 한 식이 양방향을 담는다}
식~\eqref{eq:master} 는 \emph{방전} 분기다. 충전까지 한 식으로 닫으려면 두 가지만 바꾼다: (i) 전이 중심을 분기 중심
$U_j^{\dis/\chg}$(식~\eqref{eq:hys_center})으로, (ii) 분극 부호를 방향에 따라(식~\eqref{eq:hys_obspeak}). 방향
$d\in\{\dis,\chg\}$, 부호 $\sigma_d=+1\,(\dis)/-1\,(\chg)$ (식~\eqref{eq:vapp} 의 분극 부호 $s_I$ 를 양방향으로
일반화한 것 — 방전 $s_I=\sigma_d=+1$, 충전 $\sigma_d=-1$)로 쓴다. \textbf{주의 — $\sigma_d$ 의 작용처는 분극 부호, 분기 중심 선택, 그리고 꼬리의 지지$\cdot$감쇠 방향
(피팅 알고리즘 M5) 셋이며, 평형 곡선의 $s$ 에는 작용하지 않는다}(\S\ref{sec:notation} 규약 그대로). 평형 곡선 $V_{\eq,j}(\xi)$(식~\eqref{eq:hys_Veq})는 상태함수라
충$\cdot$방전 공통이고, 그 안의 $s$ 는 방전 규약 $+1$ 로 고정된다 — 충전 분기라고 $s$ 를 뒤집지 않으며, 곡선은
방향에 불변이다. 충방전 차이는 오로지 \emph{어느
spinodal 어깨}($\xi_{s,j}^-$ vs $\xi_{s,j}^+$, 식~\eqref{eq:hys_center})에서 분리하느냐와 분극 부호 $\sigma_d$ 에서만
온다. 그래서 $U_j^\dis>U_j^\chg$ 가 한 곡선의 극대$>$극소에서 자동으로 따라온다($s$ 충돌 없음).
\begin{equation}
\boxed{\;\frac{\dd Q}{\dd V_\app}\Big|^{d}=C_\bg+\sum_{j=1}^{N_p}Q_j\bigg[\frac{\xi_{\eq,j}(1-\xi_{\eq,j})}{w_j}
+\frac{r_{a,j}}{L_{V,j}}\,e^{-(V_n-V_{a,j})/L_{V,j}}\bigg]_{U_j\to U_j^{d}}\,,}
\label{eq:hys_master}
\end{equation}
\begin{gather*}
U_j^{d}=U_j+\sigma_d\tfrac12\gamma_j\Delta U_j^\hys,\qquad V_n=V_\app-\sigma_d|I|R_n,\\
\Delta U_j^\hys=\tfrac{2}{sF}\big[\Omega_j u_j-2RT\,\mathrm{artanh}\,u_j\big],\qquad u_j=\sqrt{1-\tfrac{2RT}{\Omega_j}}.
\end{gather*}
즉 \emph{같은 항 구조에 중심과 분극 부호만 방향별로} 바뀐다($\xi_{\eq,j}$ 의 중심 $U_j$ 가 $U_j^d$ 로 옮겨갈 뿐).
\textbf{단 적용 범위}: 이 양방향성은 \emph{열역학 중심}($U_j^d$, 분기)과 \emph{분극 부호}($\sigma_d$, IR)에 대해 엄밀하다 —
\emph{꼬리항}($r_{a,j},L_{V,j}$)은 \S\ref{sec:lag} 가 방전(탈리튬화) 동역학에서만 유도했으므로(\S\ref{sec:synth} 가 충전
동역학을 범위 밖으로 명시), 충전 분기의 꼬리는 $\sigma_d$ 로 풀리는 양이 \emph{아니라} 방향별 독립 파라미터
$\{r_{a,j}^\chg,L_{V,j}^\chg\}$ 다 — 리튬화 동역학이 탈리튬화와 대칭이라는 보장이 없기 때문이다. 감쇠의
방향도 거울상이다. 충전 꼬리는 $V_n<V_{a}^{\chg}$ 쪽으로 트레일하므로 지지 조건이
$\sigma_d(V_n-V_{a,j}^{d})\ge0$ 으로 반전되며, 박스의 지수형은 방전($\sigma_d=+1$) 표기다. 본 식의 ``한 식''은 \emph{중심$\cdot$분극}에
한해 양방향이고, 꼬리는 방향별로 따로 닫는다.
\textbf{환원 검산 두 개로} 통합이 일관됨을 확인한다:
\begin{itemize}[leftmargin=1.4em,itemsep=2pt]
\item $\Omega_j\le2RT$(단상, \S\ref{sec:regsol})면 $u_j=\sqrt{1-2RT/\Omega_j}$ 가 허수 — 즉 spinodal 이 없어 분기
  자체가 미정의라 $\Delta U_j^\hys\equiv0$ 으로 둔다(\S\ref{sec:regsol} verifybox). 그러면 $U_j^\dis=U_j^\chg=U_j$ 라
  식~\eqref{eq:hys_master} 가 양방향 모두 방전식~\eqref{eq:master} 로 \emph{되돌아간다}(히스 없음) — 히스는 상분리의 산물임이 식에서 드러난다.
\item $|I|\to0$ 면 분극 항이 사라지고 남는 충방전 gap 이 정확히 $\gamma_j\Delta U_j^\hys$(열역학 잔류,
  식~\eqref{eq:hys_obsgap}). 유한 율에선 거기에 $2|I|R_n$ 이 더해진다.
\end{itemize}

\textbf{분기 영역 peak 의 형상 — 평형형이 아니라 현상학적 폭.} 한 가지 정합을 분명히 해 둔다. 분기는
$\Omega_j>2RT$ 에서만 정의되는데(위 환원 검산), 바로 그 영역에서 평형 isotherm 은 logistic 이 아니라 상분리
plateau 이고, 평형 유효 폭 식~\eqref{eq:weff} 는 음이 되어 무의미하다(\S\ref{sec:regsol}). 따라서
식~\eqref{eq:hys_master} 의 상승부 $\xi_{\eq,j}(1-\xi_{\eq,j})/w_j$ 는 그 영역에서 평형 곡선이 아니라, 관측되는
broadened peak 의 현상학적 형상 ansatz 다(피팅 폭 $w_j$). 그 유한 폭의 출처는 \S\ref{sec:hys_peak} 에서 단계로
보였듯 다입자 분포(\S\ref{sec:dist}$\cdot$\S\ref{sec:regsol} Dreyer)와 동역학(\S\ref{sec:lag} 꼬리)이다.
즉 $w_j$ 의 지위가 두 영역에서 다르다. 단상($\Omega_j\le2RT$)에서는 $w_j=w_j^\eff=(RT/F)(1-\Omega_j/2RT)$
라는 검증 가능한 평형 예측(식~\eqref{eq:weff})이고, 분기 영역($\Omega_j>2RT$)에서는 broadening 이 정하는 독립
피팅 폭이며 평형 $w^\eff$ 의 차용은 금지된다. logistic 함수형은 양 영역 공통의 형상 ansatz 일 뿐, 그 평형
유도는 단상에 한한다.

한 전이의 충방전 peak 쌍은 \emph{면적 $Q_j$ 동일}(같은 전하)$\cdot$\emph{(현상학적) 폭 $w_j$ 동일}(같은 broadening),
\emph{중심만} $\gamma_j\Delta U_j^\hys+2|I|R_n$ 갈린다 — 방전 전용 식~\eqref{eq:master} 의 모든 파라미터에
\emph{전이당 $\{\Omega_j,\gamma_j\}$ 둘만 더}하면 충방전의 중심과 분극이 닫힌다. 분기 영역의 $w_j$ 는
위에서 말한 대로 독립 피팅이고, 충전 꼬리는 적용 범위 주대로 S3$\cdot$S4 를 충전 방향 데이터로 반복해
$\{r^{\chg},L^{\chg}\}$ 로 따로 닫는다. 그것까지 마치면 완전한 양방향이다.
\emph{부분 cycle}(분기 폭의 일부 $\eta$ 만 진행)을 다룰 때는 식~\eqref{eq:hys_partial} 대로 분기 폭이 $h_j(\eta)$ 배로
줄어, 식~\eqref{eq:hys_master} 의 중심이 $U_j^{d}=U_j+\sigma_d\tfrac12\,h_j(\eta)\,\gamma_j\Delta U_j^\hys$ 로 일반화된다
(완전 cycle 은 $h_j=1$ 이라 위와 동일). \emph{식별 분리}: 절편이 $h_j(\eta)\gamma_j$ 곱이라 둘이 공선형이므로 — $\gamma_j$ 는
\emph{완전 cycle}($\eta=1,\,h_j=1$) gap 에서 먼저 닫고, $h_j(\eta)$ 는 \emph{부분 cycle} 데이터에서만 분리한다(완전
cycle 데이터가 $\gamma_j$ 를 고정하는 전제). 부분 cycle 데이터가 있을 때만 $h_j$ 보간을 더한다.

\subsection{피팅 파라미터 — 전이당 무엇을, 어느 데이터로}
\begin{center}
\renewcommand{\arraystretch}{1.3}
\begin{tabular}{p{0.10\textwidth} p{0.30\textwidth} p{0.34\textwidth} p{0.14\textwidth}}
\toprule
파라미터 & 의미 & 닫는 데이터 & 식 \\
\midrule
$U_j$ & 평형 중심 전위 & 저율 OCV peak \emph{위치} & \eqref{eq:eqpeak} \\
$w_j$ & 평형 폭 $=(RT/F)(1-\Omega_j/2RT)$(단상; 분기 영역은 현상학적 폭) & 저율 OCV peak \emph{너비}(이게 닫는 1차 미지수는 $w_j$; 그 다온도 형태는 gap-경로 $\Omega_j$ 의 \emph{검증}—독립 자유도 아님) & \eqref{eq:weff} \\
$Q_j$ & 전이 용량 & 저율 OCV peak \emph{면적} & \eqref{eq:eqpeak} \\
$R_n$ & 분극 저항 & 전류 차단 IR-drop(GITT) & \eqref{eq:vapp} \\
$\chi$ & 전달 계수 & 다율 $\partial\ln L_{q,j}/\partial V$ & \eqref{eq:LqV} \\
$\Delta H_{a,j}$ & 활성화 엔탈피 & 다온도 Arrhenius 기울기 & \eqref{eq:arrhenius} \\
$C_\bg$ & 배경 미분용량 & 저율 OCV 기저선 & \eqref{eq:dQdV} \\
($\rho_G$ 폭) & 배리어 분포 폭(선택) & 꼬리 stretched 휨 & \eqref{eq:superpose} \\
$\Omega_j$ & 정규용액 상호작용(상분리) & \emph{1차}: 충방전 gap $|I|\!\to\!0$ 절편의 \emph{온도 의존}(폭 $w_j(T)$ 는 같은 $g_j''$ 라 \emph{검증}용, 자유도 중복 셈 금지) & \eqref{eq:hys_dU} \\
$\gamma_j$ & 분기 축소(핵생성$\cdot$다입자, $\le1$) & 충방전 gap $|I|\!\to\!0$ 절편 \emph{크기}($\Omega_j$ 고정 후 분리) & \eqref{eq:hys_center} \\
\bottomrule
\end{tabular}
\end{center}
즉 \emph{전이당 핵심 $\{U_j,w_j,Q_j,\Delta H_{a,j}\}$ 와 전이 공통 $\{\chi,R_n,C_\bg\}$}; 꼬리가 휠 때만 $\rho_G$ 폭을,
\emph{충방전을 다룰 때만} 전이당 $\{\Omega_j,\gamma_j\}$ 둘을 더한다(방전만이면 불필요).

\subsection{데이터$\to$파라미터(비순환)$\to$예측}
식별은 \textbf{S0$\to$S5} 의 순서를 따른다. \S\ref{sec:synth} 의 S1--S4 에 선행 진단 S0 와 충방전 S5 를 더한
것이며, 각 단계가 앞 단계의 기지값을 써서 미지수를 하나씩만 남겨 공선형을 끊는다.
\begin{enumerate}[label=\textbf{S\arabic*.},leftmargin=2.6em,itemsep=2pt,start=0]
\item \textbf{(선행) 수송 형상 진단.} GITT 완화가 $\sqrt t$ Cottrell 형(반무한 확산 율속)인지 지수형(계면
  율속)인지 판정해(\S\ref{sec:falsify}(i)) \emph{활성화 지배}($k_{j,\mathrm{act}}\ll k_{j,\lim}$)를 확정하고
  $k_{j,\lim}$ 을 추정한다. S3 의 $\chi$ 분리가 식~\eqref{eq:LqV} 감쇠인자를 통해 이 확정에 의존하므로,
  S0 는 S3 앞에 와야 순환이 없다. \S\ref{sec:falsify}(i) 는 그 반복$\cdot$정량 점검이다.
\item \textbf{저율 OCV}($|I|\!\to\!0$). 평형 peak 기하에서 $\{U_j,w_j,Q_j,C_\bg\}$ 를 닫는다.
\item \textbf{전류 차단.} IR-drop 도약에서 $R_n$ 을 닫는다.
\item \textbf{다율.} 동일 율 안에서 여러 tail $V_\drive$ 점을 회귀해 $\chi$ 를 닫고, 두 율 사이에서는
  $L_q\propto|I|$ 로 교차검증한다.
\item \textbf{다온도 Arrhenius.} 구동항 보정 후 기울기에서 $\Delta H_{a,j}$ 를 닫는다. $\Delta S_{a,j}$ 는
  절편 조합으로만 들어가며 단독으로는 식별되지 않는다.
\item \textbf{충방전 gap}(충방전을 다룰 때만). 여러 율의 gap $\Delta V_{\mathrm{obs},j}(|I|)$ 를
  $|I|\!\to\!0$ 으로 외삽한다. 절편의 크기가 $\gamma_j$ 를, 절편의 온도 의존이 $\Omega_j$ 를 닫고, 기울기는
  $R_n$ 의 교차확인이다(식~\eqref{eq:hys_obsgap}). 단 깊은 2상측($T_{c,j}\gg$측정창)에서는 절편의 $T$-의존이
  포화해 $\Omega_j$ 는 하한만 잡히고 $\gamma_j$ 도 곱 $\gamma_j\Omega_j$ 로만 남는다 —
  \S\ref{sec:falsify}(ii) 의 regime 분기다.
\end{enumerate}
여러 전이가 겹치면 식~\eqref{eq:total} 의 합산으로 저율-anchored \emph{forward} 예측한다. 역산은 금지된다
(\S\ref{sec:overlap}).
\begin{keybox}
\textbf{이 식 하나로 끝난다.} 저율 OCV $+$ 전류 차단 $+$ 다율$\cdot$다온도(예 OCV$\cdot$0.1C$\cdot$0.2C $\times$
15/23/45$^\circ$C, 충방전 양방향)로 위 파라미터를 닫으면, \emph{같은 통합식~\eqref{eq:hys_master} 에 임의의}
$(T,|I|,V_\drive,$ 방향$)$ 를 넣어 그 조건의 \emph{방전$\cdot$충전} $\dd Q/\dd V$ 와 그 히스테리시스 gap 을
\emph{예측}한다($\Omega\le2RT$ 면 방전식~\eqref{eq:master} 로 자동 환원; 충전 \emph{꼬리 형상}까지 닫으려면 충전
데이터로 $\{r^{\chg},L^{\chg}\}$ 별도 — 중심$\cdot$분극$\cdot$gap 은 위 파라미터로 충분). 피팅은 식 하나, 파라미터는 전이당 소수,
데이터원은 분리돼 단계적. 더 정확히는 ``\emph{파라미터가 주어지면 예측이 한 식}''이며, 그 \emph{입력}(수송 형상
판정$\cdot$$R_\mathrm{ct}$ 분리$\cdot$필요 시 $\rho_G$$\cdot$절대 prefactor)은 표준 진단으로 따로 공급된다. \textbf{그리고
``하나로 끝난다''는 다음 가정의 울타리 안에서다.}
\begin{itemize}[leftmargin=1.6em,itemsep=1pt,topsep=2pt]
\item \emph{동역학} — ① 활성화 지배($k_{j,\lim}$ 무한; 아니면 \S\ref{sec:falsify}(i) 수송 진단 선행) ② 좁은
  배리어 분포($\rho_G\to\delta$; 휘면 \S\ref{sec:dist} 적분 켜짐) ③ prefactor $T$-안정(지수 앞 인자의 $T$-멱이
  Arrhenius 기울기를 오염하지 않게 절편에 흡수) ④ $q_a$ 컷 고정(꼬리 시작점 정의를 전 데이터셋에 공통 적용).
\item \emph{히스} — ⑤ site energy 의 $\xi=\tfrac12$ 대칭(분기 중심을 $U_j\pm\tfrac12\gamma_j\Delta U_j^\hys$ 로
  대칭 분할하는 근거) ⑥ $\gamma_j$ 전류$\cdot$온도(측정창) 무관(저$\sim$중율, 식~\eqref{eq:hys_obsgap}; S5
  온도-분리의 전제) ⑦ 분극 선형$\cdot$$R_n$ lumped(전이$\cdot$SOC 무관, $\eta_\mathrm{ct}\propto\ln i$ 부차).
\item \emph{겹침} — ⑧ 전이 비중첩 선형 가산(아니면 \S\ref{sec:overlap} forward-only).
\item \emph{적용역} — ⑨ $C_\bg$ 를 S1 에서 함께 고정 ⑩ \emph{반쪽전지} 환산 선행(full-cell 은 양극 분리 후).
\end{itemize}
또 $\Delta S_{a,j}$ 는 prefactor 와 결합으로만 들어가 \emph{단독 비식별}(온도 외삽엔 절편
흡수로 무해). 이 단서가 깨지면 해당 추가 절차가 켜진다 — \emph{식은 하나지만 적용역은 이 (10중) 울타리 안}.
\end{keybox}

\subsection{데이터에서 예측까지 — 한 walkthrough}
추상적 절차를 한 흐름으로 굳힌다. 한 셀의 흑연 음극을 예로(전제 \textbf{S0}: GITT 완화 형상이 지수형임을 먼저
확인해 \emph{활성화 지배}를 확정한다 — 아래 3$\cdot$4의 꼬리 회귀가 전부 이 위에 선다, \S\ref{sec:falsify}(i)):
\begin{enumerate}[label=\textbf{\arabic*.},leftmargin=2.0em,itemsep=1pt]
\item \textbf{준평형}($0.05$C 또는 GITT) 방전 $\dd Q/\dd V$ 에서 분리된 $N_p$ 개 peak 의 위치$\cdot$너비$\cdot$면적을
  읽어 $\{U_j,w_j,Q_j\}$ 와 기저선 $C_\bg$ 를 고정한다(S1). 단 \emph{히스가 보이는 전이}는 방전 peak 위치가
  $U_j^{\dis}=U_j+\tfrac12\gamma_j\Delta U_j^\hys$ 라 그대로 $U_j$ 가 아니다 — 충$\cdot$방전 저율 peak 중심의
  \emph{중점}을 $U_j$ 로 잡는다(\S\ref{sec:hys_branch} 의 대칭 분할; 중심\emph{차}는 5의 gap 으로 간다). 또 단계
  4의 구동항 보정이 각 온도의 $U_j(T)$ 를 쓰므로 이 읽기는 \emph{각 온도}(15/23/45$^\circ$C OCV)에서 반복한다(또는
  식~\eqref{eq:dUdT} 로 외삽). 분리 peak 수 $N_p$ 와 전이 경계 SOC 는 DVA 로 교차점검한다(\S\ref{sec:dva}).
\item 같은 SOC 들에서 전류 차단 IR-drop 으로 $R_n$ 을 얻는다(S2).
\item 각 율 \emph{안에서}(율별로) 동일 SOC$\cdot$tail window 의 $\ln L_{q,j}$ 를 $V_\drive$ 로 회귀해 $\chi$
  (측정 직독량은 꼬리의 전압축 로그기울기 $1/L_{V,j}$ — 측정 $V(q)$ 의 국소 기울기로
  $L_{q,j}=L_{V,j}/|\dd V_n/\dd q|_{q_a}$ 환산해 쓴다, 상세 절차는 아래 피팅 알고리즘 소절) —
  $0.1$C$\cdot$$0.2$C 두 율 \emph{간}은 $L_q\propto|I|$ 의 교차검증(합산 회귀 금지: 율 간 $\ln|I|$ 절편차가 기울기를
  오염)(S3).
\item $15/23/45^\circ$C 의 $\ln L_{q,j}$ 를 구동항 보정해 $1/T$ 로 회귀, 기울기에서 $\Delta H_{a,j}$(S4).
\item \textbf{충방전을 다룰 때만}: 같은 전이의 충$\cdot$방전 peak 중심차 $\Delta V_{\mathrm{obs},j}$ 를 여러 율에서 재어
  $|I|\to0$ 외삽 — 절편 크기로 $\gamma_j$, 절편의 $15/23/45^\circ$C 온도 의존으로 $\Omega_j$(식~\eqref{eq:hys_dU})
  를 닫는다(S5). 방전만이면 생략.
\end{enumerate}
이제 식~\eqref{eq:hys_master}(방전만이면 식~\eqref{eq:master})에 이 파라미터와 \emph{원하는} $(T,|I|,$ 방향$)$ 을
넣으면 그 조건의 곡선과 충방전 gap 이 나온다. \textbf{검증 시험}: 피팅에 \emph{쓰지 않은} 조건(예 $0.5$C$\cdot$
$5^\circ$C, 또는 다른 율의 충방전 gap)을 예측하고 그 측정과 대조한다 — 맞으면 모델이 서고, 빗나가면
\S\ref{sec:falsify} 의 경쟁 원인(확산$\cdot$전하전달)이 살아 있다는 신호다. \emph{피팅이 아니라 외삽 예측이 진짜 시험}이다.

\subsection{피팅 알고리즘 — 측정 파일에서 파라미터까지}
위 walkthrough 가 절차의 뼈대였다면, 이 소절은 그것을 실제 데이터 처리 알고리즘으로 닫는다. 목표는 분명하다 —
\emph{이 문건만 보고 방전 $\dd Q/\dd V$ 피팅 코드를 작성할 수 있어야 한다.} 아래 7 단계가 그 사양이다.

\textbf{(1) 입력 데이터 명세.} 정전류 구간별 시계열 $(t_i,\,V_{\app,i},\,I_i,\,T_i)$ 와 방향 라벨
$d\in\{\dis,\chg\}$ 를 받는다. 필요한 데이터셋은 \S\ref{sec:master} keybox 의 조합 그대로다 — 준평형
(0.05C 또는 GITT)$\cdot$0.1C$\cdot$0.2C 의 각 율을 $15/23/45^\circ$C 에서, 충방전 양방향으로, 그리고 전류 차단
transient(S2 용)를 따로 둔다. 반쪽전지 척도($V_n$ 기준)가 전제다(가정 ⑩).

\textbf{(2) 전처리.} 다음 순서로 측정 파일을 피팅 좌표로 옮긴다.
\begin{enumerate}[label=(\roman*),leftmargin=2.0em,itemsep=1pt]
\item \emph{용량 좌표}: $Q_i=\sum_{i'\le i}|I_{i'}|\,\Delta t_{i'}$ 를 누적하고 $q_i=Q_i/Q_\cell$ 로 정규화한다.
\item \emph{미분}: $\dd Q/\dd V$ 는 잡음을 증폭하므로 생(raw) 차분을 쓰지 않는다. $Q(V)$ 를 균일 $V$ 격자로
  재표집한 뒤 국소 다항(Savitzky--Golay) 평활 미분을 쓰거나, 동등하게 중심차분 후 같은 창으로 평활한다.
  \emph{평활 창은 가장 좁은 peak 의 FWHM 의 $1/5$ 이하}로 잡는다 — 과평활은 폭 $w_j$ 에 양(넓어지는 쪽)의
  계통 편향을 넣는다.
\item \emph{전위 환산}: $V_n=V_\app-\sigma_d|I|R_n$(식~\eqref{eq:vapp}). $R_n$ 은 S2 에서야 확정되므로 첫
  패스에서는 $V_\app$ 그대로 peak 검출만 하고, $R_n$ 확정 후 재환산한다. 반복은 1회로 충분하다 — $R_n$ 보정은
  위치의 평행이동이라 형상을 바꾸지 않기 때문이다.
\item \emph{온도}: 등온 구간 평균 $T$ 를 쓴다. 발열에 의한 내부 $T$ 변화의 보정은 발열 장(Chapter 3)의 소관이다.
\end{enumerate}

\textbf{(3) peak 검출과 초기값.} 저율 OCV 의 $\dd Q/\dd V$ 에서 prominence(주변 골 대비 돌출 높이) 기준으로
peak 를 검출해 $N_p$ 를 정하고, DVA 의 valley 수$\cdot$경계 위치로 교차점검한다(\S\ref{sec:dva}). 전이별
초기값은 peak 기하에서 직접 읽는다:
\begin{gather*}
U_j^{(0)}=V_{\peak,j},\qquad w_j^{(0)}=\frac{\mathrm{FWHM}_j}{3.53},\\
Q_j^{(0)}=\int(\text{배경 차감 peak}_j)\,\dd V,\qquad C_\bg^{(0)}=\text{peak 간 골의 평활 기저선}.
\end{gather*}
히스가 보이는 전이의 $U_j^{(0)}$ 는 충$\cdot$방전 저율 peak 중심의 중점으로 잡는다(walkthrough 1). 읽은
FWHM 이 이상 상한 $3.53\,RT/F\approx91$ mV($25^\circ$C)를 넘으면 꼬리 오염 신호이므로 초기값 단계에서 경고를
띄운다(\S\ref{sec:synth} worked example). 이 단계의 운용 규칙 셋을 못 박아 둔다. 첫째, \emph{$N_p$ 의 확정과
동결} — $N_p$ 는 기준 온도(예 $23^\circ$C)의 저율 OCV 방전에서 prominence 가 최대 peak 순높이의 $5\%$ 이상인
봉우리로 확정해 전 데이터셋에 동결한다. 다른 온도$\cdot$율$\cdot$방향에서는 새로 검출하지 않고 동결된 $U_j$ 의
예상 위치 근방에서 대응 peak 를 찾으며, 미검출이면 융합 플래그를 단다(\S\ref{sec:overlap}). 둘째,
\emph{충$\cdot$방전 peak 의 짝짓기} — 전위 순서대로 1:1 로 맺고, 개수가 다르면 면적이 가장 비슷한 쌍부터
맺는다. 중심차가 측정 분해능과 $2|I|R_n$ 의 합을 유의하게 넘는 전이만 ``히스가 보이는 전이''로 분류해 중점
절차를 켜며, $w_j$ 와 $Q_j$ 는 양방향 평균으로 읽는다. 면적은 물질량 보존상 양방향이 같아야 하므로
(\S\ref{sec:hys_peak}), 그 차이는 오히려 데이터 품질의 진단 플래그다. 셋째, \emph{배경의 함수형} —
$C_\bg(V_n)$ 은 peak 간 골의 anchor 점들을 지나는 저차 spline(또는 구간 상수)으로 두고 S1 에서 형상을
동결하며, 온도별 OCV 에서 각각 읽는다.

\textbf{(3$'$) 꼬리 시작점 $q_a$ 와 tail window 의 조작 규정.} \S\ref{sec:lag} 의 정의 — 평형 항의 포화 —
를 코드로 옮기면 다음과 같다. 저율 OCV 로 고정한 $\xi_{\eq,j}(V_n)$ 과 피팅 대상 율의 측정 $V_n(q)$ 로
연쇄율 $\dd\xi_{\eq,j}/\dd q=(\dd\xi_{\eq,j}/\dd V_n)(\dd V_n/\dd q)$ 를 평가하고, 그 값이 전이
정점값의 $5\%$ 로 떨어지는 정점 이후 첫 교차점을 $q_{a,j}$ 로 잡는다. 컷 분율 $5\%$ 는 전 데이터셋에 공통
고정한다(가정 ④). 교차점이 $\dd V_n/\dd q$ 가 매우 작은 plateau 구간에 떨어지면 그 구간을 제외하고 첫 단상
지점으로 민다(\S\ref{sec:lag} 의 plateau 제외 단서와 동일). tail window 는 $[q_{a,j},$ 다음 peak 의 상승
개시$)$ 로 잡고, S3 회귀를 위해 그 창을 몇 개의 sub-window 로 나눠 각 창의 로그기울기에서 국소 $L_V$ 를 얻은
뒤, 창 중심의 $V_n$ 을 $V_\drive$ 로 짝지어 $(V_\drive,\ln L_q)$ 점들을 만든다.

\textbf{(4) 모델 평가(의사코드).} 파라미터 벡터
$\bm\theta=\{U_j,w_j,Q_j,\Delta H_{a,j}\}_{j=1}^{N_p}\cup\{C_\bg,R_n,\chi,\text{Arrhenius 절편}\}$
(충방전이면 전이당 $\{\Omega_j,\gamma_j\}$ 추가)와 조건 $(T,|I|,d)$ 가 주어졌을 때, 모델 곡선
$\widehat{\dd Q/\dd V}(V_\app;\bm\theta)$ 는 다음 순서로 계산한다.
\begin{enumerate}[label=\textbf{M\arabic*.},leftmargin=2.4em,itemsep=1pt]
\item 충방전이면 $u_j=\sqrt{1-2RT/\Omega_j}$, $\Delta U_j^\hys$(식~\eqref{eq:hys_dU}),
  $U_j^d=U_j+\sigma_d\tfrac12\gamma_j\Delta U_j^\hys$(식~\eqref{eq:hys_center})를 계산한다.
  $\Omega_j\le2RT$ 면 $\Delta U_j^\hys=0$ 으로 두어 방전식으로 자동 환원된다. 방전 전용이면 $U_j^d=U_j$.
\item $V_n$ 격자에서 $\xi_{\eq,j}=[1+e^{-s(V_n-U_j^d)/w_j}]^{-1}$(식~\eqref{eq:logistic})를 평가한다.
\item 구동력 $\mathcal A_j=sF(V_n-U_j^d)$ 로 꼬리 길이
  $L_{q,j}=(|I|h/Q_\cell k_BT)\,e^{(\Delta H_{a,j}-T\Delta S_{a,j}-\chi\mathcal A_j)/RT}$ 를 평가한다.
  $T\Delta S_a$ 는 단독이 아니라 Arrhenius 절편 조합으로 들어간다(\S\ref{sec:synth} S4). 평가는 $q_a$ 한 점의
  $\mathcal A_j(V_{a,j})$ 로 하여 전이당 \emph{상수}로 동결한다 — 식~\eqref{eq:master} 의 지수 꼬리는 상수
  $L$ 의 닫힌 해이므로 격자 점별 $L(V_n)$ 대입은 금지된다(slowly-varying 검산, \S\ref{sec:lag}). $q_a$ 컷이
  $\mathcal A_j<3RT$ 에 걸리면 보편형 보정 인자 $(1+e^{-\mathcal A_j/RT})$ 를 $k_j$ 에 곱한다
  (\S\ref{sec:barrier}).
\item $L_{V,j}=|\dd V_n/\dd q|_{q_a}L_{q,j}$ 로 환산하고, 꼬리 시작점 $V_{a,j}=V_n(q_{a,j})$ 를 $q_a$ 컷
  (가정 ④)에서 읽는다. 기울기 $|\dd V_n/\dd q|_{q_a}$ 는 피팅 시에는 피팅 대상 율의 측정 $V(q)$ 평활
  기울기에서 읽고, 측정이 없는 조건을 예측할 때는 S1 파라미터로 닫힌 전하 보존식~\eqref{eq:charge} 의 평형 해
  $V_{n,\OCV}(q)$ 의 기울기로 대체한다(또는 최근접 측정 율의 $V(q)$ 를 재사용한다). 잔류 지연의 엄밀값은 \S\ref{sec:lag} 일반해의 기억 적분을 $q_a$ 에서 평가한
  $r_{a,j}=\int^{q_a}e^{-(q_a-q')/L_{q,j}}(\dd\xi_{\eq,j}/\dd q')\,\dd q'$ 이다. 실무 기본값은 regime
  으로 가른다. 지연이 작은 $L_{V,j}\ll w_j$ 에서는 극한 (i)의 접속값
  $r_{a,j}\simeq L_{q,j}\,(\dd\xi_{\eq,j}/\dd q)\big|_{q_a}$ 가 정확하다. 그러나 꼬리가 관측되는 꼬리-우세
  regime($L_{V,j}\gtrsim w_j$)에서는 이 접속값이 기억 적분을 크게 과소평가하므로 쓰지 않는다 — 그 영역에서는
  지연이 포화해 $r_{a,j}\simeq\xi_{\eq,j}(q_a)-\xi_j(q_a)$ 가 $O(1)$ 이며, $r_{a,j}\in(0,\,\xi_{\eq,j}(q_a)]$
  범위의 자유 파라미터로 풀어 준다(그 경우 $\bm\theta$ 에 추가). 분포 폭이나 컷 민감도가 클 때도 같은 자유화를
  쓴다.
\item 항을 합산한다: $\widehat y=C_\bg+\sum_jQ_j\big[\xi_{\eq,j}(1-\xi_{\eq,j})/w_j
  +(r_{a,j}^d/L_{V,j}^d)\,e^{-\sigma_d(V_n-V_{a,j}^d)/L_{V,j}^d}\,\mathbf 1\{\sigma_d(V_n-V_{a,j}^d)\ge0\}\big]$
  (식~\eqref{eq:hys_master}). 지수 안의 $\sigma_d$ 가 충전 분기의 거울상 감쇠를 처리하고 지시함수
  $\mathbf 1$ 이 단방향 지지를 처리한다 — 방전($\sigma_d=+1$)에서 본문 표기로 환원되며, 충전 꼬리 파라미터
  $\{r_{a,j}^\chg,L_{V,j}^\chg\}$ 는 방향별로 따로 닫는다(\S\ref{sec:master} 통합형의 적용 범위 주).
\item 측정축으로 되돌린다: $V_\app=V_n+\sigma_d|I|R_n$.
\end{enumerate}

\textbf{(5) 목적함수.} 전역 동시 최소화는 공선형 때문에 금지된다(\S\ref{sec:synth}). 각 단계 안에서의 잔차는
가중 최소제곱
\begin{equation*}
\chi^2(\bm\theta_{\text{free}})=\sum_i\frac{\big[y_i-\widehat y(V_i;\bm\theta)\big]^2}{\sigma_i^2}
\end{equation*}
로 두고, 해당 단계의 자유 파라미터 $\bm\theta_{\text{free}}$ 만 움직인다. 가중 $\sigma_i$ 는 미분이 증폭한
잡음을 반영해 평활 잔차에서 추정한다. $\dd V/\dd Q$ 좌표로 바꿔 피팅하는 것은 작은 $\dd Q/\dd V$ 영역에서
잡음을 $1/y^2$ 로 증폭하므로 권장하지 않는다(\S\ref{sec:dva}).

\textbf{(6) 단계별 회귀 — 입출력 표.} S0--S5 각각이 무엇을 받아 무엇을 닫는지를 한 표로 모은다. 각 단계에서
이전 단계의 출력은 \emph{동결}된다.
\begin{center}
\renewcommand{\arraystretch}{1.35}
\begin{tabular}{p{0.07\textwidth} p{0.24\textwidth} p{0.36\textwidth} p{0.22\textwidth}}
\toprule
단계 & 입력 데이터 & 회귀(자유 파라미터) & 출력(동결) \\
\midrule
S0 & GITT 완화 transient & 회귀 아님 — $\sqrt t$ 형 vs 지수형 판정 & 활성화 지배 확인 \\
S1 & 저율 OCV $\dd Q/\dd V$(각 $T$) & peak 기하 직독 $+$ 평형식 국소 회귀 & $\{U_j,w_j,Q_j\},C_\bg$ \\
S2 & 전류 차단 도약 & $\Delta V=s_I|I|R_n$ 직독 & $R_n$ \\
S3 & 등온 다전위 꼬리(율별) & $\ln L_{q,j}$ vs $V_\drive$ 기울기 $-\chi sF/RT$ & $\chi$ \\
S4 & 다온도 꼬리 & $y(T)$ vs $1/T$ 기울기 $-\Delta H_a/R$(식~\eqref{eq:arrhenius}) & $\Delta H_{a,j}$, 절편 조합 \\
S5 & 다율 충방전 gap(각 $T$) & 절편$\cdot$기울기 분해(식~\eqref{eq:hys_obsgap}) $+$ 절편의 $T$-형 회귀 & $\gamma_j,\Omega_j$($R_n$ 교차확인) \\
\bottomrule
\end{tabular}
\end{center}

\textbf{(6$'$) 회귀의 보충 규정.} S3 와 S4 는 (3$'$)의 sub-window 점들로 다음과 같이 돈다. S3 는 고정 $T$
에서 $(V_\drive,\ln L_q)$ 점들을 직선 회귀해 기울기 $-\chi sF/RT$ 에서 $\chi$ 를 얻는다. S4 는 전 온도
공통의 기준점 $q^*$ — 예컨대 각 전이의 $q_a$ — 에서 $L_q(T)$ 와
$\mathcal A_j(T)=sF(V_n(q^*;T)-U_j(T))$ 를 평가해 식~\eqref{eq:arrhenius} 의 $y(T)$ 를 만들고 $1/T$ 로
회귀한다. S5 의 절편은 온도 3점에 대해 절편$(T)=\gamma_j\,\Delta U_j^\hys(T;\Omega_j)$ 의 2-파라미터
비선형 최소제곱으로 풀며, 초기값은 $\Omega_j^{(0)}=4RT_{\mathrm{room}}$ 과
$\gamma_j^{(0)}=$절편$/\Delta U^\hys(\Omega_j^{(0)})$ 로 둔다. 단 절편의 상대 온도 기울기
$|\dd\ln(\text{절편})/\dd T|$ 가 측정창에서 그 측정 불확도보다 작으면 깊은 2상측 퇴화로 판정한다
(\S\ref{sec:falsify}(ii)). 그 경우 $(\gamma_j,\Omega_j)$ 의 분리를 포기하고, lumped gap
$\Delta_j\equiv\gamma_j\Delta U_j^\hys$ 한 개를 온도 상수로 두어 M1 에 직접 쓴다. 충전 꼬리를 닫을 때는
S3$\cdot$S4 를 충전 방향 데이터로 반복하되, 역방향 속도가 율속이 되므로 계수가 $\chi$ 에서 $1-\chi^{\chg}$
로 바뀐다. 가로축을 전위 척도 $|\mathcal A_j|/(sF)=|V_n-U_j^{\chg}|$ 로 두면 기울기가
$-(1-\chi^{\chg})sF/RT$ 인 같은 직선 회귀가 된다 — 구동이 깊어질수록($|\mathcal A_j|\uparrow$) 충전 꼬리도
짧아진다는 부호다.

\textbf{(7) 수렴과 검증.} 단계가 끝날 때마다 그 출력은 이후 단계에서 동결한다 — 뒤 단계에서 자유화하면
공선형이 되살아난다(\S\ref{sec:overlap} S$_4'$ 의 같은 규율). 예측에 필요한 온도 보간 규칙도 여기서 닫는다. $U_j(T)$ 는 세 온도의 S1 값을 선형 회귀해 외삽하며, 그 기울기가
곧 $\Delta S_j/(sF)$ 다(식~\eqref{eq:dUdT}). 분기 영역의 $w_j(T)$ 와 $C_\bg(V,T)$ 는 평형 예측이 없으므로
측정 온도 사이는 선형 보간하고, 측정창 바깥은 최근접 온도 값을 유지하며 불확도 플래그를 단다. 전 단계가 끝나면
피팅에 쓰지 않은 hold-out 조건(예: $0.5$C, $5^\circ$C, 또는 다른 율의 충방전 gap)을 예측해 측정과 대조한다 — walkthrough 의 검증 시험
그대로다. 잔차의 패턴은 진단으로 읽는다. 꼬리가 단일지수에서 벗어나 휘면 $\rho_G$ 분포 적분을 켜고
(\S\ref{sec:dist}), gap--$|I|$ 곡선이 휘면 \S\ref{sec:hys_pol} 의 곡률 진단을, Arrhenius 가 비선형이면
\S\ref{sec:falsify} 의 경쟁 원인 배제를 적용한다. 마지막으로 코드 자체의 검증에는 round-trip 시험을
권장한다. 알려진 파라미터 $\bm\theta^\ast$ 로 M-체인을 돌려 합성 $\dd Q/\dd V$ 를 만들고, 측정 수준의
잡음을 더한 뒤 전 파이프라인(전처리부터 S5 까지)을 돌려 $\bm\theta^\ast$ 가 불확도 안에서 회복되는지 본다.
단계별 동결 순서가 공선형을 실제로 끊는지, 평활 창이 $w_j$ 에 편향을 넣지 않는지가 이 시험에서 정량으로
드러난다. 이로써 측정 파일에서 파라미터, 그리고 임의 조건의 예측 곡선까지
가는 길이 전부 닫혔다. 마지막 절은 이 모델이 \emph{어떻게 틀릴 수 있는지}를 묻는다.

% ====================================================================
\section{검증과 반증(falsification)}\label{sec:falsify}
% ====================================================================
지금까지 peak 의 모양과 인자 의존을 도출하고, 겹침을 예측하고, 측정 파일에서 파라미터까지 가는 알고리즘을
닫았다. 좋은 이론은 반증 가능해야 하므로, 마지막으로 묻는다. ``꼬리는 전위가 낮춘 유효 배리어의 동역학
산물''이라는 본 장의 주장은 어떻게 틀릴 수 있는가. 그리고 같은 꼬리를 내는 경쟁 원인과는 어떻게 구별하는가.

\textbf{반증 가능한 예측.} 식~\eqref{eq:arrhenius} 에서, 꼬리 길이의 Arrhenius 기울기를 구동항으로 보정한 $y(T)$
의 기울기는 \emph{순수} $-\Delta H_a/R$ 이어야 하며 \textbf{구동 전위에 무관}해야 한다. 성립 조건을 먼저 둔다 — 해당 온도창서 $\mathcal A_j$ 가 거의 상수이고 \emph{활성화 지배}($k_{j,\mathrm{act}}\ll
k_{j,\lim}$; 아래 (i) 수송 진단으로 \emph{검증}, 미수행 시 가정)인 \emph{국소} 근사다. 그 위에서 보정 전 유효
기울기는 $-(\Delta H_a-\chi\mathcal A_j)/R$ 라 \emph{전위에 의존}한다. 일반적으로 $\mathcal A_j(T)$ 비선형성$\cdot$
유한 공율속의 직렬 율속 항이 기울기를 추가로 왜곡하지만, 핵심은 \emph{전위 의존 자체}다 — \textbf{이 전위
의존이 ``전위가 배리어를 낮춘다''의 반증 가능한 지문}이다. \emph{주의}: 반증되는 것은 \textbf{구동항 $\chi\mathcal A_j$ 를
빼낸 \emph{잔여} $y(T)$ 의 전위 의존}이지, 모형의 본질 예측인 $L_{q,j}$ 자체의 전위 의존(식~\eqref{eq:LqV}, 이건 $0$ 이
아니라 \emph{모형이 맞다는 신호})이 아니다. 또 이 보정이 깨끗이 떨어지는 것은 \emph{활성화 지배}($k_{j,\mathrm{act}}\ll
k_{j,\lim}$)에서뿐 — 유한 공율속이면 감쇠 인자가 곱해져 보정이 불완전하므로, 그 잔여를 모형 기각으로 읽기 전에 (i)
수송 진단으로 활성화 지배를 \emph{먼저} 확정해야 한다(아니면 거짓 기각). 그 선결 뒤, $\chi\mathcal A_j$ 보정 후에도
\emph{잔여} 전위 의존이 남으면 본 설명은 기각된다.

\textbf{경쟁 꼬리원의 배제.} 꼬리 길이 $L_{q,j}\propto|I|$ 와 ``전위$\uparrow\Rightarrow$꼬리$\downarrow$''
라는 거시 거동은 단순 확산과 전하전달 과전위도 공유한다. 한편 상수 접촉저항은 지수 꼬리가 아니라 강체 IR
이동이므로, S2 전류 차단의 $R_n$ 으로 이미 분리되어 있다(식~\eqref{eq:vapp}). 따라서 거시 거동만으로
barrier-lowering 에 귀속하면 오귀속이다. 다음 두 분리가 모두 필요하다:
\begin{enumerate}[label=(\roman*),leftmargin=2.0em,itemsep=2pt]
\item \textbf{확산 배제.} 고체상 확산이 율속이고 확산계수가 점유율 의존이면 transport-limited 유효 엔탈피도 전위
  의존이다. GITT 완화 transient 형상($\sqrt t$ Cottrell=반무한 확산 vs 지수=계면 율속)으로 이를 \emph{독립 판정}해
  배제한 뒤에만 위 지문이 유효하다. 이 형상 판정은 표준 진단이며 본 장 이론의 입력이다 \cite{weppner1977}.
\item \textbf{전하전달 분리.} charge-transfer 과전위 $\eta_\mathrm{ct}$ 도 전위 지수의존이라 단일 전극서 $\chi$ 와
  공선형이다. GITT/전류 차단(또는 EIS)으로 $R_\mathrm{ct}$ 를 독립 분리하는 표준 전기화학 진단 뒤에만 $\chi$ 의
  ``배리어 낮춤'' 귀속이 성립한다 \cite{weppner1977,bardfaulkner}.
\end{enumerate}
세 후보는 \emph{꼬리의 거시 거동}($\propto|I|$, 전위$\uparrow\Rightarrow$꼬리$\downarrow$)을 \emph{공유}하므로
거시 관측만으론 못 가른다 — 갈라 주는 것은 서로 다른 \emph{독립 진단}이다:
\begin{center}
\renewcommand{\arraystretch}{1.35}
\begin{tabular}{p{0.26\textwidth} p{0.30\textwidth} p{0.32\textwidth}}
\toprule
원인 & 전위 의존의 출처 & 갈라 주는 독립 진단 \\
\midrule
배리어 낮춤(본 장) & 유효 배리어 $-\chi\mathcal A_j$ & 구동항 보정 후 Arrhenius 기울기가 \emph{전위 무관} 하면 성립(식~\eqref{eq:arrhenius}) \\
고체상 확산 & $D(\theta)$ 의 점유율 의존 & GITT 완화가 $\sqrt t$(Cottrell) 형 — 계면 율속의 지수형과 다름 \\
전하전달 과전위 & $\eta_\mathrm{ct}\propto\ln i$ & EIS 의 $R_\mathrm{ct}$ 반원(또는 전류 차단)으로 독립 분리 \\
\bottomrule
\end{tabular}
\end{center}
즉 \textbf{반증의 칼날은 ``보정 후에도 남는 전위 의존''}이다 — 본 장의 barrier-lowering 이 옳다면 그것은 \emph{0}
이어야 하고, $\sqrt t$ 완화나 $R_\mathrm{ct}$ 반원이 보이면 그만큼을 먼저 떼어내야 한다. 이 분리 없이 단일
데이터셋에서 꼬리를 barrier-spectrum 으로 귀속하는 것은 금지한다.

\textbf{충방전 히스테리시스의 반증.} 본 장은 충방전 gap 을 \emph{열역학(상분리) $+$ 동역학(분극)}으로
갈랐다(식~\eqref{eq:hys_obsgap}). 이 분해도 같은 칼날로 반증된다:
\begin{enumerate}[label=(\roman*),leftmargin=2.0em,itemsep=2pt]
\item \textbf{율 외삽 절편.} 여러 율의 gap $\Delta V_{\mathrm{obs},j}(|I|)$ 를 $|I|\to0$ 으로 외삽한 \emph{절편}이 $0$ 이면
  순수 분극(상분리 없음, $\Omega_j\le2RT$)이라 본 장의 $\Omega_j>2RT$ 주장이 기각된다 — \emph{유한 절편}이 남아야
  상분리다. (분극만이면 식~\eqref{eq:hys_master} 가 방전식~\eqref{eq:master} 로 환원돼 히스가 닫혀야 한다.)
\item \textbf{온도 소멸.} 그 절편 $\gamma_j\Delta U_j^\hys(T)$ 는 $T\to T_{c,j}=\Omega_j/2R$ 서
  $\propto(T_{c,j}-T)^{3/2}$ 로 \emph{매끄럽게 소멸}해야 한다 — 식~\eqref{eq:hys_dU} 의 임계 극한
  $\Delta U_j^\hys\!\propto\!u_j^3$ 이 그것이다. 여기서 $u_j=\sqrt{1-T/T_{c,j}}$ 로 썼는데 $T_{c,j}=\Omega_j/2R$
  이므로 본문의 $u_j=\sqrt{1-2RT/\Omega_j}$ 와 \emph{같은 양}이다(항등). \emph{단 지수 $3/2$ 는 \textbf{평균장} 정규용액 값으로
  $|T-T_{c,j}|\ll T_{c,j}$ 임계근방에서만 유효}하다 — 흑연 staging 의 실측 온도창(예 15/23/45$^\circ$C)이 $T_{c,j}$ 에서
  \emph{충분히 멀면}($T_{c,j}\gg$ 창 — 창 전체가 $T_{c,j}$ 한참 \emph{아래}, 깊은 2상측) $u_j\to1$ 로 포화해 $3/2$ 멱이 안 보이므로, 그땐 정량 지수가 아니라
  \emph{절편 단조 감소와 $T_c$ 외삽 부호}의 약화 검정을 쓴다. \emph{역으로 $T_{c,j}$ 가 창 \emph{안/근처}면 $3/2$ 검정이
  오히려 유효}하다. \emph{수치 가늠으로 둘을 가른다}: 상온서 상분리가 \emph{보인다}는 관측은 $\Omega_j>2RT_{\mathrm{room}}$ 라는
  \textbf{하한}만 준다($T_{c,j}\gtrsim300$ K — 등호 아님). $T_{c,j}$ 가 측정창(288–318 K) 근처인지 훨씬 위인지는
  절편의 \emph{상대(로그) 온도 기울기} $|\dd\ln(\text{절편})/\dd T|$ 로 사전 판별한다. $T_{c,j}$ 가 근처면 이 양이
  $\sim\tfrac{3}{2(T_{c,j}-T)}$ 로 발산적으로 가파르므로 $3/2$ 정량 검정을 쓰고, 멀면 완만하므로 약화 검정을 쓴다.
  \emph{절대} 기울기 $|\dd\Delta U^\hys/\dd T|=(4R/sF)\,\mathrm{artanh}\,u_j$ 를 판별에 쓰면 안 되는 까닭은
  둘이다 — 원거리서 오히려 \emph{크고}(단조 증가), 절편엔 미지 $\gamma_j$ 가 곱해져 섞인다. 로그 기울기만이
  $\gamma_j$ 가 약분되는 $\gamma$-무관 판별량이다 —
  즉 $T_c$ 사전 추정이 어느 검정을 쓸지 정하는, 확인 가능한 gate 다. 평균장 임계지수의 정량값 자체는 본 장
  근사의 한계임을 덧붙인다. 절편이 온도에 무관하면 정규용액 상분리 그림이 틀렸다는 신호이며, 그때는 표면
  재구성$\cdot$기계적 응력 기원의 히스테리시스 같은 다른 기원을 의심한다.
\item \textbf{$\Omega_j$ 의 두 경로 내적 일관성.} 정규용액 $\Omega_j$ 는 두 관측에서 사영된다 — (a) 평형 peak
  \emph{폭} $w_j=(RT/F)(1-\Omega_j/2RT)$(식~\eqref{eq:weff}, S1 — \emph{단상 온도 구간}($T>T_{c,j}$ 측 측정)이 있을
  때만; 분기 영역의 현상학 폭은 $\Omega$ 정보가 없다), (b) 충방전 \emph{gap} 절편 식~\eqref{eq:hys_dU}(S5).
  \textbf{단, 둘은 \emph{같은} 정규용액 곡률 $g_j''$(식~\eqref{eq:regsol_gpp})의 두 평가점이라 진짜 \emph{독립} 측정이
  아니다} — 모형이 맞으면 자동 일치, 틀려도 같은 방향으로 틀린다. 따라서 이 일치는 \emph{독립 과결정}이 아니라
  \emph{단일 모형의 내적 일관성(self-consistency) 점검}이다. 둘이 불일치하면 폭 좁아짐이나 gap 중 하나가
  정규용액이 아닌 기원 — 입자 분산이나 표면 — 임을 시사한다. 독립 검정을 원하면 정규용액과 무관한 제3 입력,
  곧 열량계의 $\Omega$ 나 위상도의 $T_c$ 직접측정으로 (a)나 (b)를 치환해야 한다.
\end{enumerate}
정리하면, 히스테리시스의 반증 칼날은 ``$|I|\to0$ 절편의 유무, 임계 온도 소멸, 폭과 gap 의 $\Omega_j$
일치'' 셋이며, 모두 \S\ref{sec:regsol} 의 정규용액 $\Omega_j$ 한 양에서 정량 예측된다.

\medskip
이로써 본 장이 닫혔다. 방전 $\dd Q/\dd V$ 의 1차 원리 유도(\S\ref{sec:stage}--\S\ref{sec:dva})와 그 위의
충방전 확장(\S\ref{sec:hys_branch}--\S\ref{sec:hys_partial}), 통합 피팅식과 알고리즘(\S\ref{sec:master}),
그리고 본 절의 반증 절차까지 — 서론이 약속한 두 요구, 곧 ``실데이터 피팅에 바로 쓸 수 있을 것''과 ``1차
원리에서 빈틈없이 닫힐 것''이 각각 \S\ref{sec:master} 와 본 절에서 회수되었다.

% ====================================================================
\begin{thebibliography}{99}
\bibitem{dahn1991} J. R. Dahn, ``Phase diagram of Li$_x$C$_6$,'' \emph{Phys. Rev. B} \textbf{44}, 9170 (1991). DOI: 10.1103/PhysRevB.44.9170.
\bibitem{ohzuku1993} T. Ohzuku, Y. Iwakoshi, K. Sawai, ``Formation of Lithium-Graphite Intercalation Compounds in Nonaqueous Electrolytes,'' \emph{J. Electrochem. Soc.} \textbf{140}, 2490 (1993). DOI: 10.1149/1.2220849.
\bibitem{funabiki1999} A. Funabiki, M. Inaba, Z. Ogumi et al., ``Stage Transformation of Lithium-Graphite Intercalation Compounds Caused by Electrochemical Lithium Intercalation,'' \emph{J. Electrochem. Soc.} \textbf{146}, 2443 (1999). DOI: 10.1149/1.1391953.
\bibitem{dreyer2010} W. Dreyer \emph{et al.}, ``The thermodynamic origin of hysteresis in insertion batteries,'' \emph{Nature Materials} \textbf{9}, 448 (2010). DOI: 10.1038/nmat2730.
\bibitem{sethna1993} J. P. Sethna \emph{et al.}, ``Hysteresis and hierarchies: Dynamics of disorder-driven first-order phase transformations,'' \emph{Phys. Rev. Lett.} \textbf{70}, 3347 (1993). DOI: 10.1103/PhysRevLett.70.3347.
\bibitem{mckinnon1983} W. R. McKinnon, R. R. Haering, ``Physical Mechanisms of Intercalation,'' in \emph{Modern Aspects of Electrochemistry} \textbf{15}, 235 (1983). DOI: 10.1007/978-1-4615-7461-3\_4.
\bibitem{bazant2013} M. Z. Bazant, ``Theory of Chemical Kinetics and Charge Transfer based on Nonequilibrium Thermodynamics,'' \emph{Acc. Chem. Res.} \textbf{46}, 1144 (2013). DOI: 10.1021/ar300145c.
\bibitem{eyring1935} H. Eyring, ``The Activated Complex in Chemical Reactions,'' \emph{J. Chem. Phys.} \textbf{3}, 107 (1935). DOI: 10.1063/1.1749604.
\bibitem{marcus1956} R. A. Marcus, ``On the Theory of Oxidation-Reduction Reactions Involving Electron Transfer. I,'' \emph{J. Chem. Phys.} \textbf{24}, 966 (1956). DOI: 10.1063/1.1742723.
\bibitem{evanspolanyi1938} M. G. Evans, M. Polanyi, ``Inertia and driving force of chemical reactions,'' \emph{Trans. Faraday Soc.} \textbf{34}, 11 (1938). DOI: 10.1039/TF9383400011.
\bibitem{bardfaulkner} A. J. Bard, L. R. Faulkner, \emph{Electrochemical Methods: Fundamentals and Applications}, 2nd ed., Wiley (2001). ISBN 978-0-471-04372-0.
\bibitem{reynier2004} Y. Reynier, R. Yazami, B. Fultz, ``Thermodynamics of Lithium Intercalation into Graphites and Disordered Carbons,'' \emph{J. Electrochem. Soc.} \textbf{151}, A422 (2004). DOI: 10.1149/1.1646152.
\bibitem{lindsey1980} C. P. Lindsey, G. D. Patterson, ``Detailed comparison of the Williams--Watts and Cole--Davidson functions,'' \emph{J. Chem. Phys.} \textbf{73}, 3348 (1980). DOI: 10.1063/1.440530.
\bibitem{johnston2006} D. C. Johnston, ``Stretched exponential relaxation arising from a continuous sum of exponential decays,'' \emph{Phys. Rev. B} \textbf{74}, 184430 (2006). DOI: 10.1103/PhysRevB.74.184430.
\bibitem{svare2000} I. Svare, S. W. Martin, F. Borsa, ``Stretched exponentials with $T$-dependent exponents from fixed distributions of energy barriers for relaxation times in fast-ion conductors,'' \emph{Phys. Rev. B} \textbf{61}, 228 (2000). DOI: 10.1103/PhysRevB.61.228.
\bibitem{weppner1977} W. Weppner, R. A. Huggins, ``Determination of the Kinetic Parameters of Mixed-Conducting Electrodes (GITT),'' \emph{J. Electrochem. Soc.} \textbf{124}, 1569 (1977). DOI: 10.1149/1.2133112.
\bibitem{dubarry2012} M. Dubarry, C. Truchot, B. Y. Liaw, ``Synthesize battery degradation modes via a diagnostic and prognostic model,'' \emph{J. Power Sources} \textbf{219}, 204 (2012). DOI: 10.1016/j.jpowsour.2012.07.016.
\bibitem{newman2004} J. Newman, K. E. Thomas-Alyea, \emph{Electrochemical Systems}, 3rd ed., Wiley-Interscience (2004). ISBN 978-0-471-47756-3.
\bibitem{bloom2005} I. Bloom, A. N. Jansen, D. P. Abraham et al., ``Differential voltage analyses of high-power, lithium-ion cells. 1. Technique and application,'' \emph{J. Power Sources} \textbf{139}, 295 (2005). DOI: 10.1016/j.jpowsour.2004.07.021.
\end{thebibliography}

\end{document}
